%% ipc.tex
%%
%% Copyright (C) 2000-2018 Simone Piccardi.  Permission is granted to
%% copy, distribute and/or modify this document under the terms of the GNU Free
%% Documentation License, Version 1.1 or any later version published by the
%% Free Software Foundation; with the Invariant Sections being "Un preambolo",
%% with no Front-Cover Texts, and with no Back-Cover Texts.  A copy of the
%% license is included in the section entitled "GNU Free Documentation
%% License".
%%

\chapter{L'intercomunicazione fra processi}
\label{cha:IPC}


Uno degli aspetti fondamentali della programmazione in un sistema unix-like è
la comunicazione fra processi. In questo capitolo affronteremo solo i
meccanismi più elementari che permettono di mettere in comunicazione processi
diversi, come quelli tradizionali che coinvolgono \textit{pipe} e
\textit{fifo} e i meccanismi di intercomunicazione di System V e quelli POSIX.

Tralasceremo invece tutte le problematiche relative alla comunicazione
attraverso la rete (e le relative interfacce) che saranno affrontate in
dettaglio in un secondo tempo.  Non affronteremo neanche meccanismi più
complessi ed evoluti come le RPC (\textit{Remote Procedure Calls}) che in
genere sono implementati da un ulteriore livello di librerie sopra i
meccanismi elementari.


\section{L'intercomunicazione fra processi tradizionale}
\label{sec:ipc_unix}

Il primo meccanismo di comunicazione fra processi introdotto nei sistemi Unix,
è quello delle cosiddette \textit{pipe}; esse costituiscono una delle
caratteristiche peculiari del sistema, in particolar modo dell'interfaccia a
linea di comando. In questa sezione descriveremo le sue basi, le funzioni che
ne gestiscono l'uso e le varie forme in cui si è evoluto.


\subsection{Le \textit{pipe} standard}
\label{sec:ipc_pipes}

Le \textit{pipe} nascono sostanzialmente con Unix, e sono il primo, e tuttora
uno dei più usati, meccanismi di comunicazione fra processi. Si tratta in
sostanza di una coppia di file descriptor connessi fra di loro in modo che
quanto scrive su di uno si può rileggere dall'altro.  Si viene così a
costituire un canale di comunicazione realizzato tramite i due file
descriptor, che costituisce appunto una sorta di \textsl{tubo} (che appunto il
significato del termine inglese \textit{pipe}) attraverso cui si possono far
passare i dati.

In pratica si tratta di un buffer circolare in memoria in cui il kernel
appoggia i dati immessi nel file descriptor su cui si scrive per farli poi
riemergere dal file descriptor da cui si legge. Si tenga ben presente che in
questo passaggio di dati non è previsto nessun tipo di accesso al disco e che
nonostante l'uso dei file descriptor le \textit{pipe} non han nulla a che fare
con i file di dati di cui si è parlato al cap.~\ref{cha:file_IO_interface}.

La funzione di sistema che permette di creare questa speciale coppia di file
descriptor associati ad una \textit{pipe} è appunto \funcd{pipe}, ed il suo
prototipo è:

\begin{funcproto}{
\fhead{unistd.h}
\fdecl{int pipe(int filedes[2])}
\fdesc{Crea la coppia di file descriptor di una \textit{pipe}.} 
}

{La funzione ritorna $0$ in caso di successo e $-1$ per un errore, nel qual
  caso \var{errno} assumerà uno dei valori: 
  \begin{errlist}
  \item[\errcode{EFAULT}] \param{filedes} non è un indirizzo valido.
  \end{errlist}
  ed inoltre \errval{EMFILE} e \errval{ENFILE} nel loro significato generico.}
\end{funcproto}

La funzione restituisce la coppia di file descriptor nel vettore
\param{filedes}, il primo è aperto in lettura ed il secondo in scrittura. Come
accennato concetto di funzionamento di una \textit{pipe} è semplice: quello
che si scrive nel file descriptor aperto in scrittura viene ripresentato tale
e quale nel file descriptor aperto in lettura. 

I file descriptor infatti non sono connessi a nessun file reale, ma, come
accennato, ad un buffer nel kernel la cui dimensione è specificata dal
parametro di sistema \const{PIPE\_BUF}, (vedi
sez.~\ref{sec:sys_file_limits}). Lo schema di funzionamento di una
\textit{pipe} è illustrato in fig.~\ref{fig:ipc_pipe_singular}, in cui sono
indicati i due capi della \textit{pipe}, associati a ciascun file descriptor,
con le frecce che indicano la direzione del flusso dei dati.

\begin{figure}[!htb]
  \centering
  \includegraphics[height=5cm]{img/pipe}
  \caption{Schema della struttura di una \textit{pipe}.}
  \label{fig:ipc_pipe_singular}
\end{figure}

Della funzione di sistema esiste una seconda versione, \funcd{pipe2},
introdotta con il kernel 2.6.27 e la \acr{glibc} 2.9 e specifica di Linux
(utilizzabile solo definendo la macro \macro{\_GNU\_SOURCE}), che consente di
impostare atomicamente le caratteristiche dei file descriptor restituiti, il
suo prototipo è:

\begin{funcproto}{
\fhead{unistd.h}
\fhead{fcntl.h}
\fdecl{int pipe2(int pipefd[2], int flags)}
\fdesc{Crea la coppia di file descriptor di una \textit{pipe}.} 
}

{La funzione ritorna $0$ in caso di successo e $-1$ per un errore, nel qual
  caso \var{errno} assumerà uno dei valori: 
  \begin{errlist}
  \item[\errcode{EINVAL}] il valore di \param{flags} non valido.
  \end{errlist}
  e gli altri già visti per \func{pipe} con lo stesso significato.}
\end{funcproto}

Utilizzando un valore nullo per \param{flags} la funzione è identica a
\func{pipe}, si può però passare come valore l'OR aritmetico di uno qualunque
fra \const{O\_NONBLOCK} o \const{O\_CLOEXEC} che hanno l'effetto di impostare
su entrambi i file descriptor restituiti dalla funzione i relativi flag, già
descritti per \func{open} in tab.~\ref{tab:open_operation_flag}, che attivano
rispettivamente la modalità di accesso \textsl{non-bloccante} ed il
\textit{close-on-exec}.

Chiaramente creare una \textit{pipe} all'interno di un singolo processo non
serve a niente; se però ricordiamo quanto esposto in
sez.~\ref{sec:file_shared_access} riguardo al comportamento dei file
descriptor nei processi figli, è immediato capire come una \textit{pipe} possa
diventare un meccanismo di intercomunicazione. Un processo figlio infatti
condivide gli stessi file descriptor del padre, compresi quelli associati ad
una \textit{pipe} (secondo la situazione illustrata in
fig.~\ref{fig:ipc_pipe_fork}). In questo modo se uno dei processi scrive su un
capo della \textit{pipe}, l'altro può leggere.

\begin{figure}[!htb]
  \centering
  \includegraphics[height=5cm]{img/pipefork}
  \caption{Schema dei collegamenti ad una \textit{pipe}, condivisi fra
    processo padre e figlio dopo l'esecuzione \func{fork}.}
  \label{fig:ipc_pipe_fork}
\end{figure}

Tutto ciò ci mostra come sia immediato realizzare un meccanismo di
comunicazione fra processi attraverso una \textit{pipe}, utilizzando le
proprietà ordinarie dei file, ma ci mostra anche qual è il principale limite
nell'uso delle \textit{pipe}.\footnote{Stevens in \cite{APUE} riporta come
  limite anche il fatto che la comunicazione è unidirezionale, ma in realtà
  questo è un limite superabile usando una coppia di \textit{pipe}, anche se
  al costo di una maggiore complessità di gestione.}  È necessario infatti che
i processi possano condividere i file descriptor della \textit{pipe}, e per
questo essi devono comunque essere \textsl{parenti} (dall'inglese
\textit{siblings}), cioè o derivare da uno stesso processo padre in cui è
avvenuta la creazione della \textit{pipe}, o, più comunemente, essere nella
relazione padre/figlio.

A differenza di quanto avviene con i file normali, la lettura da una
\textit{pipe} può essere bloccante (qualora non siano presenti dati), inoltre
se si legge da una \textit{pipe} il cui capo in scrittura è stato chiuso, si
avrà la ricezione di un EOF (vale a dire che la funzione \func{read} ritornerà
restituendo 0).  Se invece si esegue una scrittura su una \textit{pipe} il cui
capo in lettura non è aperto il processo riceverà il segnale \signal{SIGPIPE},
e la funzione di scrittura restituirà un errore di \errcode{EPIPE} (al ritorno
del gestore, o qualora il segnale sia ignorato o bloccato).

La dimensione del buffer della \textit{pipe} (\const{PIPE\_BUF}) ci dà inoltre
un'altra importante informazione riguardo il comportamento delle operazioni di
lettura e scrittura su di una \textit{pipe}; esse infatti sono atomiche
fintanto che la quantità di dati da scrivere non supera questa
dimensione. Qualora ad esempio si effettui una scrittura di una quantità di
dati superiore l'operazione verrà effettuata in più riprese, consentendo
l'intromissione di scritture effettuate da altri processi.

La dimensione originale del buffer era di 4096 byte (uguale ad una pagina di
memoria) fino al kernel 2.6.11, ed è stata portata in seguito a 64kb; ma a
partire dal kernel 2.6.35 è stata resa disponibile l'operazione di controllo
\const{F\_SETPIPE\_SZ} (vedi sez.~\ref{sec:ipc_pipes}) che consente di
modificarne la dimensione.



\subsection{Un esempio dell'uso delle \textit{pipe}}
\label{sec:ipc_pipe_use}

Per capire meglio il funzionamento delle \textit{pipe} faremo un esempio di
quello che è il loro uso più comune, analogo a quello effettuato della shell,
e che consiste nell'inviare l'output di un processo (lo standard output)
sull'input di un altro. Realizzeremo il programma di esempio nella forma di un
\textit{CGI}\footnote{quella dei CGI (\textit{Common Gateway Interface}) è una
  interfaccia che consente ad un server web di eseguire un programma il cui
  output (che deve essere opportunamente formattato seguendo le specifiche
  dell'interfaccia) può essere presentato come risposta ad una richiesta HTTP
  al posto del contenuto di un file, e che ha costituito probabilmente la
  prima modalità con cui sono state create pagine HTML dinamiche.}  che genera
una immagine JPEG di un codice a barre, specificato come argomento in
ingresso.

Un programma che deve essere eseguito come \textit{CGI} deve rispondere a
delle caratteristiche specifiche, esso infatti non viene lanciato da una
shell, ma dallo stesso web server, alla richiesta di una specifica URL, che di
solito ha la forma:
\begin{Example}
http://www.sito.it/cgi-bin/programma?argomento
\end{Example}
ed il risultato dell'elaborazione deve essere presentato (con una intestazione
che ne descrive il \textit{mime-type}) sullo \textit{standard output}, in modo
che il server web possa reinviarlo al browser che ha effettuato la richiesta,
che in questo modo è in grado di visualizzarlo opportunamente.

\begin{figure}[!htb]
  \centering
  \includegraphics[height=5cm]{img/pipeuse}
  \caption{Schema dell'uso di una \textit{pipe} come mezzo di comunicazione fra
    due processi attraverso l'esecuzione una \func{fork} e la chiusura dei
    capi non utilizzati.}
  \label{fig:ipc_pipe_use}
\end{figure}

Per realizzare quanto voluto useremo in sequenza i programmi \cmd{barcode} e
\cmd{gs}, il primo infatti è in grado di generare immagini PostScript di
codici a barre corrispondenti ad una qualunque stringa, mentre il secondo
serve per poter effettuare la conversione della stessa immagine in formato
JPEG. Usando una \textit{pipe} potremo inviare l'output del primo sull'input
del secondo, secondo lo schema mostrato in fig.~\ref{fig:ipc_pipe_use}, in cui
la direzione del flusso dei dati è data dalle frecce continue.

Si potrebbe obiettare che sarebbe molto più semplice salvare il risultato
intermedio su un file temporaneo. Questo però non tiene conto del fatto che un
\textit{CGI} può essere eseguito più volte in contemporanea, e si avrebbe una
evidente \textit{race condition} in caso di accesso simultaneo a detto file da
istanze diverse. Il problema potrebbe essere superato utilizzando un sempre
diverso per il file temporaneo, che verrebbe creato all'avvio di ogni istanza,
utilizzato dai sottoprocessi, e cancellato alla fine della sua esecuzione; ma
a questo punto le cose non sarebbero più tanto semplici.  L'uso di una
\textit{pipe} invece permette di risolvere il problema in maniera semplice ed
elegante, oltre ad essere molto più efficiente, dato che non si deve scrivere
su disco.

Il programma ci servirà anche come esempio dell'uso delle funzioni di
duplicazione dei file descriptor che abbiamo trattato in
sez.~\ref{sec:file_dup}, in particolare di \func{dup2}. È attraverso queste
funzioni infatti che è possibile dirottare gli stream standard dei processi
(che abbiamo visto in tab.~\ref{tab:file_std_files} e
sez.~\ref{sec:file_stream}) sulla \textit{pipe}. In
fig.~\ref{fig:ipc_barcodepage_code} abbiamo riportato il corpo del programma,
il cui codice completo è disponibile nel file \file{BarCodePage.c} che si
trova nella directory dei sorgenti.

\begin{figure}[!htb]
  \footnotesize \centering
  \begin{minipage}[c]{\codesamplewidth}
    \includecodesample{listati/BarCodePage.c}
  \end{minipage} 
  \normalsize 
  \caption{Sezione principale del codice del \textit{CGI} 
    \file{BarCodePage.c}.}
  \label{fig:ipc_barcodepage_code}
\end{figure}

La prima operazione del programma (\texttt{\small 4-12}) è quella di creare
le due \textit{pipe} che serviranno per la comunicazione fra i due comandi
utilizzati per produrre il codice a barre; si ha cura di controllare la
riuscita della chiamata, inviando in caso di errore un messaggio invece
dell'immagine richiesta. La funzione \func{WriteMess} non è riportata in
fig.~\ref{fig:ipc_barcodepage_code}; essa si incarica semplicemente di
formattare l'uscita alla maniera dei CGI, aggiungendo l'opportuno
\textit{mime-type}, e formattando il messaggio in HTML, in modo che
quest'ultimo possa essere visualizzato correttamente da un browser.

Una volta create le \textit{pipe}, il programma può creare (\texttt{\small
  13-17}) il primo processo figlio, che si incaricherà (\texttt{\small
  19-25}) di eseguire \cmd{barcode}. Quest'ultimo legge dallo standard input
una stringa di caratteri, la converte nell'immagine PostScript del codice a
barre ad essa corrispondente, e poi scrive il risultato direttamente sullo
standard output.

Per poter utilizzare queste caratteristiche prima di eseguire \cmd{barcode} si
chiude (\texttt{\small 20}) il capo aperto in scrittura della prima
\textit{pipe}, e se ne collega (\texttt{\small 21}) il capo in lettura allo
\textit{standard input} usando \func{dup2}. Si ricordi che invocando
\func{dup2} il secondo file, qualora risulti aperto, viene, come nel caso
corrente, chiuso prima di effettuare la duplicazione. Allo stesso modo, dato
che \cmd{barcode} scrive l'immagine PostScript del codice a barre sullo
standard output, per poter effettuare una ulteriore redirezione il capo in
lettura della seconda \textit{pipe} viene chiuso (\texttt{\small 22}) mentre
il capo in scrittura viene collegato allo standard output (\texttt{\small
  23}).

In questo modo all'esecuzione (\texttt{\small 25}) di \cmd{barcode} (cui si
passa in \var{size} la dimensione della pagina per l'immagine) quest'ultimo
leggerà dalla prima \textit{pipe} la stringa da codificare che gli sarà
inviata dal padre, e scriverà l'immagine PostScript del codice a barre sulla
seconda.

Al contempo una volta lanciato il primo figlio, il processo padre prima chiude
(\texttt{\small 26}) il capo inutilizzato della prima \textit{pipe} (quello in
ingresso) e poi scrive (\texttt{\small 27}) la stringa da convertire sul capo
in uscita, così che \cmd{barcode} possa riceverla dallo \textit{standard
  input}. A questo punto l'uso della prima \textit{pipe} da parte del padre è
finito ed essa può essere definitivamente chiusa (\texttt{\small 28}), si
attende poi (\texttt{\small 29}) che l'esecuzione di \cmd{barcode} sia
completata.

Alla conclusione della sua esecuzione \cmd{barcode} avrà inviato l'immagine
PostScript del codice a barre sul capo in scrittura della seconda
\textit{pipe}; a questo punto si può eseguire la seconda conversione, da PS a
JPEG, usando il programma \cmd{gs}. Per questo si crea (\texttt{\small
  30-34}) un secondo processo figlio, che poi (\texttt{\small 35-42})
eseguirà questo programma leggendo l'immagine PostScript creata da
\cmd{barcode} dallo \textit{standard input}, per convertirla in JPEG.

Per fare tutto ciò anzitutto si chiude (\texttt{\small 37}) il capo in
scrittura della seconda \textit{pipe}, e se ne collega (\texttt{\small 38}) il
capo in lettura allo \textit{standard input}. Per poter formattare l'output
del programma in maniera utilizzabile da un browser, si provvede anche
\texttt{\small 40}) alla scrittura dell'apposita stringa di identificazione
del \textit{mime-type} in testa allo \textit{standard output}. A questo punto
si può invocare \texttt{\small 41}) \cmd{gs}, provvedendo le opportune opzioni
del comando che consentono di leggere il file da convertire dallo
\textit{standard input} e di inviare la conversione sullo \textit{standard
  output}.

Per completare le operazioni il processo padre chiude (\texttt{\small 44}) il
capo in scrittura della seconda \textit{pipe}, e attende la conclusione del
figlio (\texttt{\small 45}); a questo punto può (\texttt{\small 46})
uscire. Si tenga conto che l'operazione di chiudere il capo in scrittura della
seconda \textit{pipe} è necessaria, infatti, se non venisse chiusa, \cmd{gs},
che legge il suo \textit{standard input} da detta \textit{pipe}, resterebbe
bloccato in attesa di ulteriori dati in ingresso (l'unico modo che un
programma ha per sapere che i dati in ingresso sono terminati è rilevare che
lo \textit{standard input} è stato chiuso), e la \func{wait} non ritornerebbe.


\subsection{Le funzioni \func{popen} e \func{pclose}}
\label{sec:ipc_popen}

Come si è visto la modalità più comune di utilizzo di una \textit{pipe} è
quella di utilizzarla per fare da tramite fra output ed input di due programmi
invocati in sequenza; per questo motivo lo standard POSIX.2 ha introdotto due
funzioni che permettono di sintetizzare queste operazioni. La prima di esse si
chiama \funcd{popen} ed il suo prototipo è:


\begin{funcproto}{
\fhead{stdio.h}
\fdecl{FILE *popen(const char *command, const char *type)}
\fdesc{Esegue un programma dirottando l'uscita su una \textit{pipe}.} 
}

{La funzione ritorna l'indirizzo dello stream associato alla \textit{pipe} in
  caso di successo e \val{NULL} per un errore, nel qual caso \var{errno} potrà
  assumere i valori relativi alle sottostanti invocazioni di \func{pipe} e
  \func{fork} o \errcode{EINVAL} se \param{type} non è valido.}
\end{funcproto}

La funzione crea una \textit{pipe}, esegue una \func{fork} creando un nuovo
processo nel quale invoca il programma \param{command} attraverso la shell (in
sostanza esegue \file{/bin/sh} con il flag \code{-c}).
L'argomento \param{type} deve essere una delle due stringhe \verb|"w"| o
\verb|"r"|, per richiedere che la \textit{pipe} restituita come valore di
ritorno sia collegata allo \textit{standard input} o allo \textit{standard
  output} del comando invocato.

La funzione restituisce il puntatore ad uno stream associato alla
\textit{pipe} creata, che sarà aperto in sola lettura (e quindi associato allo
\textit{standard output} del programma indicato) in caso si sia indicato
\code{r}, o in sola scrittura (e quindi associato allo \textit{standard
  input}) in caso di \code{w}. A partire dalla versione 2.9 della \acr{glibc}
(questa è una estensione specifica di Linux) all'argomento \param{type} può
essere aggiunta la lettera ``\texttt{e}'' per impostare automaticamente il
flag di \textit{close-on-exec} sul file descriptor sottostante (si ricordi
quanto spiegato in sez.~\ref{sec:file_open_close}).

Lo \textit{stream} restituito da \func{popen} è identico a tutti gli effetti
ai \textit{file stream} visti in sez.~\ref{sec:files_std_interface}, anche se
è collegato ad una \textit{pipe} e non ad un file, e viene sempre aperto in
modalità \textit{fully-buffered} (vedi sez.~\ref{sec:file_buffering}); l'unica
differenza con gli usuali stream è che dovrà essere chiuso dalla seconda delle
due nuove funzioni, \funcd{pclose}, il cui prototipo è:

\begin{funcproto}{
\fhead{stdio.h}
\fdecl{int pclose(FILE *stream)}
\fdesc{Chiude una \textit{pipe} creata con \func{popen}.} 
}

{La funzione ritorna lo stato del processo creato da \func{popen} in caso di
  successo e $-1$ per un errore, nel qual caso \var{errno} assumerà i valori
  derivanti dalle sottostanti funzioni \func{fclose} e \func{wait4}.}
\end{funcproto}

La funzione chiude il file \param{stream} associato ad una \textit{pipe}
creato da una precedente \func{popen}, ed oltre alla chiusura dello stream si
incarica anche di attendere (tramite \func{wait4}) la conclusione del processo
creato dalla precedente \func{popen}. Se lo stato di uscita non può essere
letto la funzione restituirà per \var{errno} un errore di \errval{ECHILD}.

Per illustrare l'uso di queste due funzioni riprendiamo il problema
precedente: il programma mostrato in fig.~\ref{fig:ipc_barcodepage_code} per
quanto funzionante, è volutamente codificato in maniera piuttosto complessa,
inoltre doveva scontare un problema di \cmd{gs} che non era in grado di
riconoscere correttamente l'Encapsulated PostScript,\footnote{si fa
  riferimento alla versione di GNU Ghostscript 6.53 (2002-02-13), usata quando
  l'esempio venne scritto per la prima volta.} per cui si era utilizzato il
PostScript semplice, generando una pagina intera invece che una immagine delle
dimensioni corrispondenti al codice a barre.

Se si vuole generare una immagine di dimensioni appropriate si deve usare un
approccio diverso. Una possibilità sarebbe quella di ricorrere ad ulteriore
programma, \cmd{epstopsf}, per convertire in PDF un file EPS (che può essere
generato da \cmd{barcode} utilizzando l'opzione \cmd{-E}).  Utilizzando un PDF
al posto di un EPS \cmd{gs} esegue la conversione rispettando le dimensioni
originarie del codice a barre e produce un JPEG di dimensioni corrette.

Questo approccio però non può funzionare per via di una delle caratteristiche
principali delle \textit{pipe}. Per poter effettuare la conversione di un PDF
infatti è necessario, per la struttura del formato, potersi spostare (con
\func{lseek}) all'interno del file da convertire. Se si esegue la conversione
con \cmd{gs} su un file regolare non ci sono problemi, una \textit{pipe} però
è rigidamente sequenziale, e l'uso di \func{lseek} su di essa fallisce sempre
con un errore di \errcode{ESPIPE}, rendendo impossibile la conversione.
Questo ci dice che in generale la concatenazione di vari programmi funzionerà
soltanto quando tutti prevedono una lettura sequenziale del loro input.

Per questo motivo si è dovuto utilizzare un procedimento diverso, eseguendo
prima la conversione (sempre con \cmd{gs}) del PS in un altro formato
intermedio, il PPM,\footnote{il \textit{Portable PixMap file format} è un
  formato usato spesso come formato intermedio per effettuare conversioni, è
  infatti molto facile da manipolare, dato che usa caratteri ASCII per
  memorizzare le immagini, anche se per questo è estremamente inefficiente
  come formato di archiviazione.}  dal quale poi si può ottenere un'immagine
di dimensioni corrette attraverso vari programmi di manipolazione
(\cmd{pnmcrop}, \cmd{pnmmargin}) che può essere infine trasformata in PNG (con
\cmd{pnm2png}).

In questo caso però occorre eseguire in sequenza ben quattro comandi diversi,
inviando l'uscita di ciascuno all'ingresso del successivo, per poi ottenere il
risultato finale sullo \textit{standard output}: un caso classico di
utilizzazione delle \textit{pipe}, in cui l'uso di \func{popen} e \func{pclose}
permette di semplificare notevolmente la stesura del codice.

Nel nostro caso, dato che ciascun processo deve scrivere la sua uscita sullo
\textit{standard input} del successivo, occorrerà usare \func{popen} aprendo
la \textit{pipe} in scrittura. Il codice del nuovo programma è riportato in
fig.~\ref{fig:ipc_barcode_code}.  Come si può notare l'ordine di invocazione
dei programmi è l'inverso di quello in cui ci si aspetta che vengano
effettivamente eseguiti. Questo non comporta nessun problema dato che la
lettura su una \textit{pipe} è bloccante, per cui un processo, anche se
lanciato per primo, se non ottiene i dati che gli servono si bloccherà in
attesa sullo \textit{standard input} finché non otterrà il risultato
dell'elaborazione del processo che li deve creare, che pur essendo logicamente
precedente, viene lanciato dopo di lui.

\begin{figure}[!htb]
  \footnotesize \centering
  \begin{minipage}[c]{\codesamplewidth}
    \includecodesample{listati/BarCode.c}
  \end{minipage} 
  \normalsize 
  \caption{Codice completo del \textit{CGI} \file{BarCode.c}.}
  \label{fig:ipc_barcode_code}
\end{figure}

Nel nostro caso il primo passo (\texttt{\small 14}) è scrivere il
\textit{mime-type} sullo \textit{standard output}; a questo punto il processo
padre non necessita più di eseguire ulteriori operazioni sullo
\textit{standard output} e può tranquillamente provvedere alla redirezione.

Dato che i vari programmi devono essere lanciati in successione, si è
approntato un ciclo (\texttt{\small 15-19}) che esegue le operazioni in
sequenza: prima crea una \textit{pipe} (\texttt{\small 17}) per la scrittura
eseguendo il programma con \func{popen}, in modo che essa sia collegata allo
\textit{standard input}, e poi redirige (\texttt{\small 18}) lo
\textit{standard output} su detta \textit{pipe}.

In questo modo il primo processo ad essere invocato (che è l'ultimo della
catena) scriverà ancora sullo \textit{standard output} del processo padre, ma
i successivi, a causa di questa redirezione, scriveranno sulla \textit{pipe}
associata allo \textit{standard input} del processo invocato nel ciclo
precedente.

Alla fine tutto quello che resta da fare è lanciare (\texttt{\small 21}) il
primo processo della catena, che nel caso è \cmd{barcode}, e scrivere
(\texttt{\small 23}) la stringa del codice a barre sulla \textit{pipe}, che è
collegata al suo \textit{standard input}, infine si può eseguire
(\texttt{\small 24-27}) un ciclo che chiuda con \func{pclose}, nell'ordine
inverso rispetto a quello in cui le si sono create, tutte le \textit{pipe}
create in precedenza.


\subsection{Le \textit{pipe} con nome, o \textit{fifo}}
\label{sec:ipc_named_pipe}

Come accennato in sez.~\ref{sec:ipc_pipes} il problema delle \textit{pipe} è
che esse possono essere utilizzate solo da processi con un progenitore comune
o nella relazione padre/figlio. Per superare questo problema lo standard
POSIX.1 ha introdotto le \textit{fifo}, che hanno le stesse caratteristiche
delle \textit{pipe}, ma che invece di essere visibili solo attraverso un file
descriptor creato all'interno di un processo da una \textit{system call}
apposita, costituiscono un oggetto che risiede sul filesystem (si rammenti
quanto detto in sez.~\ref{sec:file_file_types}) che può essere aperto come un
qualunque file, così che i processi le possono usare senza dovere per forza
essere in una relazione di \textsl{parentela}.

Utilizzando una \textit{fifo} tutti i dati passeranno, come per le
\textit{pipe}, attraverso un buffer nel kernel, senza transitare dal
filesystem. Il fatto che siano associate ad un \textit{inode} presente sul
filesystem serve infatti solo a fornire un punto di accesso per i processi,
che permetta a questi ultimi di accedere alla stessa \textit{fifo} senza avere
nessuna relazione, con una semplice \func{open}. Il comportamento delle
funzioni di lettura e scrittura è identico a quello illustrato per le
\textit{pipe} in sez.~\ref{sec:ipc_pipes}.

Abbiamo già trattato in sez.~\ref{sec:file_mknod} le funzioni \func{mknod} e
\func{mkfifo} che permettono di creare una \textit{fifo}. Per utilizzarne una
un processo non avrà che da aprire il relativo file speciale o in lettura o
scrittura; nel primo caso il processo sarà collegato al capo di uscita della
\textit{fifo}, e dovrà leggere, nel secondo al capo di ingresso, e dovrà
scrivere.

Il kernel alloca un singolo buffer per ciascuna \textit{fifo} che sia stata
aperta, e questa potrà essere acceduta contemporaneamente da più processi, sia
in lettura che in scrittura. Dato che per funzionare deve essere aperta in
entrambe le direzioni, per una \textit{fifo} la funzione \func{open} di norma
si blocca se viene eseguita quando l'altro capo non è aperto.

Le \textit{fifo} però possono essere anche aperte in modalità
\textsl{non-bloccante}, nel qual caso l'apertura del capo in lettura avrà
successo solo quando anche l'altro capo è aperto, mentre l'apertura del capo
in scrittura restituirà l'errore di \errcode{ENXIO} fintanto che non verrà
aperto il capo in lettura.

In Linux è possibile aprire le \textit{fifo} anche in lettura/scrittura (lo
standard POSIX lascia indefinito il comportamento in questo caso) operazione
che avrà sempre successo immediato qualunque sia la modalità di apertura,
bloccante e non bloccante.  Questo può essere utilizzato per aprire comunque
una \textit{fifo} in scrittura anche se non ci sono ancora processi il
lettura. Infine è possibile anche usare la \textit{fifo} all'interno di un
solo processo, nel qual caso però occorre stare molto attenti alla possibili
situazioni di stallo: se si cerca di leggere da una \textit{fifo} che non
contiene dati si avrà infatti un \textit{deadlock} immediato, dato che il
processo si blocca e quindi non potrà mai eseguire le funzioni di scrittura.

Per la loro caratteristica di essere accessibili attraverso il filesystem, è
piuttosto frequente l'utilizzo di una \textit{fifo} come canale di
comunicazione nelle situazioni un processo deve ricevere informazioni da
altri. In questo caso è fondamentale che le operazioni di scrittura siano
atomiche; per questo si deve sempre tenere presente che questo è vero soltanto
fintanto che non si supera il limite delle dimensioni di \const{PIPE\_BUF} (si
ricordi quanto detto in sez.~\ref{sec:ipc_pipes}).

A parte il caso precedente, che resta probabilmente il più comune, Stevens
riporta in \cite{APUE} altre due casistiche principali per l'uso delle
\textit{fifo}:
\begin{itemize*}
\item Da parte dei comandi di shell, per evitare la creazione di file
  temporanei quando si devono inviare i dati di uscita di un processo
  sull'input di parecchi altri (attraverso l'uso del comando \cmd{tee}).  
\item Come canale di comunicazione fra un client ed un
  server (il modello \textit{client-server} è illustrato in
  sez.~\ref{sec:net_cliserv}).
\end{itemize*}

Nel primo caso quello che si fa è creare tante \textit{fifo} da usare come
\textit{standard input} quanti sono i processi a cui i vogliono inviare i
dati; questi ultimi saranno stati posti in esecuzione ridirigendo lo
\textit{standard input} dalle \textit{fifo}, si potrà poi eseguire il processo
che fornisce l'output replicando quest'ultimo, con il comando \cmd{tee}, sulle
varie \textit{fifo}.

Il secondo caso è relativamente semplice qualora si debba comunicare con un
processo alla volta, nel qual caso basta usare due \textit{fifo}, una per
leggere ed una per scrivere. Le cose diventano invece molto più complesse
quando si vuole effettuare una comunicazione fra un server ed un numero
imprecisato di client. Se il primo infatti può ricevere le richieste
attraverso una \textit{fifo} ``\textsl{nota}'', per le risposte non si può
fare altrettanto, dato che, per la struttura sequenziale delle \textit{fifo},
i client dovrebbero sapere prima di leggerli quando i dati inviati sono
destinati a loro.

Per risolvere questo problema, si può usare un'architettura come quella
illustrata in fig.~\ref{fig:ipc_fifo_server_arch} in cui i client
inviano le richieste al server su una \textit{fifo} nota mentre le
risposte vengono reinviate dal server a ciascuno di essi su una
\textit{fifo} temporanea creata per l'occasione.

\begin{figure}[!htb]
  \centering
  \includegraphics[height=9cm]{img/fifoserver}
  \caption{Schema dell'utilizzo delle \textit{fifo} nella realizzazione di una
    architettura di comunicazione client/server.}
  \label{fig:ipc_fifo_server_arch}
\end{figure}

Come esempio di uso questa architettura e dell'uso delle \textit{fifo},
abbiamo scritto un server di \textit{fortunes}, che restituisce, alle
richieste di un client, un detto a caso estratto da un insieme di frasi. Sia
il numero delle frasi dell'insieme, che i file da cui esse vengono lette
all'avvio, sono impostabili da riga di comando. Il corpo principale del
server è riportato in fig.~\ref{fig:ipc_fifo_server}, dove si è
tralasciata la parte che tratta la gestione delle opzioni a riga di comando,
che effettua l'impostazione delle variabili \var{fortunefilename}, che indica
il file da cui leggere le frasi, ed \var{n}, che indica il numero di frasi
tenute in memoria, ad un valore diverso da quelli preimpostati. Il codice
completo è nel file \file{FortuneServer.c}.

\begin{figure}[!htbp]
  \footnotesize \centering
  \begin{minipage}[c]{\codesamplewidth}
    \includecodesample{listati/FortuneServer.c}
  \end{minipage} 
  \normalsize 
  \caption{Sezione principale del codice del server di \textit{fortunes}
    basato sulle \textit{fifo}.}
  \label{fig:ipc_fifo_server}
\end{figure}

Il server richiede (\texttt{\small 12}) che sia stata impostata una dimensione
dell'insieme delle frasi non nulla, dato che l'inizializzazione del vettore
\var{fortune} avviene solo quando questa dimensione viene specificata, la
presenza di un valore nullo provoca l'uscita dal programma attraverso la
funzione (non riportata) che ne stampa le modalità d'uso.  Dopo di che
installa (\texttt{\small 13-15}) la funzione che gestisce i segnali di
interruzione (anche questa non è riportata in fig.~\ref{fig:ipc_fifo_server})
che si limita a rimuovere dal filesystem la \textit{fifo} usata dal server per
comunicare.

Terminata l'inizializzazione (\texttt{\small 16}) si effettua la chiamata alla
funzione \code{FortuneParse} che legge dal file specificato in
\var{fortunefilename} le prime \var{n} frasi e le memorizza (allocando
dinamicamente la memoria necessaria) nel vettore di puntatori \var{fortune}.
Anche il codice della funzione non è riportato, in quanto non direttamente
attinente allo scopo dell'esempio.

Il passo successivo (\texttt{\small 17-22}) è quello di creare con
\func{mkfifo} la \textit{fifo} nota sulla quale il server ascolterà le
richieste, qualora si riscontri un errore il server uscirà (escludendo
ovviamente il caso in cui la funzione \func{mkfifo} fallisce per la precedente
esistenza della \textit{fifo}).

Una volta che si è certi che la \textit{fifo} di ascolto esiste la procedura
di inizializzazione è completata. A questo punto (\texttt{\small 23}) si può
chiamare la funzione \func{daemon} per far proseguire l'esecuzione del
programma in background come demone.  Si può quindi procedere (\texttt{\small
  24-33}) alla apertura della \textit{fifo}: si noti che questo viene fatto
due volte, prima in lettura e poi in scrittura, per evitare di dover gestire
all'interno del ciclo principale il caso in cui il server è in ascolto ma non
ci sono client che effettuano richieste.  Si ricordi infatti che quando una
\textit{fifo} è aperta solo dal capo in lettura, l'esecuzione di \func{read}
ritorna con zero byte (si ha cioè una condizione di \textit{end-of-file}).

Nel nostro caso la prima apertura si bloccherà fintanto che un qualunque
client non apre a sua volta la \textit{fifo} nota in scrittura per effettuare
la sua richiesta. Pertanto all'inizio non ci sono problemi, il client però,
una volta ricevuta la risposta, uscirà, chiudendo tutti i file aperti,
compresa la \textit{fifo}.  A questo punto il server resta (se non ci sono
altri client che stanno effettuando richieste) con la \textit{fifo} chiusa sul
lato in lettura, ed in questo stato la funzione \func{read} non si bloccherà
in attesa di dati in ingresso, ma ritornerà in continuazione, restituendo una
condizione di \textit{end-of-file}.

Si è usata questa tecnica per compatibilità, Linux infatti supporta l'apertura
delle \textit{fifo} in lettura/scrittura, per cui si sarebbe potuto effettuare
una singola apertura con \const{O\_RDWR}; la doppia apertura comunque ha il
vantaggio che non si può scrivere per errore sul capo aperto in sola lettura.

Per questo motivo, dopo aver eseguito l'apertura in lettura (\texttt{\small
  24-28}),\footnote{di solito si effettua l'apertura del capo in lettura di
  una \textit{fifo} in modalità non bloccante, per evitare il rischio di uno
  stallo: se infatti nessuno apre la \textit{fifo} in scrittura il processo
  non ritornerà mai dalla \func{open}. Nel nostro caso questo rischio non
  esiste, mentre è necessario potersi bloccare in lettura in attesa di una
  richiesta.}  si esegue una seconda apertura in scrittura (\texttt{\small
  29-32}), scartando il relativo file descriptor, che non sarà mai usato, in
questo modo però la \textit{fifo} resta comunque aperta anche in scrittura,
cosicché le successive chiamate a \func{read} possono bloccarsi.

A questo punto si può entrare nel ciclo principale del programma che fornisce
le risposte ai client (\texttt{\small 34-50}); questo viene eseguito
indefinitamente (l'uscita del server viene effettuata inviando un segnale, in
modo da passare attraverso la funzione di chiusura che cancella la
\textit{fifo}).

Il server è progettato per accettare come richieste dai client delle stringhe
che contengono il nome della \textit{fifo} sulla quale deve essere inviata la
risposta.  Per cui prima (\texttt{\small 35-39}) si esegue la lettura dalla
stringa di richiesta dalla \textit{fifo} nota (che a questo punto si bloccherà
tutte le volte che non ci sono richieste). Dopo di che, una volta terminata la
stringa (\texttt{\small 40}) e selezionato (\texttt{\small 41}) un numero
casuale per ricavare la frase da inviare, si procederà (\texttt{\small
  42-46}) all'apertura della \textit{fifo} per la risposta, che poi 
(\texttt{\small 47-48}) vi sarà scritta. Infine (\texttt{\small 49}) si
chiude la \textit{fifo} di risposta che non serve più.

Il codice del client è invece riportato in fig.~\ref{fig:ipc_fifo_client},
anche in questo caso si è omessa la gestione delle opzioni e la funzione che
stampa a video le informazioni di utilizzo ed esce, riportando solo la sezione
principale del programma e le definizioni delle variabili. Il codice completo
è nel file \file{FortuneClient.c} dei sorgenti allegati.

\begin{figure}[!htbp]
  \footnotesize \centering
  \begin{minipage}[c]{\codesamplewidth}
    \includecodesample{listati/FortuneClient.c}
  \end{minipage} 
  \normalsize 
  \caption{Sezione principale del codice del client di \textit{fortunes}
    basato sulle \textit{fifo}.}
  \label{fig:ipc_fifo_client}
\end{figure}

La prima istruzione (\texttt{\small 12}) compone il nome della \textit{fifo}
che dovrà essere utilizzata per ricevere la risposta dal server.  Si usa il
\ids{PID} del processo per essere sicuri di avere un nome univoco; dopo di che
(\texttt{\small 13-18}) si procede alla creazione del relativo file, uscendo
in caso di errore (a meno che il file non sia già presente sul filesystem).

A questo punto il client può effettuare l'interrogazione del server, per
questo prima si apre la \textit{fifo} nota (\texttt{\small 19-23}), e poi ci
si scrive (\texttt{\small 24}) la stringa composta in precedenza, che contiene
il nome della \textit{fifo} da utilizzare per la risposta. Infine si richiude
la \textit{fifo} del server che a questo punto non serve più (\texttt{\small
  25}). 

Inoltrata la richiesta si può passare alla lettura della risposta; anzitutto
si apre (\texttt{\small 26-30}) la \textit{fifo} appena creata, da cui si
deve riceverla, dopo di che si effettua una lettura (\texttt{\small 31})
nell'apposito buffer; si è supposto, come è ragionevole, che le frasi inviate
dal server siano sempre di dimensioni inferiori a \const{PIPE\_BUF},
tralasciamo la gestione del caso in cui questo non è vero. Infine si stampa
(\texttt{\small 32}) a video la risposta, si chiude (\texttt{\small 33}) la
\textit{fifo} e si cancella (\texttt{\small 34}) il relativo file.  Si noti
come la \textit{fifo} per la risposta sia stata aperta solo dopo aver inviato
la richiesta, se non si fosse fatto così si avrebbe avuto uno stallo, in
quanto senza la richiesta, il server non avrebbe potuto aprirne il capo in
scrittura e l'apertura si sarebbe bloccata indefinitamente.

Verifichiamo allora il comportamento dei nostri programmi, in questo, come in
altri esempi precedenti, si fa uso delle varie funzioni di servizio, che sono
state raccolte nella libreria \file{libgapil.so}, e per poterla usare
occorrerà definire la variabile di ambiente \envvar{LD\_LIBRARY\_PATH} in modo
che il linker dinamico possa accedervi.

In generale questa variabile indica il \textit{pathname} della directory
contenente la libreria. Nell'ipotesi (che daremo sempre per verificata) che si
facciano le prove direttamente nella directory dei sorgenti (dove di norma
vengono creati sia i programmi che la libreria), il comando da dare sarà
\code{export LD\_LIBRARY\_PATH=./}; a questo punto potremo lanciare il server,
facendogli leggere una decina di frasi, con:
\begin{Console}
[piccardi@gont sources]$ \textbf{./fortuned -n10}
\end{Console}
%$

Avendo usato \func{daemon} per eseguire il server in background il comando
ritornerà immediatamente, ma potremo verificare con \cmd{ps} che in effetti il
programma resta un esecuzione in background, e senza avere associato un
terminale di controllo (si ricordi quanto detto in sez.~\ref{sec:sess_daemon}):
\begin{Console}
[piccardi@gont sources]$ \textbf{ps aux}
...
piccardi 27489  0.0  0.0  1204  356 ?        S    01:06   0:00 ./fortuned -n10
piccardi 27492  3.0  0.1  2492  764 pts/2    R    01:08   0:00 ps aux
\end{Console}
%$
e si potrà verificare anche che in \file{/tmp} è stata creata la \textit{fifo}
di ascolto \file{fortune.fifo}. A questo punto potremo interrogare il server
con il programma client; otterremo così:
\begin{Console}
[piccardi@gont sources]$ \textbf{./fortune}
Linux ext2fs has been stable for a long time, now it's time to break it
        -- Linuxkongreß '95 in Berlin
[piccardi@gont sources]$ \textbf{./fortune}
Let's call it an accidental feature.
        --Larry Wall
[piccardi@gont sources]$ \textbf{./fortune}
.........    Escape the 'Gates' of Hell
  `:::'                  .......  ......
   :::  *                  `::.    ::'
   ::: .::  .:.::.  .:: .::  `::. :'
   :::  ::   ::  ::  ::  ::    :::.
   ::: .::. .::  ::.  `::::. .:'  ::.
...:::.....................::'   .::::..
        -- William E. Roadcap
[piccardi@gont sources]$ ./fortune
Linux ext2fs has been stable for a long time, now it's time to break it
        -- Linuxkongreß '95 in Berlin
\end{Console}
%$
e ripetendo varie volte il comando otterremo, in ordine casuale, le dieci
frasi tenute in memoria dal server.

Infine per chiudere il server basterà inviargli un segnale di terminazione (ad
esempio con \cmd{killall fortuned}) e potremo verificare che il gestore del
segnale ha anche correttamente cancellato la \textit{fifo} di ascolto da
\file{/tmp}.

Benché il nostro sistema client-server funzioni, la sua struttura è piuttosto
complessa e continua ad avere vari inconvenienti\footnote{lo stesso Stevens,
  che esamina questa architettura in \cite{APUE}, nota come sia impossibile
  per il server sapere se un client è andato in crash, con la possibilità di
  far restare le \textit{fifo} temporanee sul filesystem, di come sia
  necessario intercettare \signal{SIGPIPE} dato che un client può terminare
  dopo aver fatto una richiesta, ma prima che la risposta sia inviata (cosa
  che nel nostro esempio non è stata fatta).}; in generale infatti
l'interfaccia delle \textit{fifo} non è adatta a risolvere questo tipo di
problemi, che possono essere affrontati in maniera più semplice ed efficace o
usando i socket (che tratteremo in dettaglio a partire da
cap.~\ref{cha:socket_intro}) o ricorrendo a meccanismi di comunicazione
diversi, come quelli che esamineremo in seguito.



\subsection{La funzione \func{socketpair}}
\label{sec:ipc_socketpair}

Un meccanismo di comunicazione molto simile alle \textit{pipe}, ma che non
presenta il problema della unidirezionalità del flusso dei dati, è quello dei
cosiddetti \textsl{socket locali} (o \textit{Unix domain socket}).  Tratteremo
in generale i socket in cap.~\ref{cha:socket_intro}, nell'ambito
dell'interfaccia che essi forniscono per la programmazione di rete, e vedremo
anche (in~sez.~\ref{sec:sock_sa_local}) come si possono utilizzare i file
speciali di tipo socket, analoghi a quelli associati alle \textit{fifo} (si
rammenti sez.~\ref{sec:file_file_types}) cui si accede però attraverso quella
medesima interfaccia; vale però la pena esaminare qui una modalità di uso dei
socket locali che li rende sostanzialmente identici ad una \textit{pipe}
bidirezionale.

La funzione di sistema \funcd{socketpair}, introdotta da BSD ma supportata in
genere da qualunque sistema che fornisca l'interfaccia dei socket ed inclusa
in POSIX.1-2001, consente infatti di creare una coppia di file descriptor
connessi fra loro (tramite un socket, appunto) senza dover ricorrere ad un
file speciale sul filesystem. I descrittori sono del tutto analoghi a quelli
che si avrebbero con una chiamata a \func{pipe}, con la sola differenza è che
in questo caso il flusso dei dati può essere effettuato in entrambe le
direzioni. Il prototipo della funzione è:

\begin{funcproto}{
\fhead{sys/types.h} 
\fhead{sys/socket.h}
\fdecl{int socketpair(int domain, int type, int protocol, int sv[2])}
\fdesc{Crea una coppia di socket connessi fra loro.} 
}

{La funzione ritorna $0$ in caso di successo e $-1$ per un errore, nel qual
  caso \var{errno} assumerà uno dei valori: 
  \begin{errlist}
  \item[\errcode{EAFNOSUPPORT}] i socket locali non sono supportati.
  \item[\errcode{EOPNOTSUPP}] il protocollo specificato non supporta la
  creazione di coppie di socket.
  \item[\errcode{EPROTONOSUPPORT}] il protocollo specificato non è supportato.
  \end{errlist}
  ed inoltre \errval{EFAULT}, \errval{EMFILE} e \errval{ENFILE} nel loro
  significato generico.}
\end{funcproto}

La funzione restituisce in \param{sv} la coppia di descrittori connessi fra di
loro: quello che si scrive su uno di essi sarà ripresentato in input
sull'altro e viceversa. Gli argomenti \param{domain}, \param{type} e
\param{protocol} derivano dall'interfaccia dei socket (vedi
sez.~\ref{sec:sock_creation}) che è quella che fornisce il substrato per
connettere i due descrittori, ma in questo caso i soli valori validi che
possono essere specificati sono rispettivamente \const{AF\_UNIX},
\const{SOCK\_STREAM} e \val{0}.  

A partire dal kernel 2.6.27 la funzione supporta nell'indicazione del tipo di
socket anche i due flag \const{SOCK\_NONBLOCK} e \const{SOCK\_CLOEXEC}
(trattati in sez.~\ref{sec:sock_type}), con effetto identico agli analoghi
\const{O\_CLOEXEC} e \const{O\_NONBLOCK} di una \func{open} (vedi
tab.~\ref{tab:open_operation_flag}).

L'utilità di chiamare questa funzione per evitare due chiamate a \func{pipe}
può sembrare limitata; in realtà l'utilizzo di questa funzione (e dei socket
locali in generale) permette di trasmettere attraverso le linea non solo dei
dati, ma anche dei file descriptor: si può cioè passare da un processo ad un
altro un file descriptor, con una sorta di duplicazione dello stesso non
all'interno di uno stesso processo, ma fra processi distinti (torneremo su
questa funzionalità in sez.~\ref{sec:sock_fd_passing}).


\section{L'intercomunicazione fra processi di System V}
\label{sec:ipc_sysv}

Benché le \textit{pipe} e le \textit{fifo} siano ancora ampiamente usate, esse
scontano il limite fondamentale che il meccanismo di comunicazione che
forniscono è rigidamente sequenziale: una situazione in cui un processo scrive
qualcosa che molti altri devono poter leggere non può essere implementata con
una \textit{pipe}.

Per questo nello sviluppo di System V vennero introdotti una serie di nuovi
oggetti per la comunicazione fra processi ed una nuova interfaccia di
programmazione, poi inclusa anche in POSIX.1-2001, che fossero in grado di
garantire una maggiore flessibilità.  In questa sezione esamineremo come Linux
supporta quello che viene chiamato il \textsl{Sistema di comunicazione fra
  processi} di System V, cui da qui in avanti faremo riferimento come
\textit{SysV-IPC} (dove IPC è la sigla di \textit{Inter-Process
  Comunication}).



\subsection{Considerazioni generali}
\label{sec:ipc_sysv_generic}

La principale caratteristica del \textit{SysV-IPC} è quella di essere basato
su oggetti permanenti che risiedono nel kernel. Questi, a differenza di quanto
avviene per i file descriptor, non mantengono un contatore dei riferimenti, e
non vengono cancellati dal sistema una volta che non sono più in uso.  Questo
comporta due problemi: il primo è che, al contrario di quanto avviene per
\textit{pipe} e \textit{fifo}, la memoria allocata per questi oggetti non
viene rilasciata automaticamente quando non c'è più nessuno che li utilizzi ed
essi devono essere cancellati esplicitamente, se non si vuole che restino
attivi fino al riavvio del sistema. Il secondo problema è, dato che non c'è
come per i file un contatore del numero di riferimenti che ne indichi l'essere
in uso, che essi possono essere cancellati anche se ci sono dei processi che
li stanno utilizzando, con tutte le conseguenze (ovviamente assai sgradevoli)
del caso.

Un'ulteriore caratteristica negativa è che gli oggetti usati nel
\textit{SysV-IPC} vengono creati direttamente dal kernel, e sono accessibili
solo specificando il relativo \textsl{identificatore}. Questo è un numero
progressivo (un po' come il \ids{PID} dei processi) che il kernel assegna a
ciascuno di essi quanto vengono creati (sul procedimento di assegnazione
torneremo in sez.~\ref{sec:ipc_sysv_id_use}). L'identificatore viene
restituito dalle funzioni che creano l'oggetto, ed è quindi locale al processo
che le ha eseguite. Dato che l'identificatore viene assegnato dinamicamente
dal kernel non è possibile prevedere quale sarà, né utilizzare un qualche
valore statico, si pone perciò il problema di come processi diversi possono
accedere allo stesso oggetto.

Per risolvere il problema nella struttura \struct{ipc\_perm} che il kernel
associa a ciascun oggetto, viene mantenuto anche un campo apposito che
contiene anche una \textsl{chiave}, identificata da una variabile del tipo
primitivo \type{key\_t}, da specificare in fase di creazione dell'oggetto, e
tramite la quale è possibile ricavare l'identificatore.\footnote{in sostanza
  si sposta il problema dell'accesso dalla classificazione in base
  all'identificatore alla classificazione in base alla chiave, una delle tante
  complicazioni inutili presenti nel \textit{SysV-IPC}.} Oltre la chiave, la
struttura, la cui definizione è riportata in fig.~\ref{fig:ipc_ipc_perm},
mantiene varie proprietà ed informazioni associate all'oggetto.

\begin{figure}[!htb]
  \footnotesize \centering
  \begin{minipage}[c]{0.80\textwidth}
    \includestruct{listati/ipc_perm.h}
  \end{minipage} 
  \normalsize 
  \caption{La struttura \structd{ipc\_perm}, come definita in
    \headfiled{sys/ipc.h}.}
  \label{fig:ipc_ipc_perm}
\end{figure}

Usando la stessa chiave due processi diversi possono ricavare l'identificatore
associato ad un oggetto ed accedervi. Il problema che sorge a questo punto è
come devono fare per accordarsi sull'uso di una stessa chiave. Se i processi
sono \textsl{imparentati} la soluzione è relativamente semplice, in tal caso
infatti si può usare il valore speciale \constd{IPC\_PRIVATE} per creare un
nuovo oggetto nel processo padre, l'identificatore così ottenuto sarà
disponibile in tutti i figli, e potrà essere passato come argomento attraverso
una \func{exec}.

Però quando i processi non sono \textsl{imparentati} (come capita tutte le
volte che si ha a che fare con un sistema client-server) tutto questo non è
possibile; si potrebbe comunque salvare l'identificatore su un file noto, ma
questo ovviamente comporta lo svantaggio di doverselo andare a rileggere.  Una
alternativa più efficace è quella che i programmi usino un valore comune per
la chiave (che ad esempio può essere dichiarato in un header comune), ma c'è
sempre il rischio che questa chiave possa essere stata già utilizzata da
qualcun altro.  Dato che non esiste una convenzione su come assegnare queste
chiavi in maniera univoca l'interfaccia mette a disposizione una funzione
apposita, \funcd{ftok}, che permette di ottenere una chiave specificando il
nome di un file ed un numero di versione; il suo prototipo è:

\begin{funcproto}{
\fhead{sys/types.h} 
\fhead{sys/ipc.h} 
\fdecl{key\_t ftok(const char *pathname, int proj\_id)}
\fdesc{Restituisce una chiave per identificare un oggetto del \textit{SysV
    IPC}.}}

{La funzione ritorna la chiave in caso di successo e $-1$ per un errore, nel
  qual caso \var{errno} assumerà uno dei possibili codici di errore di
  \func{stat}.}
\end{funcproto}

La funzione determina un valore della chiave sulla base di \param{pathname},
che deve specificare il \textit{pathname} di un file effettivamente esistente
e di un numero di progetto \param{proj\_id)}, che di norma viene specificato
come carattere, dato che ne vengono utilizzati solo gli 8 bit meno
significativi. Nelle \acr{libc4} e \acr{libc5}, come avviene in SunOS,
l'argomento \param{proj\_id} è dichiarato tipo \ctyp{char}, la \acr{glibc} usa
il prototipo specificato da XPG4, ma vengono lo stesso utilizzati gli 8 bit
meno significativi.

Il problema è che anche così non c'è la sicurezza che il valore della chiave
sia univoco, infatti esso è costruito combinando il byte di \param{proj\_id)}
con i 16 bit meno significativi dell'inode del file \param{pathname} (che
vengono ottenuti attraverso \func{stat}, da cui derivano i possibili errori),
e gli 8 bit meno significativi del numero del dispositivo su cui è il file.
Diventa perciò relativamente facile ottenere delle collisioni, specie se i
file sono su dispositivi con lo stesso \textit{minor number}, come
\file{/dev/hda1} e \file{/dev/sda1}.

In genere quello che si fa è utilizzare un file comune usato dai programmi che
devono comunicare (ad esempio un header comune, o uno dei programmi che devono
usare l'oggetto in questione), utilizzando il numero di progetto per ottenere
le chiavi che interessano. In ogni caso occorre sempre controllare, prima di
creare un oggetto, che la chiave non sia già stata utilizzata. Se questo va
bene in fase di creazione, le cose possono complicarsi per i programmi che
devono solo accedere, in quanto, a parte gli eventuali controlli sugli altri
attributi di \struct{ipc\_perm}, non esiste una modalità semplice per essere
sicuri che l'oggetto associato ad una certa chiave sia stato effettivamente
creato da chi ci si aspetta.

Questo è, insieme al fatto che gli oggetti sono permanenti e non mantengono un
contatore di riferimenti per la cancellazione automatica, il principale
problema del \textit{SysV-IPC}. Non esiste infatti una modalità chiara per
identificare un oggetto, come sarebbe stato se lo si fosse associato ad in
file, e tutta l'interfaccia è inutilmente complessa.  Per questo se ne
sconsiglia assolutamente l'uso nei nuovi programmi, considerato che è ormai
disponibile una revisione completa dei meccanismi di IPC fatta secondo quanto
indicato dallo standard POSIX.1b, che presenta una realizzazione più sicura ed
una interfaccia più semplice, che tratteremo in sez.~\ref{sec:ipc_posix}.


\subsection{Il controllo di accesso}
\label{sec:ipc_sysv_access_control}

Oltre alle chiavi, abbiamo visto in fig.~\ref{fig:ipc_ipc_perm} che ad ogni
oggetto sono associate in \struct{ipc\_perm} ulteriori informazioni, come gli
identificatori del creatore (nei campi \var{cuid} e \var{cgid}) e del
proprietario (nei campi \var{uid} e \var{gid}) dello stesso, e un insieme di
permessi (nel campo \var{mode}). In questo modo è possibile definire un
controllo di accesso sugli oggetti di IPC, simile a quello che si ha per i
file (vedi sez.~\ref{sec:file_perm_overview}).

Benché questo controllo di accesso sia molto simile a quello dei file, restano
delle importanti differenze. La prima è che il permesso di esecuzione non
esiste (e se specificato viene ignorato), per cui si può parlare solo di
permessi di lettura e scrittura (nel caso dei semafori poi quest'ultimo è più
propriamente un permesso di modifica). I valori di \var{mode} sono gli stessi
ed hanno lo stesso significato di quelli riportati in
tab.~\ref{tab:file_mode_flags} e come per i file definiscono gli accessi per
il proprietario, il suo gruppo e tutti gli altri. 

Se però si vogliono usare le costanti simboliche di
tab.~\ref{tab:file_mode_flags} occorrerà includere anche il file
\headfile{sys/stat.h}; alcuni sistemi definiscono le costanti \constd{MSG\_R}
(il valore ottale \texttt{0400}) e \constd{MSG\_W} (il valore ottale
\texttt{0200}) per indicare i permessi base di lettura e scrittura per il
proprietario, da utilizzare, con gli opportuni shift, pure per il gruppo e gli
altri. In Linux, visto la loro scarsa utilità, queste costanti non sono
definite.

Quando l'oggetto viene creato i campi \var{cuid} e \var{uid} di
\struct{ipc\_perm} ed i campi \var{cgid} e \var{gid} vengono impostati
rispettivamente al valore dell'\ids{UID} e del \ids{GID} effettivo del processo
che ha chiamato la funzione, ma, mentre i campi \var{uid} e \var{gid} possono
essere cambiati, i campi \var{cuid} e \var{cgid} restano sempre gli stessi.

Il controllo di accesso è effettuato a due livelli. Il primo livello è nelle
funzioni che richiedono l'identificatore di un oggetto data la chiave. Queste
specificano tutte un argomento \param{flag}, in tal caso quando viene
effettuata la ricerca di una chiave, qualora \param{flag} specifichi dei
permessi, questi vengono controllati e l'identificatore viene restituito solo
se corrispondono a quelli dell'oggetto. Se ci sono dei permessi non presenti
in \var{mode} l'accesso sarà negato. Questo controllo però è di utilità
indicativa, dato che è sempre possibile specificare per \param{flag} un valore
nullo, nel qual caso l'identificatore sarà restituito comunque.

Il secondo livello di controllo è quello delle varie funzioni che accedono
direttamente (in lettura o scrittura) all'oggetto. In tal caso lo schema dei
controlli è simile a quello dei file, ed avviene secondo questa sequenza:
\begin{itemize*}
\item se il processo ha i privilegi di amministratore (più precisamente
  \const{CAP\_IPC\_OWNER}) l'accesso è sempre consentito.
\item se l'\ids{UID} effettivo del processo corrisponde o al valore del campo
  \var{cuid} o a quello del campo \var{uid} ed il permesso per il proprietario
  in \var{mode} è appropriato\footnote{per appropriato si intende che è
    impostato il permesso di scrittura per le operazioni di scrittura e quello
    di lettura per le operazioni di lettura.} l'accesso è consentito.
\item se il \ids{GID} effettivo del processo corrisponde o al
  valore del campo \var{cgid} o a quello del campo \var{gid} ed il permesso
  per il gruppo in \var{mode} è appropriato l'accesso è consentito.
\item se il permesso per gli altri è appropriato l'accesso è consentito.
\end{itemize*}
solo se tutti i controlli elencati falliscono l'accesso è negato. Si noti che
a differenza di quanto avviene per i permessi dei file, fallire in uno dei
passi elencati non comporta il fallimento dell'accesso. Un'ulteriore
differenza rispetto a quanto avviene per i file è che per gli oggetti di IPC
il valore di \textit{umask} (si ricordi quanto esposto in
sez.~\ref{sec:file_perm_management}) non ha alcun significato.


\subsection{Gli identificatori ed il loro utilizzo}
\label{sec:ipc_sysv_id_use}

L'unico campo di \struct{ipc\_perm} del quale non abbiamo ancora parlato è
\var{seq}, che in fig.~\ref{fig:ipc_ipc_perm} è qualificato con un criptico
``\textsl{numero di sequenza}'', ne parliamo adesso dato che esso è
strettamente attinente alle modalità con cui il kernel assegna gli
identificatori degli oggetti del sistema di IPC.

Quando il sistema si avvia, alla creazione di ogni nuovo oggetto di IPC viene
assegnato un numero progressivo, pari al numero di oggetti di quel tipo
esistenti. Se il comportamento fosse sempre questo sarebbe identico a quello
usato nell'assegnazione dei file descriptor nei processi, ed i valori degli
identificatori tenderebbero ad essere riutilizzati spesso e restare di piccole
dimensioni (inferiori al numero massimo di oggetti disponibili).

Questo va benissimo nel caso dei file descriptor, che sono locali ad un
processo, ma qui il comportamento varrebbe per tutto il sistema, e per
processi del tutto scorrelati fra loro. Così si potrebbero avere situazioni
come quella in cui un server esce e cancella le sue code di messaggi, ed il
relativo identificatore viene immediatamente assegnato a quelle di un altro
server partito subito dopo, con la possibilità che i client del primo non
facciano in tempo ad accorgersi dell'avvenuto, e finiscano con l'interagire
con gli oggetti del secondo, con conseguenze imprevedibili.

Proprio per evitare questo tipo di situazioni il sistema usa il valore di
\var{seq} per provvedere un meccanismo che porti gli identificatori ad
assumere tutti i valori possibili, rendendo molto più lungo il periodo in cui
un identificatore può venire riutilizzato.

Il sistema dispone sempre di un numero fisso di oggetti di IPC, fino al kernel
2.2.x questi erano definiti dalle costanti \const{MSGMNI}, \const{SEMMNI} e
\const{SHMMNI}, e potevano essere cambiati (come tutti gli altri limiti
relativi al \textit{SysV-IPC}) solo con una ricompilazione del kernel.  A
partire dal kernel 2.4.x è possibile cambiare questi valori a sistema attivo
scrivendo sui file \sysctlrelfile{kernel}{shmmni},
\sysctlrelfile{kernel}{msgmni} e \sysctlrelfiled{kernel}{sem} di
\file{/proc/sys/kernel} o con l'uso di \func{sysctl}.

\begin{figure}[!htb]
  \footnotesize \centering
  \begin{minipage}[c]{\codesamplewidth}
    \includecodesample{listati/IPCTestId.c}
  \end{minipage} 
  \normalsize 
  \caption{Sezione principale del programma di test per l'assegnazione degli
    identificatori degli oggetti di IPC \file{IPCTestId.c}.}
  \label{fig:ipc_sysv_idtest}
\end{figure}

Per ciascun tipo di oggetto di IPC viene mantenuto in \var{seq} un numero di
sequenza progressivo che viene incrementato di uno ogni volta che l'oggetto
viene cancellato. Quando l'oggetto viene creato usando uno spazio che era già
stato utilizzato in precedenza, per restituire il nuovo identificatore al
numero di oggetti presenti viene sommato il valore corrente del campo
\var{seq}, moltiplicato per il numero massimo di oggetti di quel tipo.

Questo in realtà è quanto avveniva fino ai kernel della serie 2.2, dalla serie
2.4 viene usato lo stesso fattore di moltiplicazione per qualunque tipo di
oggetto, utilizzando il valore dalla costante \constd{IPCMNI} (definita in
\file{include/linux/ipc.h}), che indica il limite massimo complessivo per il
numero di tutti gli oggetti presenti nel \textit{SysV-IPC}, ed il cui default
è 32768.  Si evita così il riutilizzo degli stessi numeri, e si fa sì che
l'identificatore assuma tutti i valori possibili.

In fig.~\ref{fig:ipc_sysv_idtest} è riportato il codice di un semplice
programma di test che si limita a creare un oggetto di IPC (specificato con
una opzione a riga di comando), stamparne il numero di identificatore, e
cancellarlo, il tutto un numero di volte specificato tramite una seconda
opzione.  La figura non riporta il codice di selezione delle opzioni, che
permette di inizializzare i valori delle variabili \var{type} al tipo di
oggetto voluto, e \var{n} al numero di volte che si vuole effettuare il ciclo
di creazione, stampa, cancellazione.

I valori di default sono per l'uso delle code di messaggi e per 5 ripetizioni
del ciclo. Per questo motivo se non si utilizzano opzioni verrà eseguito per
cinque volte il ciclo (\texttt{\small 7-11}), in cui si crea una coda di
messaggi (\texttt{\small 8}), se ne stampa l'identificativo (\texttt{\small
  9}) e la si rimuove (\texttt{\small 10}). Non stiamo ad approfondire adesso
il significato delle funzioni utilizzate, che verranno esaminate nelle
prossime sezioni.

Quello che ci interessa infatti è verificare l'allocazione degli
identificativi associati agli oggetti; lanciando il comando si otterrà
pertanto qualcosa del tipo:
\begin{Console}
piccardi@gont sources]$ \textbf{./ipctestid}
Identifier Value 0 
Identifier Value 32768 
Identifier Value 65536 
Identifier Value 98304 
Identifier Value 131072
\end{Console}
%$
il che ci mostra che stiamo lavorando con un kernel posteriore alla serie 2.2
nel quale non avevamo ancora usato nessuna coda di messaggi (il valore nullo
del primo identificativo indica che il campo \var{seq} era zero). Ripetendo il
comando, e quindi eseguendolo in un processo diverso, in cui non può esistere
nessuna traccia di quanto avvenuto in precedenza, otterremo come nuovo
risultato:
\begin{Console}
[piccardi@gont sources]$ \textbf{./ipctestid}
Identifier Value 163840
Identifier Value 196608 
Identifier Value 229376 
Identifier Value 262144 
Identifier Value 294912 
\end{Console}
%$
in cui la sequenza numerica prosegue, cosa che ci mostra come il valore di
\var{seq} continui ad essere incrementato e costituisca in effetti una
quantità mantenuta all'interno del sistema ed indipendente dai processi.


\subsection{Code di messaggi}
\label{sec:ipc_sysv_mq}

Il primo oggetto introdotto dal \textit{SysV-IPC} è quello delle code di
messaggi.  Le code di messaggi sono oggetti analoghi alle \textit{pipe} o alle
\textit{fifo} ed il loro scopo principale è quello di fornire a processi
diversi un meccanismo con cui scambiarsi dei dati in forma di messaggio. Dato
che le \textit{pipe} e le \textit{fifo} costituiscono una ottima alternativa,
ed in genere sono molto più semplici da usare, le code di messaggi sono il
meno utilizzato degli oggetti introdotti dal \textit{SysV-IPC}.

La funzione di sistema che permette di ottenere l'identificativo di una coda
di messaggi esistente per potervi accedere, oppure di creare una nuova coda
qualora quella indicata non esista ancora, è \funcd{msgget}, e il suo
prototipo è:

\begin{funcproto}{
\fhead{sys/types.h}
\fhead{sys/ipc.h} 
\fhead{sys/msg.h} 
\fdecl{int msgget(key\_t key, int flag)}
\fdesc{Ottiene o crea una coda di messaggi.} 
}

{La funzione ritorna l'identificatore (un intero positivo) in caso di successo
  e $-1$ per un errore, nel qual caso \var{errno} assumerà uno dei valori:
  \begin{errlist}
  \item[\errcode{EACCES}] il processo chiamante non ha i privilegi per
    accedere alla coda richiesta.
  \item[\errcode{EEXIST}] si è richiesta la creazione di una coda che già
    esiste, ma erano specificati sia \const{IPC\_CREAT} che \const{IPC\_EXCL}.
  \item[\errcode{EIDRM}] la coda richiesta è marcata per essere cancellata
    (solo fino al kernel 2.3.20).
  \item[\errcode{ENOENT}] si è cercato di ottenere l'identificatore di una coda
    di messaggi specificando una chiave che non esiste e \const{IPC\_CREAT}
    non era specificato.
  \item[\errcode{ENOSPC}] si è cercato di creare una coda di messaggi quando è
    stato superato il limite massimo di code (\const{MSGMNI}).
 \end{errlist}
 ed inoltre \errval{ENOMEM} nel suo significato generico.}
\end{funcproto}

Le funzione (come le analoghe che si usano per gli altri oggetti) serve sia a
ottenere l'identificatore di una coda di messaggi esistente, che a crearne una
nuova. L'argomento \param{key} specifica la chiave che è associata
all'oggetto, eccetto il caso in cui si specifichi il valore
\const{IPC\_PRIVATE}, nel qual caso la coda è creata ex-novo e non vi è
associata alcuna chiave (per questo viene detta \textsl{privata}), ed il
processo e i suoi eventuali figli potranno farvi riferimento solo attraverso
l'identificatore.

Se invece si specifica un valore diverso da \const{IPC\_PRIVATE} (in Linux
questo significa un valore diverso da zero) l'effetto della funzione dipende
dal valore di \param{flag}, se questo è nullo la funzione si limita ad
effettuare una ricerca sugli oggetti esistenti, restituendo l'identificatore
se trova una corrispondenza, o fallendo con un errore di \errcode{ENOENT} se
non esiste o di \errcode{EACCES} se si sono specificati dei permessi non
validi.

Se invece si vuole creare una nuova coda di messaggi \param{flag} non può
essere nullo e deve essere fornito come maschera binaria, impostando il bit
corrispondente al valore \constd{IPC\_CREAT}. In questo caso i nove bit meno
significativi di \param{flag} saranno usati come permessi per il nuovo
oggetto, secondo quanto illustrato in sez.~\ref{sec:ipc_sysv_access_control}.
Se si imposta anche il bit corrispondente a \constd{IPC\_EXCL} la funzione avrà
successo solo se l'oggetto non esiste già, fallendo con un errore di
\errcode{EEXIST} altrimenti.

Si tenga conto che l'uso di \const{IPC\_PRIVATE} non impedisce ad altri
processi di accedere alla coda, se hanno privilegi sufficienti, una volta che
questi possano indovinare o ricavare, ad esempio per tentativi,
l'identificatore ad essa associato. Per come sono implementati gli oggetti di
IPC infatti non esiste alcun modo in cui si possa garantire l'accesso
esclusivo ad una coda di messaggi.  Usare \const{IPC\_PRIVATE} o
\const{IPC\_CREAT} e \const{IPC\_EXCL} per \param{flag} comporta solo la
creazione di una nuova coda.

\begin{table}[htb]
  \footnotesize
  \centering
  \begin{tabular}[c]{|c|r|l|l|}
    \hline
    \textbf{Costante} & \textbf{Valore} & \textbf{File in \file{/proc}}
    & \textbf{Significato} \\
    \hline
    \hline
    \constd{MSGMNI}&   16& \file{msgmni} & Numero massimo di code di
                                          messaggi.\\
    \constd{MSGMAX}& 8192& \file{msgmax} & Dimensione massima di un singolo
                                          messaggio.\\
    \constd{MSGMNB}&16384& \file{msgmnb} & Dimensione massima del contenuto di 
                                          una coda.\\
    \hline
  \end{tabular}
  \caption{Valori delle costanti associate ai limiti delle code di messaggi.}
  \label{tab:ipc_msg_limits}
\end{table}

Le code di messaggi sono caratterizzate da tre limiti fondamentali, un tempo
definiti staticamente e corrispondenti alle prime tre costanti riportate in
tab.~\ref{tab:ipc_msg_limits}.  Come accennato però con tutte le versioni più
recenti del kernel con Linux è possibile modificare questi limiti attraverso
l'uso di \func{sysctl} o scrivendo nei file \sysctlrelfiled{kernel}{msgmax},
\sysctlrelfiled{kernel}{msgmnb} e \sysctlrelfiled{kernel}{msgmni} di
\file{/proc/sys/kernel/}.

\itindbeg{linked~list}

Una coda di messaggi è costituita da una \textit{linked list}.\footnote{una
  \textit{linked list} è una tipica struttura di dati, organizzati in una
  lista in cui ciascun elemento contiene un puntatore al successivo. In questo
  modo la struttura è veloce nell'estrazione ed immissione dei dati dalle
  estremità dalla lista (basta aggiungere un elemento in testa o in coda ed
  aggiornare un puntatore), e relativamente veloce da attraversare in ordine
  sequenziale (seguendo i puntatori), è invece relativamente lenta
  nell'accesso casuale e nella ricerca.}  I nuovi messaggi vengono inseriti in
coda alla lista e vengono letti dalla cima, in fig.~\ref{fig:ipc_mq_schema} si
è riportato uno schema semplificato con cui queste strutture vengono mantenute
dal kernel. Lo schema illustrato in realtà è una semplificazione di quello
usato fino ai kernel della serie 2.2. A partire della serie 2.4 la gestione
delle code di messaggi è effettuata in maniera diversa (e non esiste una
struttura \kstruct{msqid\_ds} nel kernel), ma abbiamo mantenuto lo schema
precedente dato che illustra in maniera più che adeguata i principi di
funzionamento delle code di messaggi.

\itindend{linked~list}

\begin{figure}[!htb]
  \centering \includegraphics[width=13cm]{img/mqstruct}
  \caption{Schema delle strutture di una coda di messaggi
    (\kstructd{msqid\_ds} e \kstructd{msg}).}
  \label{fig:ipc_mq_schema}
\end{figure}


A ciascuna coda è associata una struttura \kstruct{msqid\_ds} la cui
definizione è riportata in fig.~\ref{fig:ipc_msqid_ds} ed a cui si accede
includendo \headfiled{sys/msg.h};
%
% INFO: sotto materiale obsoleto e non interessante
% In questa struttura il
% kernel mantiene le principali informazioni riguardo lo stato corrente della
% coda. Come accennato questo vale fino ai kernel della serie 2.2, essa viene
% usata nei kernel della serie 2.4 solo per compatibilità in quanto è quella
% restituita dalle funzioni dell'interfaccia; si noti come ci sia una differenza
% con i campi mostrati nello schema di fig.~\ref{fig:ipc_mq_schema} che sono
% presi dalla definizione di \file{include/linux/msg.h}, e fanno riferimento
% alla definizione della omonima struttura usata nel kernel. 
%In fig.~\ref{fig:ipc_msqid_ds} sono elencati i campi definiti in
%\headfile{sys/msg.h};  
si tenga presente che il campo \var{\_\_msg\_cbytes} non è previsto dallo
standard POSIX.1-2001 e che alcuni campi fino al kernel 2.2 erano definiti
come \ctyp{short}.

\begin{figure}[!htb]
  \footnotesize \centering
  \begin{minipage}[c]{.91\textwidth}
    \includestruct{listati/msqid_ds.h}
  \end{minipage} 
  \normalsize 
  \caption{La struttura \structd{msqid\_ds}, associata a ciascuna coda di
    messaggi.}
  \label{fig:ipc_msqid_ds}
\end{figure}

Quando si crea una nuova coda con \func{msgget} questa struttura viene
inizializzata,\footnote{in realtà viene inizializzata una struttura interna al
  kernel, ma i dati citati sono gli stessi.} in particolare il campo
\var{msg\_perm} che esprime i permessi di accesso viene inizializzato nella
modalità illustrata in sez.~\ref{sec:ipc_sysv_access_control}. Per quanto
riguarda gli altri campi invece:
\begin{itemize*}
\item il campo \var{msg\_qnum}, che esprime il numero di messaggi presenti
  sulla coda, viene inizializzato a 0.
\item i campi \var{msg\_lspid} e \var{msg\_lrpid}, che esprimono
  rispettivamente il \ids{PID} dell'ultimo processo che ha inviato o ricevuto
  un messaggio sulla coda, sono inizializzati a 0.
\item i campi \var{msg\_stime} e \var{msg\_rtime}, che esprimono
  rispettivamente il tempo in cui è stato inviato o ricevuto l'ultimo
  messaggio sulla coda, sono inizializzati a 0.
\item il campo \var{msg\_ctime}, che esprime il tempo di ultima modifica della
  coda, viene inizializzato al tempo corrente.
\item il campo \var{msg\_qbytes}, che esprime la dimensione massima del
  contenuto della coda (in byte) viene inizializzato al valore preimpostato
  del sistema (\const{MSGMNB}).
\item il campo \var{\_\_msg\_cbytes}, che esprime la dimensione in byte dei
  messaggi presenti sulla coda, viene inizializzato a zero.
% i campi \var{msg\_first} e \var{msg\_last} che esprimono l'indirizzo del
%   primo e ultimo messaggio sono inizializzati a \val{NULL} e
%   \var{msg\_cbytes}, che esprime la dimensione in byte dei messaggi presenti è
%   inizializzato a zero. Questi campi sono ad uso interno dell'implementazione
%   e non devono essere utilizzati da programmi in \textit{user space}).
\end{itemize*}

Una volta creata una coda di messaggi le operazioni di controllo vengono
effettuate con la funzione di sistema \funcd{msgctl}, che, come le analoghe
\func{semctl} e \func{shmctl}, fa le veci di quello che \func{ioctl} è per i
file; il suo prototipo è:

\begin{funcproto}{
\fhead{sys/types.h}
\fhead{sys/ipc.h}
\fhead{sys/msg.h}
\fdecl{int msgctl(int msqid, int cmd, struct msqid\_ds *buf)}
\fdesc{Esegue una operazione su una coda.} 
}

{La funzione ritorna $0$ in caso di successo e $-1$ per un errore, nel qual
  caso \var{errno} assumerà uno dei valori: 
  \begin{errlist}
  \item[\errcode{EACCES}] si è richiesto \const{IPC\_STAT} ma processo
    chiamante non ha i privilegi di lettura sulla coda.
  \item[\errcode{EIDRM}] la coda richiesta è stata cancellata.
  \item[\errcode{EPERM}] si è richiesto \const{IPC\_SET} o \const{IPC\_RMID} ma
    il processo non ha i privilegi, o si è richiesto di aumentare il valore di
    \var{msg\_qbytes} oltre il limite \const{MSGMNB} senza essere
    amministratore.
  \end{errlist}
  ed inoltre \errval{EFAULT} ed \errval{EINVAL} nel loro significato
  generico.}
\end{funcproto}

La funzione permette di eseguire una operazione di controllo per la coda
specificata dall'identificatore \param{msqid}, utilizzando i valori della
struttura \struct{msqid\_ds}, mantenuta all'indirizzo \param{buf}. Il
comportamento della funzione dipende dal valore dell'argomento \param{cmd},
che specifica il tipo di azione da eseguire. I valori possibili
per \param{cmd} sono:
\begin{basedescript}{\desclabelwidth{1.6cm}\desclabelstyle{\nextlinelabel}}
\item[\constd{IPC\_STAT}] Legge le informazioni riguardo la coda nella
  struttura \struct{msqid\_ds} indicata da \param{buf}. Occorre avere il
  permesso di lettura sulla coda.
\item[\constd{IPC\_RMID}] Rimuove la coda, cancellando tutti i dati, con
  effetto immediato. Tutti i processi che cercheranno di accedere alla coda
  riceveranno un errore di \errcode{EIDRM}, e tutti processi in attesa su
  funzioni di lettura o di scrittura sulla coda saranno svegliati ricevendo
  il medesimo errore. Questo comando può essere eseguito solo da un processo
  con \ids{UID} effettivo corrispondente al creatore o al proprietario della
  coda, o all'amministratore.
\item[\constd{IPC\_SET}] Permette di modificare i permessi ed il proprietario
  della coda, ed il limite massimo sulle dimensioni del totale dei messaggi in
  essa contenuti (\var{msg\_qbytes}). I valori devono essere passati in una
  struttura \struct{msqid\_ds} puntata da \param{buf}.  Per modificare i
  valori di \var{msg\_perm.mode}, \var{msg\_perm.uid} e \var{msg\_perm.gid}
  occorre essere il proprietario o il creatore della coda, oppure
  l'amministratore e lo stesso vale per \var{msg\_qbytes}. Infine solo
  l'amministratore (più precisamente un processo con la capacità
  \const{CAP\_IPC\_RESOURCE}) ha la facoltà di incrementarne il valore a
  limiti superiori a \const{MSGMNB}. Se eseguita con successo la funzione
  aggiorna anche il campo \var{msg\_ctime}.
\end{basedescript}

A questi tre valori, che sono quelli previsti dallo standard, su Linux se ne
affiancano altri tre (\constd{IPC\_INFO}, \constd{MSG\_STAT} e
\constd{MSG\_INFO}) introdotti ad uso del programma \cmd{ipcs} per ottenere le
informazioni generali relative alle risorse usate dalle code di
messaggi. Questi potranno essere modificati o rimossi in favore dell'uso di
\texttt{/proc}, per cui non devono essere usati e non li tratteremo.

Una volta che si abbia a disposizione l'identificatore, per inviare un
messaggio su una coda si utilizza la funzione di sistema \funcd{msgsnd}, il
cui prototipo è:

\begin{funcproto}{
\fhead{sys/types.h}
\fhead{sys/ipc.h}
\fhead{sys/msg.h}
\fdecl{int msgsnd(int msqid, struct msgbuf *msgp, size\_t msgsz, int msgflg)}
\fdesc{Invia un messaggio su una coda.}
}

{La funzione ritorna $0$ in caso di successo e $-1$ per un errore, nel qual
  caso \var{errno} assumerà uno dei valori: 
  \begin{errlist}
  \item[\errcode{EACCES}] non si hanno i privilegi di accesso sulla coda.
  \item[\errcode{EAGAIN}] il messaggio non può essere inviato perché si è
    superato il limite \var{msg\_qbytes} sul numero massimo di byte presenti
    sulla coda, e si è richiesto \const{IPC\_NOWAIT} in \param{flag}.
  \item[\errcode{EIDRM}] la coda è stata cancellata.
  \item[\errcode{EINVAL}] si è specificato un \param{msgid} invalido, o un
    valore non positivo per \param{mtype}, o un valore di \param{msgsz}
    maggiore di \const{MSGMAX}.
  \end{errlist}
  ed inoltre \errval{EFAULT}, \errval{EINTR} e \errval{ENOMEM} nel loro
  significato generico.}
\end{funcproto}

La funzione inserisce il messaggio sulla coda specificata da \param{msqid}; il
messaggio ha lunghezza specificata da \param{msgsz} ed è passato attraverso il
l'argomento \param{msgp}.  Quest'ultimo deve venire passato sempre come
puntatore ad una struttura \struct{msgbuf} analoga a quella riportata in
fig.~\ref{fig:ipc_msbuf} che è quella che deve contenere effettivamente il
messaggio.  La dimensione massima per il testo di un messaggio non può
comunque superare il limite \const{MSGMAX}.

\begin{figure}[!htb]
  \footnotesize \centering
  \begin{minipage}[c]{0.8\textwidth}
    \includestruct{listati/msgbuf.h}
  \end{minipage} 
  \normalsize 
  \caption{Schema della struttura \structd{msgbuf}, da utilizzare come
    argomento per inviare/ricevere messaggi.}
  \label{fig:ipc_msbuf}
\end{figure}

La struttura di fig.~\ref{fig:ipc_msbuf} è comunque solo un modello, tanto che
la definizione contenuta in \headfile{sys/msg.h} usa esplicitamente per il
secondo campo il valore \code{mtext[1]}, che non è di nessuna utilità ai fini
pratici.  La sola cosa che conta è che la struttura abbia come primo membro un
campo \var{mtype} come nell'esempio; esso infatti serve ad identificare il
tipo di messaggio e deve essere sempre specificato come intero positivo di
tipo \ctyp{long}.  Il campo \var{mtext} invece può essere di qualsiasi tipo e
dimensione, e serve a contenere il testo del messaggio.

In generale pertanto per inviare un messaggio con \func{msgsnd} si usa
ridefinire una struttura simile a quella di fig.~\ref{fig:ipc_msbuf}, adattando
alle proprie esigenze il campo \var{mtype}, (o ridefinendo come si vuole il
corpo del messaggio, anche con più campi o con strutture più complesse) avendo
però la cura di mantenere nel primo campo un valore di tipo \ctyp{long} che ne
indica il tipo.

Si tenga presente che la lunghezza che deve essere indicata in questo
argomento è solo quella del messaggio, non quella di tutta la struttura, se
cioè \var{message} è una propria struttura che si passa alla funzione,
\param{msgsz} dovrà essere uguale a \code{sizeof(message)-sizeof(long)}, (se
consideriamo il caso dell'esempio in fig.~\ref{fig:ipc_msbuf}, \param{msgsz}
dovrà essere pari a \var{LENGTH}).

Per capire meglio il funzionamento della funzione riprendiamo in
considerazione la struttura della coda illustrata in
fig.~\ref{fig:ipc_mq_schema}. Alla chiamata di \func{msgsnd} il nuovo
messaggio sarà aggiunto in fondo alla lista inserendo una nuova struttura
\kstruct{msg}, il puntatore \var{msg\_last} di \kstruct{msqid\_ds} verrà
aggiornato, come pure il puntatore al messaggio successivo per quello che era
il precedente ultimo messaggio; il valore di \var{mtype} verrà mantenuto in
\var{msg\_type} ed il valore di \param{msgsz} in \var{msg\_ts}; il testo del
messaggio sarà copiato all'indirizzo specificato da \var{msg\_spot}.

Il valore dell'argomento \param{flag} permette di specificare il comportamento
della funzione. Di norma, quando si specifica un valore nullo, la funzione
ritorna immediatamente a meno che si sia ecceduto il valore di
\var{msg\_qbytes}, o il limite di sistema sul numero di messaggi, nel qual
caso si blocca.  Se si specifica per \param{flag} il valore
\constd{IPC\_NOWAIT} la funzione opera in modalità non-bloccante, ed in questi
casi ritorna immediatamente con un errore di \errcode{EAGAIN}.

Se non si specifica \const{IPC\_NOWAIT} la funzione resterà bloccata fintanto
che non si liberano risorse sufficienti per poter inserire nella coda il
messaggio, nel qual caso ritornerà normalmente. La funzione può ritornare con
una condizione di errore anche in due altri casi: quando la coda viene rimossa
(nel qual caso si ha un errore di \errcode{EIDRM}) o quando la funzione viene
interrotta da un segnale (nel qual caso si ha un errore di \errcode{EINTR}).

Una volta completato con successo l'invio del messaggio sulla coda, la
funzione aggiorna i dati mantenuti in \struct{msqid\_ds}, in particolare
vengono modificati:
\begin{itemize*}
\item Il valore di \var{msg\_lspid}, che viene impostato al \ids{PID} del
  processo chiamante.
\item Il valore di \var{msg\_qnum}, che viene incrementato di uno.
\item Il valore \var{msg\_stime}, che viene impostato al tempo corrente.
\end{itemize*}

La funzione di sistema che viene utilizzata per estrarre un messaggio da una
coda è \funcd{msgrcv}, ed il suo prototipo è:

\begin{funcproto}{
\fhead{sys/types.h}
\fhead{sys/ipc.h} 
\fhead{sys/msg.h}
\fdecl{ssize\_t msgrcv(int msqid, struct msgbuf *msgp, size\_t msgsz,
    long msgtyp, int msgflg)}
\fdesc{Legge un messaggio da una coda.} 
}

{La funzione ritorna il numero di byte letti in caso di successo e $-1$ per un
  errore, nel qual caso \var{errno} assumerà uno dei valori:
  \begin{errlist}
  \item[\errcode{E2BIG}] il testo del messaggio è più lungo di \param{msgsz} e
    non si è specificato \const{MSG\_NOERROR} in \param{msgflg}.
  \item[\errcode{EACCES}] non si hanno i privilegi di accesso sulla coda.
  \item[\errcode{EIDRM}] la coda è stata cancellata.
  \item[\errcode{EINTR}] la funzione è stata interrotta da un segnale mentre
    era in attesa di ricevere un messaggio.
  \item[\errcode{EINVAL}] si è specificato un \param{msgid} invalido o un
    valore di \param{msgsz} negativo.
  \end{errlist}
  ed inoltre \errval{EFAULT} nel suo significato generico.}
\end{funcproto}

La funzione legge un messaggio dalla coda specificata da \param{msqid},
scrivendolo sulla struttura puntata da \param{msgp}, che dovrà avere un
formato analogo a quello di fig.~\ref{fig:ipc_msbuf}.  Una volta estratto, il
messaggio sarà rimosso dalla coda.  L'argomento \param{msgsz} indica la
lunghezza massima del testo del messaggio (equivalente al valore del parametro
\var{LENGTH} nell'esempio di fig.~\ref{fig:ipc_msbuf}).

Se il testo del messaggio ha lunghezza inferiore a \param{msgsz} esso viene
rimosso dalla coda; in caso contrario, se \param{msgflg} è impostato a
\constd{MSG\_NOERROR}, il messaggio viene troncato e la parte in eccesso viene
perduta, altrimenti il messaggio non viene estratto e la funzione ritorna con
un errore di \errcode{E2BIG}.

L'argomento \param{msgtyp} permette di restringere la ricerca ad un
sottoinsieme dei messaggi presenti sulla coda; la ricerca infatti è fatta con
una scansione della struttura mostrata in fig.~\ref{fig:ipc_mq_schema},
restituendo il primo messaggio incontrato che corrisponde ai criteri
specificati (che quindi, visto come i messaggi vengono sempre inseriti dalla
coda, è quello meno recente); in particolare:
\begin{itemize*}
\item se \param{msgtyp} è 0 viene estratto il messaggio in cima alla coda, cioè
  quello fra i presenti che è stato inserito per primo. 
\item se \param{msgtyp} è positivo viene estratto il primo messaggio il cui
  tipo (il valore del campo \var{mtype}) corrisponde al valore di
  \param{msgtyp}.
\item se \param{msgtyp} è negativo viene estratto il primo fra i messaggi con
  il valore più basso del tipo, fra tutti quelli il cui tipo ha un valore
  inferiore al valore assoluto di \param{msgtyp}.
\end{itemize*}

Il valore di \param{msgflg} permette di controllare il comportamento della
funzione, esso può essere nullo o una maschera binaria composta da uno o più
valori.  Oltre al precedente \const{MSG\_NOERROR}, sono possibili altri due
valori: \constd{MSG\_EXCEPT}, che permette, quando \param{msgtyp} è positivo,
di leggere il primo messaggio nella coda con tipo diverso da \param{msgtyp}, e
\const{IPC\_NOWAIT} che causa il ritorno immediato della funzione quando non
ci sono messaggi sulla coda.

Il comportamento usuale della funzione infatti, se non ci sono messaggi
disponibili per la lettura, è di bloccare il processo. Nel caso però si sia
specificato \const{IPC\_NOWAIT} la funzione ritorna immediatamente con un
errore \errcode{ENOMSG}. Altrimenti la funzione ritorna normalmente non appena
viene inserito un messaggio del tipo desiderato, oppure ritorna con errore
qualora la coda sia rimossa (con \var{errno} impostata a \errcode{EIDRM}) o se
il processo viene interrotto da un segnale (con \var{errno} impostata a
\errcode{EINTR}).

Una volta completata con successo l'estrazione del messaggio dalla coda, la
funzione aggiorna i dati mantenuti in \struct{msqid\_ds}, in particolare
vengono modificati:
\begin{itemize*}
\item Il valore di \var{msg\_lrpid}, che viene impostato al \ids{PID} del
  processo chiamante.
\item Il valore di \var{msg\_qnum}, che viene decrementato di uno.
\item Il valore \var{msg\_rtime}, che viene impostato al tempo corrente.
\end{itemize*}

Le code di messaggi presentano il solito problema di tutti gli oggetti del
\textit{SysV-IPC} che essendo questi permanenti restano nel sistema occupando
risorse anche quando un processo è terminato, al contrario delle \textit{pipe}
per le quali tutte le risorse occupate vengono rilasciate quanto l'ultimo
processo che le utilizzava termina. Questo comporta che in caso di errori si
può saturare il sistema, e che devono comunque essere esplicitamente previste
delle funzioni di rimozione in caso di interruzioni o uscite dal programma
(come vedremo in fig.~\ref{fig:ipc_mq_fortune_server}).

L'altro problema è che non facendo uso di file descriptor le tecniche di
\textit{I/O multiplexing} descritte in sez.~\ref{sec:file_multiplexing} non
possono essere utilizzate, e non si ha a disposizione niente di analogo alle
funzioni \func{select} e \func{poll}. Questo rende molto scomodo usare più di
una di queste strutture alla volta; ad esempio non si può scrivere un server
che aspetti un messaggio su più di una coda senza fare ricorso ad una tecnica
di \textit{polling} che esegua un ciclo di attesa su ciascuna di esse.

Come esempio dell'uso delle code di messaggi possiamo riscrivere il nostro
server di \textit{fortunes} usando queste al posto delle \textit{fifo}. In
questo caso useremo una sola coda di messaggi, usando il tipo di messaggio per
comunicare in maniera indipendente con client diversi.

\begin{figure}[!htbp]
  \footnotesize \centering
  \begin{minipage}[c]{\codesamplewidth}
    \includecodesample{listati/MQFortuneServer.c}
  \end{minipage} 
  \normalsize 
  \caption{Sezione principale del codice del server di \textit{fortunes}
    basato sulle \textit{message queue}.}
  \label{fig:ipc_mq_fortune_server}
\end{figure}

In fig.~\ref{fig:ipc_mq_fortune_server} si è riportato un estratto delle parti
principali del codice del nuovo server (il codice completo è nel file
\file{MQFortuneServer.c} nei sorgenti allegati). Il programma è basato su un
uso accorto della caratteristica di poter associate un ``tipo'' ai messaggi
per permettere una comunicazione indipendente fra il server ed i vari client,
usando il \ids{PID} di questi ultimi come identificativo. Questo è possibile
in quanto, al contrario di una \textit{fifo}, la lettura di una coda di
messaggi può non essere sequenziale, proprio grazie alla classificazione dei
messaggi sulla base del loro tipo.

Il programma, oltre alle solite variabili per il nome del file da cui leggere
le \textit{fortunes} e per il vettore di stringhe che contiene le frasi,
definisce due strutture appositamente per la comunicazione; con
\var{msgbuf\_read} vengono passate (\texttt{\small 8-11}) le richieste mentre
con \var{msgbuf\_write} vengono restituite (\texttt{\small 12-15}) le frasi.

La gestione delle opzioni si è al solito omessa, essa si curerà di impostare
nella variabile \var{n} il numero di frasi da leggere specificato a linea di
comando ed in \var{fortunefilename} il file da cui leggerle. Dopo aver
installato (\texttt{\small 19-21}) i gestori dei segnali per trattare
l'uscita dal server, viene prima controllato (\texttt{\small 22}) il numero di
frasi richieste abbia senso (cioè sia maggiore di zero), le quali poi vengono
lette (\texttt{\small 23}) nel vettore in memoria con la stessa funzione
\code{FortuneParse} usata anche per il server basato sulle \textit{fifo}.

Una volta inizializzato il vettore di stringhe coi messaggi presi dal file
delle \textit{fortune} si procede (\texttt{\small 25}) con la generazione di
una chiave per identificare la coda di messaggi (si usa il nome del file dei
sorgenti del server) con la quale poi si esegue (\texttt{\small 26}) la
creazione della stessa (si noti come si sia chiamata \func{msgget} con un
valore opportuno per l'argomento \param{flag}), avendo cura di abortire il
programma (\texttt{\small 27-29}) in caso di errore.

Finita la fase di inizializzazione il server prima (\texttt{\small 32}) chiama
la funzione \func{daemon} per andare in background e poi esegue in permanenza
il ciclo principale (\texttt{\small 33-40}). Questo inizia (\texttt{\small
  34}) con il porsi in attesa di un messaggio di richiesta da parte di un
client. Si noti infatti come \func{msgrcv} richieda un messaggio con
\var{mtype} uguale a 1, questo è il valore usato per le richieste dato che
corrisponde al \ids{PID} di \cmd{init}, che non può essere un client. L'uso
del flag \const{MSG\_NOERROR} è solo per sicurezza, dato che i messaggi di
richiesta sono di dimensione fissa (e contengono solo il \ids{PID} del
client).

Se non sono presenti messaggi di richiesta \func{msgrcv} si bloccherà,
ritornando soltanto in corrispondenza dell'arrivo sulla coda di un messaggio
di richiesta da parte di un client, in tal caso il ciclo prosegue
(\texttt{\small 35}) selezionando una frase a caso, copiandola (\texttt{\small
  36}) nella struttura \var{msgbuf\_write} usata per la risposta e
calcolandone (\texttt{\small 37}) la dimensione.

Per poter permettere a ciascun client di ricevere solo la risposta indirizzata
a lui il tipo del messaggio in uscita viene inizializzato (\texttt{\small 38})
al valore del \ids{PID} del client ricevuto nel messaggio di richiesta.
L'ultimo passo del ciclo (\texttt{\small 39}) è inviare sulla coda il
messaggio di risposta. Si tenga conto che se la coda è piena anche questa
funzione potrà bloccarsi fintanto che non venga liberato dello spazio.

Si noti che il programma può terminare solo grazie ad una interruzione da
parte di un segnale; in tal caso verrà eseguito (\texttt{\small 45-48}) il
gestore \code{HandSIGTERM}, che semplicemente si limita a cancellare la coda
(\texttt{\small 46}) ed ad uscire (\texttt{\small 47}).

\begin{figure}[!htbp]
  \footnotesize \centering
  \begin{minipage}[c]{\codesamplewidth}
    \includecodesample{listati/MQFortuneClient.c}
  \end{minipage} 
  \normalsize 
  \caption{Sezione principale del codice del client di \textit{fortunes}
    basato sulle \textit{message queue}.}
  \label{fig:ipc_mq_fortune_client}
\end{figure}

In fig.~\ref{fig:ipc_mq_fortune_client} si è riportato un estratto il codice
del programma client.  Al solito il codice completo è con i sorgenti allegati,
nel file \file{MQFortuneClient.c}.  Come sempre si sono rimosse le parti
relative alla gestione delle opzioni, ed in questo caso, anche la
dichiarazione delle variabili, che, per la parte relative alle strutture usate
per la comunicazione tramite le code, sono le stesse viste in
fig.~\ref{fig:ipc_mq_fortune_server}.

Il client in questo caso è molto semplice; la prima parte del programma
(\texttt{\small 4-9}) si occupa di accedere alla coda di messaggi, ed è
identica a quanto visto per il server, solo che in questo caso \func{msgget}
non viene chiamata con il flag di creazione in quanto la coda deve essere
preesistente. In caso di errore (ad esempio se il server non è stato avviato)
il programma termina immediatamente. 

Una volta acquisito l'identificatore della coda il client compone
(\texttt{\small 12-13}) il messaggio di richiesta in \var{msg\_read}, usando
1 per il tipo ed inserendo il proprio \ids{PID} come dato da passare al
server.  Calcolata (\texttt{\small 14}) la dimensione, provvede
(\texttt{\small 15}) ad immettere la richiesta sulla coda.

A questo punto non resta che rileggere la risposta (\texttt{\small 16}) dalla
coda del server richiedendo a \func{msgrcv} di selezionare i messaggi di tipo
corrispondente al valore del \ids{PID} inviato nella richiesta. L'ultimo passo
(\texttt{\small 17}) prima di uscire è quello di stampare a video il messaggio
ricevuto.
 
Proviamo allora il nostro nuovo sistema, al solito occorre definire
\code{LD\_LIBRARY\_PATH} per accedere alla libreria \file{libgapil.so}, dopo
di che, in maniera del tutto analoga a quanto fatto con il programma che usa
le \textit{fifo}, potremo far partire il server con:
\begin{Console}
[piccardi@gont sources]$ \textbf{./mqfortuned -n10}
\end{Console}
%$
come nel caso precedente, avendo eseguito il server in background, il comando
ritornerà immediatamente; potremo però verificare con \cmd{ps} che il
programma è effettivamente in esecuzione, e che ha creato una coda di
messaggi:
\begin{Console}
[piccardi@gont sources]$ \textbf{ipcs}

------ Shared Memory Segments --------
key        shmid      owner      perms      bytes      nattch     status      

------ Semaphore Arrays --------
key        semid      owner      perms      nsems     

------ Message Queues --------
key        msqid      owner      perms      used-bytes   messages    
0x0102dc6a 0          piccardi   666        0            0           
\end{Console}
%$
a questo punto potremo usare il client per ottenere le nostre frasi:
\begin{Console}
[piccardi@gont sources]$ \textbf{./mqfortune}
Linux ext2fs has been stable for a long time, now it's time to break it
        -- Linuxkongreß '95 in Berlin
[piccardi@gont sources]$ \textbf{./mqfortune}
Let's call it an accidental feature.
        --Larry Wall
\end{Console} 
con un risultato del tutto equivalente al precedente. Infine potremo chiudere
il server inviando il segnale di terminazione con il comando \code{killall
  mqfortuned}, verificando che effettivamente la coda di messaggi venga
rimossa.

Benché funzionante questa architettura risente dello stesso inconveniente
visto anche nel caso del precedente server basato sulle \textit{fifo}; se il
client viene interrotto dopo l'invio del messaggio di richiesta e prima della
lettura della risposta, quest'ultima resta nella coda (così come per le
\textit{fifo} si aveva il problema delle \textit{fifo} che restavano nel
filesystem). In questo caso però il problemi sono maggiori, sia perché è molto
più facile esaurire la memoria dedicata ad una coda di messaggi che gli
\textit{inode} di un filesystem, sia perché, con il riutilizzo dei \ids{PID}
da parte dei processi, un client eseguito in un momento successivo potrebbe
ricevere un messaggio non indirizzato a lui.


\subsection{I semafori}
\label{sec:ipc_sysv_sem}

I semafori non sono propriamente meccanismi di intercomunicazione come
\textit{pipe}, \textit{fifo} e code di messaggi, poiché non consentono di
scambiare dati fra processi, ma servono piuttosto come meccanismi di
sincronizzazione o di protezione per le \textsl{sezioni critiche} del codice
(si ricordi quanto detto in sez.~\ref{sec:proc_race_cond}).  Un semaforo
infatti non è altro che un contatore mantenuto nel kernel che determina se
consentire o meno la prosecuzione dell'esecuzione di un programma. In questo
modo si può controllare l'accesso ad una risorsa condivisa da più processi,
associandovi un semaforo che assicuri che non possa essere usata da più di un
processo alla volta.

Il concetto di semaforo è uno dei concetti base nella programmazione ed è
assolutamente generico, così come del tutto generali sono modalità con cui lo
si utilizza. Un processo che deve accedere ad una risorsa condivisa eseguirà
un controllo del semaforo: se questo è positivo il suo valore sarà
decrementato, indicando che si è consumato una unità della risorsa, ed il
processo potrà proseguire nell'utilizzo di quest'ultima, provvedendo a
rilasciarla, una volta completate le operazioni volute, reincrementando il
semaforo.

Se al momento del controllo il valore del semaforo è nullo la risorsa viene
considerata non disponibile, ed il processo si bloccherà fin quando chi la sta
utilizzando non la rilascerà, incrementando il valore del semaforo. Non appena
il semaforo diventa positivo, indicando che la risorsa è tornata disponibile,
il processo bloccato in attesa riprenderà l'esecuzione, e potrà operare come
nel caso precedente (decremento del semaforo, accesso alla risorsa, incremento
del semaforo).

Per poter implementare questo tipo di logica le operazioni di controllo e
decremento del contatore associato al semaforo devono essere atomiche,
pertanto una realizzazione di un oggetto di questo tipo è necessariamente
demandata al kernel. La forma più semplice di semaforo è quella del
\textsl{semaforo binario}, o \textit{mutex}, in cui un valore diverso da zero
(normalmente 1) indica la libertà di accesso, e un valore nullo l'occupazione
della risorsa. In generale però si possono usare semafori con valori interi,
utilizzando il valore del contatore come indicatore del ``numero di risorse''
ancora disponibili.

Il sistema di intercomunicazione di \textit{SysV-IPC} prevede anche una
implementazione dei semafori, ma gli oggetti utilizzati sono tuttavia non
semafori singoli, ma gruppi (più propriamente \textsl{insiemi}) di semafori
detti ``\textit{semaphore set}''. La funzione di sistema che permette di
creare o ottenere l'identificatore di un insieme di semafori è \funcd{semget},
ed il suo prototipo è:

\begin{funcproto}{
\fhead{sys/types.h}
\fhead{sys/ipc.h}
\fhead{sys/sem.h}
\fdecl{int semget(key\_t key, int nsems, int flag)}
\fdesc{Restituisce l'identificatore di un insieme di semafori.} 
}

{La funzione ritorna l'identificatore (un intero positivo) in caso di successo
  e $-1$ per un errore, nel qual caso \var{errno} assumerà uno dei valori:
  \begin{errlist}
  \item[\errcode{ENOSPC}] si è superato il limite di sistema per il numero
    totale di semafori (\const{SEMMNS}) o di insiemi (\const{SEMMNI}).
  \item[\errcode{EINVAL}] \param{nsems} è minore di zero o maggiore del limite
    sul numero di semafori di un insieme (\const{SEMMSL}), o se l'insieme già
    esiste, maggiore del numero di semafori che contiene.
  \item[\errcode{ENOMEM}] il sistema non ha abbastanza memoria per poter
    contenere le strutture per un nuovo insieme di semafori.
  \end{errlist}
  ed inoltre \errval{EACCES}, \errval{EEXIST}, \errval{EIDRM} e
  \errval{ENOENT} con lo stesso significato che hanno per \func{msgget}.}
\end{funcproto}

La funzione è del tutto analoga a \func{msgget}, solo che in questo caso
restituisce l'identificatore di un insieme di semafori, in particolare è
identico l'uso degli argomenti \param{key} e \param{flag}, per cui non
ripeteremo quanto detto al proposito in sez.~\ref{sec:ipc_sysv_mq}. L'argomento
\param{nsems} permette di specificare quanti semafori deve contenere l'insieme
quando se ne richieda la creazione, e deve essere nullo quando si effettua una
richiesta dell'identificatore di un insieme già esistente.

Purtroppo questa implementazione complica inutilmente lo schema elementare che
abbiamo descritto, dato che non è possibile definire un singolo semaforo, ma
se ne deve creare per forza un insieme.  Ma questa in definitiva è solo una
complicazione inutile dell'interfaccia, il problema è che i semafori forniti
dal \textit{SysV-IPC} soffrono di altri due difetti progettuali molto più
gravi.

Il primo difetto è che non esiste una funzione che permetta di creare ed
inizializzare un semaforo in un'unica chiamata; occorre prima creare l'insieme
dei semafori con \func{semget} e poi inizializzarlo con \func{semctl}, si
perde così ogni possibilità di eseguire l'operazione atomicamente. Eventuali
accessi che possono avvenire fra la creazione e l'inizializzazione potranno
avere effetti imprevisti.

Il secondo difetto deriva dalla caratteristica generale degli oggetti del
\textit{SysV-IPC} di essere risorse globali di sistema, che non vengono
cancellate quando nessuno le usa più. In questo caso il problema è più grave
perché ci si a trova a dover affrontare esplicitamente il caso in cui un
processo termina per un qualche errore lasciando un semaforo occupato, che
resterà tale fino al successivo riavvio del sistema. Come vedremo esistono
delle modalità per evitare tutto ciò, ma diventa necessario indicare
esplicitamente che si vuole il ripristino del semaforo all'uscita del
processo, e la gestione diventa più complicata.

\begin{figure}[!htb]
  \footnotesize \centering
  \begin{minipage}[c]{.85\textwidth}
    \includestruct{listati/semid_ds.h}
  \end{minipage} 
  \normalsize 
  \caption{La struttura \structd{semid\_ds}, associata a ciascun insieme di
    semafori.}
  \label{fig:ipc_semid_ds}
\end{figure}

A ciascun insieme di semafori è associata una struttura \struct{semid\_ds},
riportata in fig.~\ref{fig:ipc_semid_ds}.\footnote{anche in questo caso in
  realtà il kernel usa una sua specifica struttura interna, ma i dati
  significativi sono sempre quelli citati.}  Come nel caso delle code di
messaggi quando si crea un nuovo insieme di semafori con \func{semget} questa
struttura viene inizializzata. In particolare il campo \var{sem\_perm}, che
esprime i permessi di accesso, viene inizializzato come illustrato in
sez.~\ref{sec:ipc_sysv_access_control} (si ricordi che in questo caso il
permesso di scrittura è in realtà permesso di alterare il semaforo), per
quanto riguarda gli altri campi invece:
\begin{itemize*}
\item il campo \var{sem\_nsems}, che esprime il numero di semafori
  nell'insieme, viene inizializzato al valore di \param{nsems}.
\item il campo \var{sem\_ctime}, che esprime il tempo di ultimo cambiamento
  dell'insieme, viene inizializzato al tempo corrente.
\item il campo \var{sem\_otime}, che esprime il tempo dell'ultima operazione
  effettuata, viene inizializzato a zero.
\end{itemize*}

\begin{figure}[!htb]
  \footnotesize \centering
  \begin{minipage}[c]{.85\textwidth}
    \includestruct{listati/sem.h}
  \end{minipage} 
  \normalsize 
  \caption{La struttura \structd{sem}, che contiene i dati di un singolo
    semaforo.} 
  \label{fig:ipc_sem}
\end{figure}

Ciascun semaforo dell'insieme è realizzato come una struttura di tipo
\struct{sem} che ne contiene i dati essenziali, la cui definizione è riportata
in fig.~\ref{fig:ipc_sem}.\footnote{in realtà in fig~\ref{fig:ipc_sem} si è
  riportata la definizione originaria del kernel 1.0, che contiene la prima
  realizzazione del \textit{SysV-IPC} in Linux; ormai questa struttura è
  ridotta ai soli due primi membri, e gli altri vengono calcolati
  dinamicamente, la si è usata solo a scopo di esempio, perché indica tutti i
  valori associati ad un semaforo, restituiti dalle funzioni di controllo, e
  citati dalle pagine di manuale.}  Questa struttura non è accessibile
direttamente dallo \textit{user space}, ma i valori in essa specificati
possono essere letti in maniera indiretta, attraverso l'uso delle opportune
funzioni di controllo.  I dati mantenuti nella struttura, ed elencati in
fig.~\ref{fig:ipc_sem}, indicano rispettivamente:
\begin{basedescript}{\desclabelwidth{2cm}\desclabelstyle{\nextlinelabel}}
\item[\var{semval}] il valore numerico del semaforo.
\item[\var{sempid}] il \ids{PID} dell'ultimo processo che ha eseguito una
  operazione sul semaforo.
\item[\var{semncnt}] il numero di processi in attesa che esso venga
  incrementato.
\item[\var{semzcnt}] il numero di processi in attesa che esso si annulli.
\end{basedescript}

Come per le code di messaggi anche per gli insiemi di semafori esistono una
serie di limiti, i cui valori sono associati ad altrettante costanti, che si
sono riportate in tab.~\ref{tab:ipc_sem_limits}. Alcuni di questi limiti sono
al solito accessibili e modificabili attraverso \func{sysctl} o scrivendo
direttamente nel file \sysctlfile{kernel/sem}.

\begin{table}[htb]
  \footnotesize
  \centering
  \begin{tabular}[c]{|c|r|p{8cm}|}
    \hline
    \textbf{Costante} & \textbf{Valore} & \textbf{Significato} \\
    \hline
    \hline
    \constd{SEMMNI}&          128 & Numero massimo di insiemi di semafori.\\
    \constd{SEMMSL}&          250 & Numero massimo di semafori per insieme.\\
    \constd{SEMMNS}&\const{SEMMNI}*\const{SEMMSL}& Numero massimo di semafori
                                                   nel sistema.\\
    \constd{SEMVMX}&        32767 & Massimo valore per un semaforo.\\
    \constd{SEMOPM}&           32 & Massimo numero di operazioni per chiamata a
                                    \func{semop}. \\
    \constd{SEMMNU}&\const{SEMMNS}& Massimo numero di strutture di ripristino.\\
    \constd{SEMUME}&\const{SEMOPM}& Massimo numero di voci di ripristino.\\
    \constd{SEMAEM}&\const{SEMVMX}& Valore massimo per l'aggiustamento
                                    all'uscita. \\
    \hline
  \end{tabular}
  \caption{Valori delle costanti associate ai limiti degli insiemi di
    semafori, definite in \file{linux/sem.h}.} 
  \label{tab:ipc_sem_limits}
\end{table}


La funzione di sistema che permette di effettuare le varie operazioni di
controllo sui semafori fra le quali, come accennato, è impropriamente compresa
anche la loro inizializzazione, è \funcd{semctl}; il suo prototipo è:

\begin{funcproto}{
\fhead{sys/types.h}
\fhead{sys/ipc.h}
\fhead{sys/sem.h}
\fdecl{int semctl(int semid, int semnum, int cmd)}
\fdecl{int semctl(int semid, int semnum, int cmd, union semun arg)}
\fdesc{Esegue le operazioni di controllo su un semaforo o un insieme di
  semafori.}
}

{La funzione ritorna in caso di successo un valore positivo quanto usata con
  tre argomenti ed un valore nullo quando usata con quattro e $-1$ per un
  errore, nel qual caso \var{errno} assumerà uno dei valori:
  \begin{errlist}
  \item[\errcode{EACCES}] i permessi assegnati al semaforo non consentono
    l'operazione di lettura o scrittura richiesta e non si hanno i privilegi
    di amministratore.
    \item[\errcode{EIDRM}] l'insieme di semafori è stato cancellato.
    \item[\errcode{EPERM}] si è richiesto \const{IPC\_SET} o \const{IPC\_RMID}
      ma il processo non è né il creatore né il proprietario del semaforo e
      non ha i privilegi di amministratore.
    \item[\errcode{ERANGE}] si è richiesto \const{SETALL} \const{SETVAL} ma il
      valore a cui si vuole impostare il semaforo è minore di zero o maggiore
      di \const{SEMVMX}.
   \end{errlist}
   ed inoltre \errval{EFAULT} ed \errval{EINVAL} nel loro significato
   generico.}
\end{funcproto}

La funzione può avere tre o quattro argomenti, a seconda dell'operazione
specificata con \param{cmd}, ed opera o sull'intero insieme specificato da
\param{semid} o sul singolo semaforo di un insieme, specificato da
\param{semnum}. 

\begin{figure}[!htb]
  \footnotesize \centering
  \begin{minipage}[c]{0.80\textwidth}
    \includestruct{listati/semun.h}
  \end{minipage} 
  \normalsize 
  \caption{La definizione dei possibili valori di una \dirct{union}
    \structd{semun}, usata come quarto argomento della funzione
    \func{semctl}.}
  \label{fig:ipc_semun}
\end{figure}

Qualora la funzione operi con quattro argomenti \param{arg} è un argomento
generico, che conterrà un dato diverso a seconda dell'azione richiesta; per
unificare detto argomento esso deve essere passato come una unione
\struct{semun}, la cui definizione, con i possibili valori che può assumere, è
riportata in fig.~\ref{fig:ipc_semun}.

Nelle versioni più vecchie della \acr{glibc} questa unione veniva definita in
\file{sys/sem.h}, ma nelle versioni più recenti questo non avviene più in
quanto lo standard POSIX.1-2001 richiede che sia sempre definita a cura del
chiamante. In questa seconda evenienza la \acr{glibc} definisce però la
macro \macrod{\_SEM\_SEMUN\_UNDEFINED} che può essere usata per controllare la
situazione.

Come già accennato sia il comportamento della funzione che il numero di
argomenti con cui deve essere invocata dipendono dal valore dell'argomento
\param{cmd}, che specifica l'azione da intraprendere. Per questo argomento i
valori validi, quelli cioè che non causano un errore di \errcode{EINVAL}, sono
i seguenti:
\begin{basedescript}{\desclabelwidth{1.6cm}\desclabelstyle{\nextlinelabel}}
\item[\const{IPC\_STAT}] Legge i dati dell'insieme di semafori, copiandone i
  valori nella struttura \struct{semid\_ds} posta all'indirizzo specificato
  con \var{arg.buf}. Occorre avere il permesso di lettura.
  L'argomento \param{semnum} viene ignorato.
\item[\const{IPC\_RMID}] Rimuove l'insieme di semafori e le relative strutture
  dati, con effetto immediato. Tutti i processi che erano bloccati in attesa
  vengono svegliati, ritornando con un errore di \errcode{EIDRM}.  L'\ids{UID}
  effettivo del processo deve corrispondere o al creatore o al proprietario
  dell'insieme, o all'amministratore. L'argomento
  \param{semnum} viene ignorato.
\item[\const{IPC\_SET}] Permette di modificare i permessi ed il proprietario
  dell'insieme. I valori devono essere passati in una struttura
  \struct{semid\_ds} puntata da \param{arg.buf} di cui saranno usati soltanto
  i campi \var{sem\_perm.uid}, \var{sem\_perm.gid} e i nove bit meno
  significativi di \var{sem\_perm.mode}. La funziona aggiorna anche il campo
  \var{sem\_ctime}.  L'\ids{UID} effettivo del processo deve corrispondere o
  al creatore o al proprietario dell'insieme, o all'amministratore.
  L'argomento \param{semnum} viene ignorato.
\item[\constd{GETALL}] Restituisce il valore corrente di ciascun semaforo
  dell'insieme (corrispondente al campo \var{semval} di \struct{sem}) nel
  vettore indicato da \param{arg.array}. Occorre avere il permesso di lettura.
  L'argomento \param{semnum} viene ignorato.
\item[\constd{GETNCNT}] Restituisce come valore di ritorno della funzione il
  numero di processi in attesa che il semaforo \param{semnum} dell'insieme
  \param{semid} venga incrementato (corrispondente al campo \var{semncnt} di
  \struct{sem}). Va invocata con tre argomenti.  Occorre avere il permesso di
  lettura.
\item[\constd{GETPID}] Restituisce come valore di ritorno della funzione il
  \ids{PID} dell'ultimo processo che ha compiuto una operazione sul semaforo
  \param{semnum} dell'insieme \param{semid} (corrispondente al campo
  \var{sempid} di \struct{sem}). Va invocata con tre argomenti.  Occorre avere
  il permesso di lettura.
\item[\constd{GETVAL}] Restituisce come valore di ritorno della funzione il 
  valore corrente del semaforo \param{semnum} dell'insieme \param{semid}
  (corrispondente al campo \var{semval} di \struct{sem}). Va invocata con tre
  argomenti.  Occorre avere il permesso di lettura.
\item[\constd{GETZCNT}] Restituisce come valore di ritorno della funzione il
  numero di processi in attesa che il valore del semaforo \param{semnum}
  dell'insieme \param{semid} diventi nullo (corrispondente al campo
  \var{semncnt} di \struct{sem}). Va invocata con tre argomenti.  Occorre
  avere il permesso di lettura.
\item[\constd{SETALL}] Inizializza il valore di tutti i semafori dell'insieme,
  aggiornando il campo \var{sem\_ctime} di \struct{semid\_ds}. I valori devono
  essere passati nel vettore indicato da \param{arg.array}.  Si devono avere i
  privilegi di scrittura.  L'argomento \param{semnum} viene ignorato.
\item[\constd{SETVAL}] Inizializza il semaforo \param{semnum} al valore passato
  dall'argomento \param{arg.val}, aggiornando il campo \var{sem\_ctime} di
  \struct{semid\_ds}.  Si devono avere i privilegi di scrittura.
\end{basedescript}

Come per \func{msgctl} esistono tre ulteriori valori, \const{IPC\_INFO},
\constd{SEM\_STAT} e \constd{SEM\_INFO}, specifici di Linux e fuori da ogni
standard, creati specificamente ad uso del comando \cmd{ipcs}. Dato che anche
questi potranno essere modificati o rimossi, non devono essere utilizzati e
pertanto non li tratteremo.

Quando si imposta il valore di un semaforo (sia che lo si faccia per tutto
l'insieme con \const{SETALL}, che per un solo semaforo con \const{SETVAL}), i
processi in attesa su di esso reagiscono di conseguenza al cambiamento di
valore.  Inoltre la coda delle operazioni di ripristino viene cancellata per
tutti i semafori il cui valore viene modificato.

\begin{table}[htb]
  \footnotesize
  \centering
  \begin{tabular}[c]{|c|l|}
    \hline
    \textbf{Operazione}  & \textbf{Valore restituito} \\
    \hline
    \hline
    \const{GETNCNT}& Valore di \var{semncnt}.\\
    \const{GETPID} & Valore di \var{sempid}.\\
    \const{GETVAL} & Valore di \var{semval}.\\
    \const{GETZCNT}& Valore di \var{semzcnt}.\\
    \hline
  \end{tabular}
  \caption{Valori di ritorno della funzione \func{semctl}.} 
  \label{tab:ipc_semctl_returns}
\end{table}

Il valore di ritorno della funzione in caso di successo dipende
dall'operazione richiesta; per tutte le operazioni che richiedono quattro
argomenti esso è sempre nullo, per le altre operazioni, elencate in
tab.~\ref{tab:ipc_semctl_returns} viene invece restituito il valore richiesto,
corrispondente al campo della struttura \struct{sem} indicato nella seconda
colonna della tabella.

Le operazioni ordinarie sui semafori, come l'acquisizione o il rilascio degli
stessi (in sostanza tutte quelle non comprese nell'uso di \func{semctl})
vengono effettuate con la funzione di sistema \funcd{semop}, il cui prototipo
è:

\begin{funcproto}{
\fhead{sys/types.h}
\fhead{sys/ipc.h}
\fhead{sys/sem.h}
\fdecl{int semop(int semid, struct sembuf *sops, unsigned nsops)}
\fdesc{Esegue operazioni ordinarie su un semaforo o un insieme di semafori.} 
}

{La funzione ritorna $0$ in caso di successo e $-1$ per un errore, nel qual
  caso \var{errno} assumerà uno dei valori: 
  \begin{errlist}
    \item[\errcode{E2BIG}] l'argomento \param{nsops} è maggiore del numero
      massimo di operazioni \const{SEMOPM}.
    \item[\errcode{EACCES}] il processo non ha i permessi per eseguire
      l'operazione richiesta e non ha i privilegi di amministratore.
    \item[\errcode{EAGAIN}] un'operazione comporterebbe il blocco del processo,
      ma si è specificato \const{IPC\_NOWAIT} in \var{sem\_flg}.
    \item[\errcode{EFBIG}] il valore del campo \var{sem\_num} è negativo o
      maggiore o uguale al numero di semafori dell'insieme.
    \item[\errcode{EIDRM}] l'insieme di semafori è stato cancellato.
    \item[\errcode{EINTR}] la funzione, bloccata in attesa dell'esecuzione
      dell'operazione, viene interrotta da un segnale.
    \item[\errcode{ENOMEM}] si è richiesto un \const{SEM\_UNDO} ma il sistema
      non ha le risorse per allocare la struttura di ripristino.
    \item[\errcode{ERANGE}] per alcune operazioni il valore risultante del
      semaforo viene a superare il limite massimo \const{SEMVMX}.
  \end{errlist}
  ed inoltre \errval{EFAULT} ed \errval{EINVAL} nel loro significato
  generico.}
\end{funcproto}

La funzione permette di eseguire operazioni multiple sui singoli semafori di
un insieme. La funzione richiede come primo argomento l'identificatore
\param{semid} dell'insieme su cui si vuole operare, il numero di operazioni da
effettuare viene specificato con l'argomento \param{nsop}, mentre il loro
contenuto viene passato con un puntatore ad un vettore di strutture
\struct{sembuf} nell'argomento \param{sops}. Le operazioni richieste vengono
effettivamente eseguite se e soltanto se è possibile effettuarle tutte quante,
ed in tal caso vengono eseguite nella sequenza passata nel
vettore \param{sops}.

Con lo standard POSIX.1-2001 è stata introdotta una variante di \func{semop}
che consente di specificare anche un tempo massimo di attesa. La nuova
funzione di sistema, disponibile a partire dal kernel 2.4.22 e dalla
\acr{glibc} 2.3.3, ed utilizzabile solo dopo aver definito la macro
\macro{\_GNU\_SOURCE}, è \funcd{semtimedop}, ed il suo prototipo è:

\begin{funcproto}{
\fhead{sys/types.h}
\fhead{sys/ipc.h}
\fhead{sys/sem.h}
\fdecl{int semtimedop(int semid, struct sembuf *sops, unsigned nsops,
                      struct timespec *timeout)}
\fdesc{Esegue operazioni ordinarie su un semaforo o un insieme di semafori.} 
}

{La funzione ritorna $0$ in caso di successo e $-1$ per un errore, nel qual
  caso \var{errno} assumerà uno dei valori: 
  \begin{errlist}
  \item[\errcode{EAGAIN}] l'operazione comporterebbe il blocco del processo,
    ma si è specificato \const{IPC\_NOWAIT} in \var{sem\_flg} oppure si è
    atteso oltre quanto indicato da \param{timeout}.
  \end{errlist}
  e gli altri valori già visti per \func{semop}, con lo stesso significato.}
\end{funcproto}

Rispetto a \func{semop} la funzione consente di specificare un tempo massimo
di attesa, indicato con una struttura \struct{timespec} (vedi
fig.~\ref{fig:sys_timespec_struct}), per le operazioni che verrebbero
bloccate. Alla scadenza di detto tempo la funzione ritorna comunque con un
errore di \errval{EAGAIN} senza che nessuna delle operazioni richieste venga
eseguita. 

Si tenga presente che la precisione della temporizzazione è comunque limitata
dalla risoluzione dell'orologio di sistema, per cui il tempo di attesa verrà
arrotondato per eccesso. In caso si passi un valore \val{NULL}
per \param{timeout} il comportamento di \func{semtimedop} è identico a quello
di \func{semop}.


\begin{figure}[!htb]
  \footnotesize \centering
  \begin{minipage}[c]{.80\textwidth}
    \includestruct{listati/sembuf.h}
  \end{minipage} 
  \normalsize 
  \caption{La struttura \structd{sembuf}, usata per le operazioni sui
    semafori.}
  \label{fig:ipc_sembuf}
\end{figure}

Come indicato il contenuto di ciascuna operazione deve essere specificato
attraverso una struttura \struct{sembuf} (la cui definizione è riportata in
fig.~\ref{fig:ipc_sembuf}) che il programma chiamante deve avere cura di
allocare in un opportuno vettore. La struttura permette di indicare il
semaforo su cui operare, il tipo di operazione, ed un flag di controllo.  

Il campo \var{sem\_num} serve per indicare a quale semaforo dell'insieme fa
riferimento l'operazione. Si ricordi che i semafori sono numerati come gli
elementi di un vettore, per cui il primo semaforo di un insieme corrisponde ad
un valore nullo di \var{sem\_num}.

Il campo \var{sem\_flg} è un flag, mantenuto come maschera binaria, per il
quale possono essere impostati i due valori \const{IPC\_NOWAIT} e
\constd{SEM\_UNDO}.  Impostando \const{IPC\_NOWAIT} si fa sì che in tutti quei
casi in cui l'esecuzione di una operazione richiederebbe di porre il processo
vada nello stato di \textit{sleep}, invece di bloccarsi \func{semop} ritorni
immediatamente (abortendo così le eventuali operazioni restanti) con un errore
di \errcode{EAGAIN}.  Impostando \const{SEM\_UNDO} si richiede invece che
l'operazione in questione venga registrata, in modo che il valore del semaforo
possa essere ripristinato all'uscita del processo.

Infine \var{sem\_op} è il campo che controlla qual'è l'operazione che viene
eseguita e determina in generale il comportamento della chiamata a
\func{semop}. I casi possibili per il valore di questo campo sono tre:
\begin{basedescript}{\desclabelwidth{1.8cm}}
\item[\var{sem\_op} $>0$] In questo caso il valore viene aggiunto al valore
  corrente di \var{semval} per il semaforo indicato. Questa operazione non
  causa mai un blocco del processo, ed eventualmente \func{semop} ritorna
  immediatamente con un errore di \errcode{ERANGE} qualora si sia superato il
  limite \const{SEMVMX}. Se l'operazione ha successo si passa immediatamente
  alla successiva.  Specificando \const{SEM\_UNDO} si aggiorna il contatore
  per il ripristino del valore del semaforo. Al processo chiamante è richiesto
  il privilegio di alterazione (scrittura) sull'insieme di semafori.
  
\item[\var{sem\_op} $=0$] Nel caso \var{semval} sia zero l'operazione ha
  successo immediato, e o si passa alla successiva o \func{semop} ritorna con
  successo se questa era l'ultima. Se \var{semval} è diverso da zero il
  comportamento è controllato da \var{sem\_flg}, se è stato impostato
  \const{IPC\_NOWAIT} \func{semop} ritorna immediatamente abortendo tutte le
  operazioni con un errore di \errcode{EAGAIN}, altrimenti viene incrementato
  \var{semzcnt} di uno ed il processo viene bloccato fintanto che non si
  verifica una delle condizioni seguenti:
  \begin{itemize*}
  \item \var{semval} diventa zero, nel qual caso \var{semzcnt} viene
    decrementato di uno, l'operazione ha successo e si passa alla successiva,
    oppure \func{semop} ritorna con successo se questa era l'ultima.
  \item l'insieme di semafori viene rimosso, nel qual caso \func{semop}
    ritorna abortendo tutte le operazioni con un errore di \errcode{EIDRM}.
  \item il processo chiamante riceve un segnale, nel qual caso \var{semzcnt}
    viene decrementato di uno e \func{semop} ritorna abortendo tutte le
    operazioni con un errore di \errcode{EINTR}.
  \end{itemize*}
  Al processo chiamante è richiesto soltanto il privilegio di lettura
  dell'insieme dei semafori.
  
\item[\var{sem\_op} $<0$] Nel caso in cui \var{semval} è maggiore o uguale del
  valore assoluto di \var{sem\_op} (se cioè la somma dei due valori resta
  positiva o nulla) i valori vengono sommati e l'operazione ha successo e si
  passa alla successiva, oppure \func{semop} ritorna con successo se questa
  era l'ultima. Qualora si sia impostato \const{SEM\_UNDO} viene anche
  aggiornato il contatore per il ripristino del valore del semaforo. In caso
  contrario (quando cioè la somma darebbe luogo ad un valore di \var{semval}
  negativo) se si è impostato \const{IPC\_NOWAIT} \func{semop} ritorna
  immediatamente abortendo tutte le operazioni con un errore di
  \errcode{EAGAIN}, altrimenti viene incrementato di uno \var{semncnt} ed il
  processo resta in stato di \textit{sleep} fintanto che non si ha una delle
  condizioni seguenti:
  \begin{itemize*}
  \item \var{semval} diventa maggiore o uguale del valore assoluto di
    \var{sem\_op}, nel qual caso \var{semncnt} viene decrementato di uno, il
    valore di \var{sem\_op} viene sommato a \var{semval}, e se era stato
    impostato \const{SEM\_UNDO} viene aggiornato il contatore per il
    ripristino del valore del semaforo.
  \item l'insieme di semafori viene rimosso, nel qual caso \func{semop}
    ritorna abortendo tutte le operazioni con un errore di \errcode{EIDRM}.
  \item il processo chiamante riceve un segnale, nel qual caso \var{semncnt}
    viene decrementato di uno e \func{semop} ritorna abortendo tutte le
    operazioni con un errore di \errcode{EINTR}.
  \end{itemize*}    
  Al processo chiamante è richiesto il privilegio di alterazione (scrittura)
  sull'insieme di semafori.
\end{basedescript}

Qualora si sia usato \func{semtimedop} alle condizioni di errore precedenti si
aggiunge anche quella di scadenza del tempo di attesa indicato
con \param{timeout} che farà abortire la funzione, qualora resti bloccata
troppo a lungo nell'esecuzione delle operazioni richieste, con un errore di
\errcode{EAGAIN}.

In caso di successo (sia per \func{semop} che per \func{semtimedop}) per ogni
semaforo modificato verrà aggiornato il campo \var{sempid} al valore del
\ids{PID} del processo chiamante; inoltre verranno pure aggiornati al tempo
corrente i campi \var{sem\_otime} e \var{sem\_ctime}.

Dato che, come già accennato in precedenza, in caso di uscita inaspettata i
semafori possono restare occupati, abbiamo visto come \func{semop} (e
\func{semtimedop}) permetta di attivare un meccanismo di ripristino attraverso
l'uso del flag \const{SEM\_UNDO}. Il meccanismo è implementato tramite una
apposita struttura \kstruct{sem\_undo}, associata ad ogni processo per ciascun
semaforo che esso ha modificato; all'uscita i semafori modificati vengono
ripristinati, e le strutture disallocate.  Per mantenere coerente il
comportamento queste strutture non vengono ereditate attraverso una
\func{fork} (altrimenti si avrebbe un doppio ripristino), mentre passano
inalterate nell'esecuzione di una \func{exec} (altrimenti non si avrebbe
ripristino).

Tutto questo però ha un problema di fondo. Per capire di cosa si tratta
occorre fare riferimento all'implementazione usata in Linux, che è riportata
in maniera semplificata nello schema di fig.~\ref{fig:ipc_sem_schema}.  Si è
presa come riferimento l'architettura usata fino al kernel 2.2.x che è più
semplice (ed illustrata in dettaglio in \cite{tlk}). Nel kernel 2.4.x la
struttura del \textit{SysV-IPC} è stata modificata, ma le definizioni relative
a queste strutture restano per compatibilità (in particolare con le vecchie
versioni delle librerie del C, come le \acr{libc5}).

\begin{figure}[!htb]
  \centering \includegraphics[width=12cm]{img/semtruct}
  \caption{Schema delle varie strutture di un insieme di semafori
    (\kstructd{semid\_ds}, \kstructd{sem}, \kstructd{sem\_queue} e
    \kstructd{sem\_undo}).}
  \label{fig:ipc_sem_schema}
\end{figure}

Alla creazione di un nuovo insieme viene allocata una nuova strutture
\kstruct{semid\_ds} ed il relativo vettore di strutture \kstruct{sem}. Quando
si richiede una operazione viene anzitutto verificato che tutte le operazioni
possono avere successo; se una di esse comporta il blocco del processo il
kernel crea una struttura \kstruct{sem\_queue} che viene aggiunta in fondo
alla coda di attesa associata a ciascun insieme di semafori, che viene
referenziata tramite i campi \var{sem\_pending} e \var{sem\_pending\_last} di
\kstruct{semid\_ds}.  Nella struttura viene memorizzato il riferimento alle
operazioni richieste (nel campo \var{sops}, che è un puntatore ad una
struttura \struct{sembuf}) e al processo corrente (nel campo \var{sleeper})
poi quest'ultimo viene messo stato di attesa e viene invocato lo
\textit{scheduler} per passare all'esecuzione di un altro processo.

Se invece tutte le operazioni possono avere successo vengono eseguite
immediatamente, dopo di che il kernel esegue una scansione della coda di
attesa (a partire da \var{sem\_pending}) per verificare se qualcuna delle
operazioni sospese in precedenza può essere eseguita, nel qual caso la
struttura \kstruct{sem\_queue} viene rimossa e lo stato del processo associato
all'operazione (\var{sleeper}) viene riportato a \textit{running}; il tutto
viene ripetuto fin quando non ci sono più operazioni eseguibili o si è
svuotata la coda.  Per gestire il meccanismo del ripristino tutte le volte che
per un'operazione si è specificato il flag \const{SEM\_UNDO} viene mantenuta
per ciascun insieme di semafori una apposita struttura \kstruct{sem\_undo} che
contiene (nel vettore puntato dal campo \var{semadj}) un valore di
aggiustamento per ogni semaforo cui viene sommato l'opposto del valore usato
per l'operazione.

Queste strutture sono mantenute in due liste (rispettivamente attraverso i due
campi \var{id\_next} e \var{proc\_next}) una associata all'insieme di cui fa
parte il semaforo, che viene usata per invalidare le strutture se questo viene
cancellato o per azzerarle se si è eseguita una operazione con \func{semctl},
l'altra associata al processo che ha eseguito l'operazione, attraverso il
campo \var{semundo} di \kstruct{task\_struct}, come mostrato in
\ref{fig:ipc_sem_schema}. Quando un processo termina, la lista ad esso
associata viene scandita e le operazioni applicate al semaforo.  Siccome un
processo può accumulare delle richieste di ripristino per semafori differenti
attraverso diverse chiamate a \func{semop}, si pone il problema di come
eseguire il ripristino dei semafori all'uscita del processo, ed in particolare
se questo può essere fatto atomicamente.

Il punto è cosa succede quando una delle operazioni previste per il ripristino
non può essere eseguita immediatamente perché ad esempio il semaforo è
occupato; in tal caso infatti, se si pone il processo in stato di
\textit{sleep} aspettando la disponibilità del semaforo (come faceva
l'implementazione originaria) si perde l'atomicità dell'operazione. La scelta
fatta dal kernel è pertanto quella di effettuare subito le operazioni che non
prevedono un blocco del processo e di ignorare silenziosamente le altre;
questo però comporta il fatto che il ripristino non è comunque garantito in
tutte le occasioni.

Come esempio di uso dell'interfaccia dei semafori vediamo come implementare
con essa dei semplici \textit{mutex} (cioè semafori binari), tutto il codice
in questione, contenuto nel file \file{Mutex.c} allegato ai sorgenti, è
riportato in fig.~\ref{fig:ipc_mutex_create}. Utilizzeremo l'interfaccia per
creare un insieme contenente un singolo semaforo, per il quale poi useremo un
valore unitario per segnalare la disponibilità della risorsa, ed un valore
nullo per segnalarne l'indisponibilità. 

\begin{figure}[!htbp]
  \footnotesize \centering
  \begin{minipage}[c]{\codesamplewidth}
    \includecodesample{listati/Mutex.c}
  \end{minipage} 
  \normalsize 
  \caption{Il codice delle funzioni che permettono di creare o recuperare
    l'identificatore di un semaforo da utilizzare come \textit{mutex}.}
  \label{fig:ipc_mutex_create}
\end{figure}

La prima funzione (\texttt{\small 2-15}) è \func{MutexCreate} che data una
chiave crea il semaforo usato per il mutex e lo inizializza, restituendone
l'identificatore. Il primo passo (\texttt{\small 6}) è chiamare \func{semget}
con \const{IPC\_CREATE} per creare il semaforo qualora non esista,
assegnandogli i privilegi di lettura e scrittura per tutti. In caso di errore
(\texttt{\small 7-9}) si ritorna subito il risultato di \func{semget},
altrimenti (\texttt{\small 10}) si inizializza il semaforo chiamando
\func{semctl} con il comando \const{SETVAL}, utilizzando l'unione
\struct{semunion} dichiarata ed avvalorata in precedenza (\texttt{\small 4})
ad 1 per significare che risorsa è libera. In caso di errore (\texttt{\small
  11-13}) si restituisce il valore di ritorno di \func{semctl}, altrimenti
(\texttt{\small 14}) si ritorna l'identificatore del semaforo.

La seconda funzione (\texttt{\small 17-20}) è \func{MutexFind}, che, data una
chiave, restituisce l'identificatore del semaforo ad essa associato. La
comprensione del suo funzionamento è immediata in quanto essa è soltanto un
\textit{wrapper}\footnote{si chiama così una funzione usata per fare da
  \textsl{involucro} alla chiamata di un altra, usata in genere per
  semplificare un'interfaccia (come in questo caso) o per utilizzare con la
  stessa funzione diversi substrati (librerie, ecc.)  che possono fornire le
  stesse funzionalità.} di una chiamata a \func{semget} per cercare
l'identificatore associato alla chiave, il valore di ritorno di quest'ultima
viene passato all'indietro al chiamante.

La terza funzione (\texttt{\small 22-25}) è \func{MutexRead} che, dato un
identificatore, restituisce il valore del semaforo associato al mutex. Anche
in questo caso la funzione è un \textit{wrapper} per una chiamata a
\func{semctl} con il comando \const{GETVAL}, che permette di restituire il
valore del semaforo.

La quarta e la quinta funzione (\texttt{\small 36-44}) sono \func{MutexLock},
e \func{MutexUnlock}, che permettono rispettivamente di bloccare e sbloccare
il mutex. Entrambe fanno da wrapper per \func{semop}, utilizzando le due
strutture \var{sem\_lock} e \var{sem\_unlock} definite in precedenza
(\texttt{\small 27-34}). Si noti come per queste ultime si sia fatto uso
dell'opzione \const{SEM\_UNDO} per evitare che il semaforo resti bloccato in
caso di terminazione imprevista del processo.

L'ultima funzione (\texttt{\small 46-49}) della serie, è \func{MutexRemove},
che rimuove il mutex. Anche in questo caso si ha un wrapper per una chiamata a
\func{semctl} con il comando \const{IPC\_RMID}, che permette di cancellare il
semaforo; il valore di ritorno di quest'ultima viene passato all'indietro.

Chiamare \func{MutexLock} decrementa il valore del semaforo: se questo è
libero (ha già valore 1) sarà bloccato (valore nullo), se è bloccato la
chiamata a \func{semop} si bloccherà fintanto che la risorsa non venga
rilasciata. Chiamando \func{MutexUnlock} il valore del semaforo sarà
incrementato di uno, sbloccandolo qualora fosse bloccato.  

Si noti che occorre eseguire sempre prima \func{MutexLock} e poi
\func{MutexUnlock}, perché se per un qualche errore si esegue più volte
quest'ultima il valore del semaforo crescerebbe oltre 1, e \func{MutexLock}
non avrebbe più l'effetto aspettato (bloccare la risorsa quando questa è
considerata libera).  Infine si tenga presente che usare \func{MutexRead} per
controllare il valore dei mutex prima di proseguire in una operazione di
sblocco non servirebbe comunque, dato che l'operazione non sarebbe atomica.
Vedremo in sez.~\ref{sec:ipc_lock_file} come sia possibile ottenere
un'interfaccia analoga a quella appena illustrata, senza incorrere in questi
problemi, usando il \textit{file locking}.


\subsection{Memoria condivisa}
\label{sec:ipc_sysv_shm}

Il terzo oggetto introdotto dal \textit{SysV-IPC} è quello dei segmenti di
memoria condivisa. La funzione di sistema che permette di ottenerne uno è
\funcd{shmget}, ed il suo prototipo è:

\begin{funcproto}{
\fhead{sys/types.h}
\fhead{sys/ipc.h}
\fhead{sys/shm.h}
\fdecl{int shmget(key\_t key, int size, int flag)}
\fdesc{Ottiene o crea una memoria condivisa.} 
}

{La funzione ritorna l'identificatore (un intero positivo) in caso di successo
  e $-1$ per un errore, nel qual caso \var{errno} assumerà uno dei valori:
  \begin{errlist}
    \item[\errcode{ENOSPC}] si è superato il limite (\const{SHMMNI}) sul numero
      di segmenti di memoria nel sistema, o cercato di allocare un segmento le
      cui dimensioni fanno superare il limite di sistema (\const{SHMALL}) per
      la memoria ad essi riservata.
    \item[\errcode{EINVAL}] si è richiesta una dimensione per un nuovo segmento
      maggiore di \const{SHMMAX} o minore di \const{SHMMIN}, o se il segmento
      già esiste \param{size} è maggiore delle sue dimensioni.
    \item[\errcode{ENOMEM}] il sistema non ha abbastanza memoria per poter
      contenere le strutture per un nuovo segmento di memoria condivisa.
    \item[\errcode{ENOMEM}] si è specificato \const{IPC\_HUGETLB} ma non si
      hanno i privilegi di amministratore.
   \end{errlist}
   ed inoltre \errval{EACCES}, \errval{ENOENT}, \errval{EEXIST},
   \errval{EIDRM}, con lo stesso significato che hanno per \func{msgget}.}
\end{funcproto}


La funzione, come \func{semget}, è analoga a \func{msgget}, ed identico è
l'uso degli argomenti \param{key} e \param{flag} per cui non ripeteremo quanto
detto al proposito in sez.~\ref{sec:ipc_sysv_mq}.  A partire dal kernel 2.6
però sono stati introdotti degli ulteriori bit di controllo per
l'argomento \param{flag}, specifici di \func{shmget}, attinenti alle modalità
di gestione del segmento di memoria condivisa in relazione al sistema della
memoria virtuale.

Il primo dei due flag è \constd{SHM\_HUGETLB} che consente di richiedere la
creazione del segmento usando una \textit{huge page}, le pagine di memoria di
grandi dimensioni introdotte con il kernel 2.6 per ottimizzare le prestazioni
nei sistemi più recenti che hanno grandi quantità di memoria. L'operazione è
privilegiata e richiede che il processo abbia la \textit{capability}
\const{CAP\_IPC\_LOCK}. Questa funzionalità è specifica di Linux e non è
portabile.

Il secondo flag aggiuntivo, introdotto a partire dal kernel 2.6.15, è
\constd{SHM\_NORESERVE}, ed ha lo stesso scopo del flag \const{MAP\_NORESERVE}
di \func{mmap} (vedi sez.~\ref{sec:file_memory_map}): non vengono riservate
delle pagine di swap ad uso del meccanismo del \textit{copy on write} per
mantenere le modifiche fatte sul segmento. Questo significa che caso di
scrittura sul segmento quando non c'è più memoria disponibile, si avrà
l'emissione di un \signal{SIGSEGV}.

Infine l'argomento \param{size} specifica la dimensione del segmento di
memoria condivisa; il valore deve essere specificato in byte, ma verrà
comunque arrotondato al multiplo superiore di \const{PAGE\_SIZE}. Il valore
deve essere specificato quando si crea un nuovo segmento di memoria con
\const{IPC\_CREAT} o \const{IPC\_PRIVATE}, se invece si accede ad un segmento
di memoria condivisa esistente non può essere maggiore del valore con cui esso
è stato creato.

La memoria condivisa è la forma più veloce di comunicazione fra due processi,
in quanto permette agli stessi di vedere nel loro spazio di indirizzi una
stessa sezione di memoria.  Pertanto non è necessaria nessuna operazione di
copia per trasmettere i dati da un processo all'altro, in quanto ciascuno può
accedervi direttamente con le normali operazioni di lettura e scrittura dei
dati in memoria.

Ovviamente tutto questo ha un prezzo, ed il problema fondamentale della
memoria condivisa è la sincronizzazione degli accessi. È evidente infatti che
se un processo deve scambiare dei dati con un altro, si deve essere sicuri che
quest'ultimo non acceda al segmento di memoria condivisa prima che il primo
non abbia completato le operazioni di scrittura, inoltre nel corso di una
lettura si deve essere sicuri che i dati restano coerenti e non vengono
sovrascritti da un accesso in scrittura sullo stesso segmento da parte di un
altro processo. Per questo in genere la memoria condivisa viene sempre
utilizzata in abbinamento ad un meccanismo di sincronizzazione, il che, di
norma, significa insieme a dei semafori.

\begin{figure}[!htb]
  \footnotesize \centering
  \begin{minipage}[c]{0.80\textwidth}
    \includestruct{listati/shmid_ds.h}
  \end{minipage} 
  \normalsize 
  \caption{La struttura \structd{shmid\_ds}, associata a ciascun segmento di
    memoria condivisa.}
  \label{fig:ipc_shmid_ds}
\end{figure}

A ciascun segmento di memoria condivisa è associata una struttura
\struct{shmid\_ds}, riportata in fig.~\ref{fig:ipc_shmid_ds}.  Come nel caso
delle code di messaggi quando si crea un nuovo segmento di memoria condivisa
con \func{shmget} questa struttura viene inizializzata, in particolare il
campo \var{shm\_perm} viene inizializzato come illustrato in
sez.~\ref{sec:ipc_sysv_access_control}, e valgono le considerazioni ivi fatte
relativamente ai permessi di accesso; per quanto riguarda gli altri campi
invece:
\begin{itemize*}
\item il campo \var{shm\_segsz}, che esprime la dimensione del segmento, viene
  inizializzato al valore di \param{size}.
\item il campo \var{shm\_ctime}, che esprime il tempo di creazione del
  segmento, viene inizializzato al tempo corrente.
\item i campi \var{shm\_atime} e \var{shm\_dtime}, che esprimono
  rispettivamente il tempo dell'ultima volta che il segmento è stato
  agganciato o sganciato da un processo, vengono inizializzati a zero.
\item il campo \var{shm\_lpid}, che esprime il \ids{PID} del processo che ha
  eseguito l'ultima operazione, viene inizializzato a zero.
\item il campo \var{shm\_cpid}, che esprime il \ids{PID} del processo che ha
  creato il segmento, viene inizializzato al \ids{PID} del processo chiamante.
\item il campo \var{shm\_nattac}, che esprime il numero di processi agganciati
  al segmento viene inizializzato a zero.
\end{itemize*}

Come per le code di messaggi e gli insiemi di semafori, anche per i segmenti
di memoria condivisa esistono una serie di limiti imposti dal sistema.  Alcuni
di questi limiti sono al solito accessibili e modificabili attraverso
\func{sysctl} o scrivendo direttamente nei rispettivi file di
\file{/proc/sys/kernel/}. 

In tab.~\ref{tab:ipc_shm_limits} si sono riportate le
costanti simboliche associate a ciascuno di essi, il loro significato, i
valori preimpostati, e, quando presente, il file in \file{/proc/sys/kernel/}
che permettono di cambiarne il valore. 


\begin{table}[htb]
  \footnotesize
  \centering
  \begin{tabular}[c]{|c|r|c|p{7cm}|}
    \hline
    \textbf{Costante} & \textbf{Valore} & \textbf{File in \texttt{proc}}
    & \textbf{Significato} \\
    \hline
    \hline
    \constd{SHMALL}& 0x200000&\sysctlrelfiled{kernel}{shmall}
                             & Numero massimo di pagine che 
                               possono essere usate per i segmenti di
                               memoria condivisa.\\
    \constd{SHMMAX}&0x2000000&\sysctlrelfiled{kernel}{shmmax} 
                             & Dimensione massima di un segmento di memoria
                               condivisa.\\ 
    \constd{SHMMNI}&     4096&\sysctlrelfiled{kernel}{shmmni}
                             & Numero massimo di segmenti di memoria condivisa
                              presenti nel kernel.\\ 
    \constd{SHMMIN}&        1& ---         & Dimensione minima di un segmento di
                                             memoria condivisa.\\
    \constd{SHMLBA}&\const{PAGE\_SIZE}&--- & Limite inferiore per le dimensioni
                                             minime di un segmento (deve essere
                                             allineato alle dimensioni di una
                                             pagina di memoria).\\
    \constd{SHMSEG}&   ---   &     ---     & Numero massimo di segmenti di
                                             memoria condivisa per ciascun
                                             processo (l'implementazione non
                                             prevede l'esistenza di questo
                                             limite).\\


    \hline
  \end{tabular}
  \caption{Valori delle costanti associate ai limiti dei segmenti di memoria
    condivisa, insieme al relativo file in \file{/proc/sys/kernel/} ed al
    valore preimpostato presente nel sistema.} 
  \label{tab:ipc_shm_limits}
\end{table}

Al solito la funzione di sistema che permette di effettuare le operazioni di
controllo su un segmento di memoria condivisa è \funcd{shmctl}; il suo
prototipo è:

\begin{funcproto}{
\fhead{sys/ipc.h}
\fhead{sys/shm.h}
\fdecl{int shmctl(int shmid, int cmd, struct shmid\_ds *buf)}

\fdesc{Esegue le operazioni di controllo su un segmento di memoria condivisa.}
}

{La funzione ritorna $0$ in caso di successo e $-1$ per un errore, nel qual
  caso \var{errno} assumerà uno dei valori:
  \begin{errlist}
    \item[\errcode{EACCES}] si è richiesto \const{IPC\_STAT} ma i permessi non
      consentono l'accesso in lettura al segmento.
    \item[\errcode{EFAULT}] l'indirizzo specificato con \param{buf} non è
      valido.
    \item[\errcode{EIDRM}] l'argomento \param{shmid} fa riferimento ad un
      segmento che è stato cancellato.
    \item[\errcode{EINVAL}] o \param{shmid} non è un identificatore valido o
      \param{cmd} non è un comando valido.
    \item[\errcode{ENOMEM}] si è richiesto un \textit{memory lock} di
      dimensioni superiori al massimo consentito.
    \item[\errcode{EOVERFLOW}] si è tentato il comando \const{IPC\_STAT} ma il
      valore del \ids{GID} o dell'\ids{UID} è troppo grande per essere
      memorizzato nella struttura puntata da \param{buf}.
    \item[\errcode{EPERM}] si è specificato un comando con \const{IPC\_SET} o
      \const{IPC\_RMID} senza i permessi necessari.
  \end{errlist}
}  
\end{funcproto}

Il comando specificato attraverso l'argomento \param{cmd} determina i diversi
effetti della funzione. Nello standard POSIX.1-2001 i valori che esso può
assumere, ed il corrispondente comportamento della funzione, sono i seguenti:

\begin{basedescript}{\desclabelwidth{2.2cm}\desclabelstyle{\nextlinelabel}}
\item[\const{IPC\_STAT}] Legge le informazioni riguardo il segmento di memoria
  condivisa nella struttura \struct{shmid\_ds} puntata da \param{buf}. Occorre
  che il processo chiamante abbia il permesso di lettura sulla segmento.
\item[\const{IPC\_RMID}] Marca il segmento di memoria condivisa per la
  rimozione, questo verrà cancellato effettivamente solo quando l'ultimo
  processo ad esso agganciato si sarà staccato. Questo comando può essere
  eseguito solo da un processo con \ids{UID} effettivo corrispondente o al
  creatore del segmento, o al proprietario del segmento, o all'amministratore.
\item[\const{IPC\_SET}] Permette di modificare i permessi ed il proprietario
  del segmento.  Per modificare i valori di \var{shm\_perm.mode},
  \var{shm\_perm.uid} e \var{shm\_perm.gid} occorre essere il proprietario o
  il creatore del segmento, oppure l'amministratore. Compiuta l'operazione
  aggiorna anche il valore del campo \var{shm\_ctime}.
\end{basedescript}

Oltre ai precedenti su Linux sono definiti anche degli ulteriori comandi, che
consentono di estendere le funzionalità, ovviamente non devono essere usati se
si ha a cuore la portabilità. Questi comandi aggiuntivi sono:

\begin{basedescript}{\desclabelwidth{2.2cm}\desclabelstyle{\nextlinelabel}}
\item[\constd{SHM\_LOCK}] Abilita il \textit{memory locking} sul segmento di
  memoria condivisa, impedendo che la memoria usata per il segmento venga
  salvata su disco dal meccanismo della memoria virtuale. Come illustrato in
  sez.~\ref{sec:proc_mem_lock} fino al kernel 2.6.9 solo l'amministratore
  poteva utilizzare questa capacità,\footnote{che richiedeva la
    \textit{capability} \const{CAP\_IPC\_LOCK}.} a partire dal kernel 2.6.10
  anche gli utenti normali possono farlo fino al limite massimo determinato da
  \const{RLIMIT\_MEMLOCK} (vedi sez.~\ref{sec:sys_resource_limit}).
\item[\constd{SHM\_UNLOCK}] Disabilita il \textit{memory locking} sul segmento
  di memoria condivisa.  Fino al kernel 2.6.9 solo l'amministratore poteva
  utilizzare questo comando in corrispondenza di un segmento da lui bloccato.
\end{basedescript}

A questi due, come per \func{msgctl} e \func{semctl}, si aggiungono tre
ulteriori valori, \const{IPC\_INFO}, \constd{SHM\_STAT} e \constd{SHM\_INFO},
introdotti ad uso del programma \cmd{ipcs} per ottenere le informazioni
generali relative alle risorse usate dai segmenti di memoria condivisa. Dato
che potranno essere modificati o rimossi in favore dell'uso di \texttt{/proc},
non devono essere usati e non li tratteremo.

L'argomento \param{buf} viene utilizzato solo con i comandi \const{IPC\_STAT}
e \const{IPC\_SET} nel qual caso esso dovrà puntare ad una struttura
\struct{shmid\_ds} precedentemente allocata, in cui nel primo caso saranno
scritti i dati del segmento di memoria restituiti dalla funzione e da cui, nel
secondo caso, verranno letti i dati da impostare sul segmento.

Una volta che lo si è creato, per utilizzare un segmento di memoria condivisa
l'interfaccia prevede due funzioni, \funcd{shmat} e \func{shmdt}. La prima di
queste serve ad agganciare un segmento al processo chiamante, in modo che
quest'ultimo possa inserirlo nel suo spazio di indirizzi per potervi accedere;
il suo prototipo è:

\begin{funcproto}{
\fhead{sys/types.h} 
\fhead{sys/shm.h}
\fdecl{void *shmat(int shmid, const void *shmaddr, int shmflg)}

\fdesc{Aggancia un segmento di memoria condivisa al processo chiamante.}
}

{La funzione ritorna l'indirizzo del segmento in caso di successo e $-1$ (in
  un cast a \ctyp{void *}) per un errore, nel qual caso \var{errno} assumerà
  uno dei valori:
  \begin{errlist}
    \item[\errcode{EACCES}] il processo non ha i privilegi per accedere al
      segmento nella modalità richiesta.
    \item[\errcode{EINVAL}] si è specificato un identificatore invalido per
      \param{shmid}, o un indirizzo non allineato sul confine di una pagina
      per \param{shmaddr} o il valore \val{NULL} indicando \const{SHM\_REMAP}.
  \end{errlist}
  ed inoltre \errval{ENOMEM} nel suo significato generico.
}  
\end{funcproto}

La funzione inserisce un segmento di memoria condivisa all'interno dello
spazio di indirizzi del processo, in modo che questo possa accedervi
direttamente, la situazione dopo l'esecuzione di \func{shmat} è illustrata in
fig.~\ref{fig:ipc_shmem_layout} (per la comprensione del resto dello schema si
ricordi quanto illustrato al proposito in sez.~\ref{sec:proc_mem_layout}). In
particolare l'indirizzo finale del segmento dati (quello impostato da
\func{brk}, vedi sez.~\ref{sec:proc_mem_alloc}) non viene influenzato.
Si tenga presente infine che la funzione ha successo anche se il segmento è
stato marcato per la cancellazione.

\begin{figure}[!htb]
  \centering \includegraphics[height=10cm]{img/sh_memory_layout}
  \caption{Disposizione dei segmenti di memoria di un processo quando si è
    agganciato un segmento di memoria condivisa.}
  \label{fig:ipc_shmem_layout}
\end{figure}

L'argomento \param{shmaddr} specifica a quale indirizzo\footnote{lo standard
  SVID prevede che l'argomento \param{shmaddr} sia di tipo \ctyp{char *}, così
  come il valore di ritorno della funzione; in Linux è stato così con la
  \acr{libc4} e la \acr{libc5}, con il passaggio alla \acr{glibc} il tipo di
  \param{shmaddr} è divenuto un \ctyp{const void *} e quello del valore di
  ritorno un \ctyp{void *} seguendo POSIX.1-2001.} deve essere associato il
segmento, se il valore specificato è \val{NULL} è il sistema a scegliere
opportunamente un'area di memoria libera (questo è il modo più portabile e
sicuro di usare la funzione).  Altrimenti il kernel aggancia il segmento
all'indirizzo specificato da \param{shmaddr}; questo però può avvenire solo se
l'indirizzo coincide con il limite di una pagina, cioè se è un multiplo esatto
del parametro di sistema \const{SHMLBA}, che in Linux è sempre uguale
\const{PAGE\_SIZE}.

Si tenga presente però che quando si usa \val{NULL} come valore di
\param{shmaddr}, l'indirizzo restituito da \func{shmat} può cambiare da
processo a processo; pertanto se nell'area di memoria condivisa si salvano
anche degli indirizzi, si deve avere cura di usare valori relativi (in genere
riferiti all'indirizzo di partenza del segmento).

L'argomento \param{shmflg} permette di cambiare il comportamento della
funzione; esso va specificato come maschera binaria, i bit utilizzati al
momento sono tre e sono identificati dalle costanti \const{SHM\_RND},
\const{SHM\_RDONLY} e \const{SHM\_REMAP} che vanno combinate con un OR
aritmetico.

Specificando \constd{SHM\_RND} si evita che \func{shmat} ritorni un errore
quando \param{shmaddr} non è allineato ai confini di una pagina. Si può quindi
usare un valore qualunque per \param{shmaddr}, e il segmento verrà comunque
agganciato, ma al più vicino multiplo di \const{SHMLBA}; il nome della
costante sta infatti per \textit{rounded}, e serve per specificare un
indirizzo come arrotondamento.

L'uso di \constd{SHM\_RDONLY} permette di agganciare il segmento in sola
lettura (si ricordi che anche le pagine di memoria hanno dei permessi), in tal
caso un tentativo di scrivere sul segmento comporterà una violazione di
accesso con l'emissione di un segnale di \signal{SIGSEGV}. Il comportamento
usuale di \func{shmat} è quello di agganciare il segmento con l'accesso in
lettura e scrittura (ed il processo deve aver questi permessi in
\var{shm\_perm}), non è prevista la possibilità di agganciare un segmento in
sola scrittura.

Infine \constd{SHM\_REMAP} è una estensione specifica di Linux (quindi non
portabile) che indica che la mappatura del segmento deve rimpiazzare ogni
precedente mappatura esistente nell'intervallo iniziante
all'indirizzo \param{shmaddr} e di dimensione pari alla lunghezza del
segmento. In condizioni normali questo tipo di richiesta fallirebbe con un
errore di \errval{EINVAL}. Ovviamente usando \const{SHM\_REMAP}
l'argomento \param{shmaddr} non può essere nullo.

In caso di successo la funzione \func{shmat} aggiorna anche i seguenti campi
della struttura \struct{shmid\_ds}:
\begin{itemize*}
\item il tempo \var{shm\_atime} dell'ultima operazione di aggancio viene
  impostato al tempo corrente.
\item il \ids{PID} \var{shm\_lpid} dell'ultimo processo che ha operato sul
  segmento viene impostato a quello del processo corrente.
\item il numero \var{shm\_nattch} di processi agganciati al segmento viene
  aumentato di uno.
\end{itemize*}

Come accennato in sez.~\ref{sec:proc_fork} un segmento di memoria condivisa
agganciato ad un processo viene ereditato da un figlio attraverso una
\func{fork}, dato che quest'ultimo riceve una copia dello spazio degli
indirizzi del padre. Invece, dato che attraverso una \func{exec} viene
eseguito un diverso programma con uno spazio di indirizzi completamente
diverso, tutti i segmenti agganciati al processo originario vengono
automaticamente sganciati. Lo stesso avviene all'uscita del processo
attraverso una \func{exit}.

Una volta che un segmento di memoria condivisa non serve più, si può
sganciarlo esplicitamente dal processo usando l'altra funzione
dell'interfaccia, \funcd{shmdt}, il cui prototipo è:

\begin{funcproto}{
\fhead{sys/types.h} 
\fhead{sys/shm.h}
\fdecl{int shmdt(const void *shmaddr)}

\fdesc{Sgancia dal processo un segmento di memoria condivisa.}
}

{La funzione ritorna $0$ in caso di successo e $-1$ per un errore, la funzione
  fallisce solo quando non c'è un segmento agganciato
  all'indirizzo \param{shmaddr}, con \var{errno} che assume il valore
  \errval{EINVAL}.  
}  
\end{funcproto}

La funzione sgancia dallo spazio degli indirizzi del processo un segmento di
memoria condivisa; questo viene identificato con l'indirizzo \param{shmaddr}
restituito dalla precedente chiamata a \func{shmat} con il quale era stato
agganciato al processo.

In caso di successo la funzione aggiorna anche i seguenti campi di
\struct{shmid\_ds}:
\begin{itemize*}
\item il tempo \var{shm\_dtime} dell'ultima operazione di sganciamento viene
  impostato al tempo corrente.
\item il \ids{PID} \var{shm\_lpid} dell'ultimo processo che ha operato sul
  segmento viene impostato a quello del processo corrente.
\item il numero \var{shm\_nattch} di processi agganciati al segmento viene
  decrementato di uno.
\end{itemize*} 
inoltre la regione di indirizzi usata per il segmento di memoria condivisa
viene tolta dallo spazio di indirizzi del processo.

\begin{figure}[!htbp]
  \footnotesize \centering
  \begin{minipage}[c]{\codesamplewidth}
    \includecodesample{listati/SharedMem.c}
  \end{minipage} 
  \normalsize 
  \caption{Il codice delle funzioni che permettono di creare, trovare e
    rimuovere un segmento di memoria condivisa.}
  \label{fig:ipc_sysv_shm_func}
\end{figure}

Come esempio di uso di queste funzioni vediamo come implementare una serie di
funzioni di libreria che ne semplifichino l'uso, automatizzando le operazioni
più comuni; il codice, contenuto nel file \file{SharedMem.c}, è riportato in
fig.~\ref{fig:ipc_sysv_shm_func}.

La prima funzione (\texttt{\small 1-16}) è \func{ShmCreate} che, data una
chiave, crea il segmento di memoria condivisa restituendo il puntatore allo
stesso. La funzione comincia (\texttt{\small 6}) con il chiamare
\func{shmget}, usando il flag \const{IPC\_CREATE} per creare il segmento
qualora non esista, ed assegnandogli i privilegi specificati dall'argomento
\var{perm} e la dimensione specificata dall'argomento \var{shm\_size}.  In
caso di errore (\texttt{\small 7-9}) si ritorna immediatamente un puntatore
nullo, altrimenti (\texttt{\small 10}) si prosegue agganciando il segmento di
memoria condivisa al processo con \func{shmat}. In caso di errore
(\texttt{\small 11-13}) si restituisce di nuovo un puntatore nullo, infine
(\texttt{\small 14}) si inizializza con \func{memset} il contenuto del
segmento al valore costante specificato dall'argomento \var{fill}, e poi si
ritorna il puntatore al segmento stesso.

La seconda funzione (\texttt{\small 17-31}) è \func{ShmFind}, che, data una
chiave, restituisce l'indirizzo del segmento ad essa associato. Anzitutto
(\texttt{\small 22}) si richiede l'identificatore del segmento con
\func{shmget}, ritornando (\texttt{\small 23-25}) un puntatore nullo in caso
di errore. Poi si prosegue (\texttt{\small 26}) agganciando il segmento al
processo con \func{shmat}, restituendo (\texttt{\small 27-29}) di nuovo un
puntatore nullo in caso di errore, se invece non ci sono errori si restituisce
il puntatore ottenuto da \func{shmat}.

La terza funzione (\texttt{\small 32-51}) è \func{ShmRemove} che, data la
chiave ed il puntatore associati al segmento di memoria condivisa, prima lo
sgancia dal processo e poi lo rimuove. Il primo passo (\texttt{\small 37}) è
la chiamata a \func{shmdt} per sganciare il segmento, restituendo
(\texttt{\small 38-39}) un valore -1 in caso di errore. Il passo successivo
(\texttt{\small 41}) è utilizzare \func{shmget} per ottenere l'identificatore
associato al segmento data la chiave \var{key}. Al solito si restituisce un
valore di -1 (\texttt{\small 42-45}) in caso di errore, mentre se tutto va
bene si conclude restituendo un valore nullo.

Benché la memoria condivisa costituisca il meccanismo di intercomunicazione
fra processi più veloce, essa non è sempre il più appropriato, dato che, come
abbiamo visto, si avrà comunque la necessità di una sincronizzazione degli
accessi.  Per questo motivo, quando la comunicazione fra processi è
sequenziale, altri meccanismi come le \textit{pipe}, le \textit{fifo} o i
socket, che non necessitano di sincronizzazione esplicita, sono da
preferire. Essa diventa l'unico meccanismo possibile quando la comunicazione
non è sequenziale\footnote{come accennato in sez.~\ref{sec:ipc_sysv_mq} per la
  comunicazione non sequenziale si possono usare le code di messaggi,
  attraverso l'uso del campo \var{mtype}, ma solo se quest'ultima può essere
  effettuata in forma di messaggio.} o quando non può avvenire secondo una
modalità predefinita.

Un esempio classico di uso della memoria condivisa è quello del
``\textit{monitor}'', in cui viene per scambiare informazioni fra un processo
server, che vi scrive dei dati di interesse generale che ha ottenuto, e i
processi client interessati agli stessi dati che così possono leggerli in
maniera completamente asincrona.  Con questo schema di funzionamento da una
parte si evita che ciascun processo client debba compiere l'operazione,
potenzialmente onerosa, di ricavare e trattare i dati, e dall'altra si evita
al processo server di dover gestire l'invio a tutti i client di tutti i dati
(non potendo il server sapere quali di essi servono effettivamente al singolo
client).

Nel nostro caso implementeremo un ``\textsl{monitor}'' di una directory: un
processo si incaricherà di tenere sotto controllo alcuni parametri relativi ad
una directory (il numero dei file contenuti, la dimensione totale, quante
directory, link simbolici, file normali, ecc.) che saranno salvati in un
segmento di memoria condivisa cui altri processi potranno accedere per
ricavare la parte di informazione che interessa.

In fig.~\ref{fig:ipc_dirmonitor_main} si è riportata la sezione principale del
corpo del programma server, insieme alle definizioni delle altre funzioni
usate nel programma e delle variabili globali, omettendo tutto quello che
riguarda la gestione delle opzioni e la stampa delle istruzioni di uso a
video; al solito il codice completo si trova con i sorgenti allegati nel file
\file{DirMonitor.c}.

\begin{figure}[!htbp]
  \footnotesize \centering
  \begin{minipage}[c]{\codesamplewidth}
    \includecodesample{listati/DirMonitor.c}
  \end{minipage} 
  \normalsize 
  \caption{Codice della funzione principale del programma \file{DirMonitor.c}.}
  \label{fig:ipc_dirmonitor_main}
\end{figure}

Il programma usa delle variabili globali (\texttt{\small 2-14}) per mantenere
i valori relativi agli oggetti usati per la comunicazione inter-processo; si è
definita inoltre una apposita struttura \struct{DirProp} che contiene i dati
relativi alle proprietà che si vogliono mantenere nella memoria condivisa, per
l'accesso da parte dei client.

Il programma, dopo la sezione, omessa, relativa alla gestione delle opzioni da
riga di comando (che si limitano alla eventuale stampa di un messaggio di
aiuto a video ed all'impostazione della durata dell'intervallo con cui viene
ripetuto il calcolo delle proprietà della directory) controlla (\texttt{\small
  20-23}) che sia stato specificato l'argomento necessario contenente il nome
della directory da tenere sotto controllo, senza il quale esce immediatamente
con un messaggio di errore.

Poi, per verificare che l'argomento specifichi effettivamente una directory,
si esegue (\texttt{\small 24-26}) su di esso una \func{chdir}, uscendo
immediatamente in caso di errore.  Questa funzione serve anche per impostare
la directory di lavoro del programma nella directory da tenere sotto
controllo, in vista del successivo uso della funzione \func{daemon}. Si noti
come si è potuta fare questa scelta, nonostante le indicazioni illustrate in
sez.~\ref{sec:sess_daemon}, per il particolare scopo del programma, che
necessita comunque di restare all'interno di una directory.

Infine (\texttt{\small 27-29}) si installano i gestori per i vari segnali di
terminazione che, avendo a che fare con un programma che deve essere eseguito
come server, sono il solo strumento disponibile per concluderne l'esecuzione.

Il passo successivo (\texttt{\small 30-39}) è quello di creare gli oggetti di
intercomunicazione necessari. Si inizia costruendo (\texttt{\small 30}) la
chiave da usare come riferimento con il nome del programma,\footnote{si è
  usato un riferimento relativo alla home dell'utente, supposto che i sorgenti
  di GaPiL siano stati installati direttamente in essa; qualora si effettui
  una installazione diversa si dovrà correggere il programma.} dopo di che si
richiede (\texttt{\small 31}) la creazione di un segmento di memoria condivisa
con usando la funzione \func{ShmCreate} illustrata in precedenza (una pagina
di memoria è sufficiente per i dati che useremo), uscendo (\texttt{\small
  32-35}) qualora la creazione ed il successivo agganciamento al processo non
abbia successo. Con l'indirizzo \var{shmptr} così ottenuto potremo poi
accedere alla memoria condivisa, che, per come abbiamo lo abbiamo definito,
sarà vista nella forma data da \struct{DirProp}. Infine (\texttt{\small
  36-39}) utilizzando sempre la stessa chiave, si crea, tramite le funzioni
di interfaccia già descritte in sez.~\ref{sec:ipc_sysv_sem}, anche un mutex,
che utilizzeremo per regolare l'accesso alla memoria condivisa.

Completata l'inizializzazione e la creazione degli oggetti di
intercomunicazione il programma entra nel ciclo principale (\texttt{\small
  40-49}) dove vengono eseguite indefinitamente le attività di monitoraggio.
Il primo passo (\texttt{\small 41}) è eseguire \func{daemon} per proseguire
con l'esecuzione in background come si conviene ad un programma demone; si
noti che si è mantenuta, usando un valore non nullo del primo argomento, la
directory di lavoro corrente.  Una volta che il programma è andato in
background l'esecuzione prosegue all'interno di un ciclo infinito
(\texttt{\small 42-48}).

Si inizia (\texttt{\small 43}) bloccando il mutex con \func{MutexLock} per
poter accedere alla memoria condivisa (la funzione si bloccherà
automaticamente se qualche client sta leggendo), poi (\texttt{\small 44}) si
cancellano i valori precedentemente immagazzinati nella memoria condivisa con
\func{memset}, e si esegue (\texttt{\small 45}) un nuovo calcolo degli stessi
utilizzando la funzione \myfunc{dir\_scan}; infine (\texttt{\small 46}) si
sblocca il mutex con \func{MutexUnlock}, e si attende (\texttt{\small 47}) per
il periodo di tempo specificato a riga di comando con l'opzione \code{-p}
usando una \func{sleep}.

Si noti come per il calcolo dei valori da mantenere nella memoria condivisa si
sia usata ancora una volta la funzione \myfunc{dir\_scan}, già utilizzata (e
descritta in dettaglio) in sez.~\ref{sec:file_dir_read}, che ci permette di
effettuare la scansione delle voci della directory, chiamando per ciascuna di
esse la funzione \func{ComputeValues}, che esegue tutti i calcoli necessari.

\begin{figure}[!htbp]
  \footnotesize \centering
  \begin{minipage}[c]{\codesamplewidth}
    \includecodesample{listati/ComputeValues.c}
  \end{minipage} 
  \normalsize 
  \caption{Codice delle funzioni ausiliarie usate da \file{DirMonitor.c}.}
  \label{fig:ipc_dirmonitor_sub}
\end{figure}


Il codice di quest'ultima è riportato in fig.~\ref{fig:ipc_dirmonitor_sub}.
Come si vede la funzione (\texttt{\small 2-16}) è molto semplice e si limita a
chiamare (\texttt{\small 5}) la funzione \func{stat} sul file indicato da
ciascuna voce, per ottenerne i dati, che poi utilizza per incrementare i vari
contatori nella memoria condivisa, cui accede grazie alla variabile globale
\var{shmptr}.

Dato che la funzione è chiamata da \myfunc{dir\_scan}, si è all'interno del
ciclo principale del programma, con un mutex acquisito, perciò non è
necessario effettuare nessun controllo e si può accedere direttamente alla
memoria condivisa usando \var{shmptr} per riempire i campi della struttura
\struct{DirProp}; così prima (\texttt{\small 6-7}) si sommano le dimensioni
dei file ed il loro numero, poi, utilizzando le macro di
tab.~\ref{tab:file_type_macro}, si contano (\texttt{\small 8-14}) quanti ce
ne sono per ciascun tipo.

In fig.~\ref{fig:ipc_dirmonitor_sub} è riportato anche il codice
(\texttt{\small 17-23}) del gestore dei segnali di terminazione, usato per
chiudere il programma. Esso, oltre a provocare l'uscita del programma, si
incarica anche di cancellare tutti gli oggetti di intercomunicazione non più
necessari.  Per questo anzitutto (\texttt{\small 19}) acquisisce il mutex con
\func{MutexLock}, per evitare di operare mentre un client sta ancora leggendo
i dati, dopo di che (\texttt{\small 20}) distacca e rimuove il segmento di
memoria condivisa usando \func{ShmRemove}.  Infine (\texttt{\small 21})
rimuove il mutex con \func{MutexRemove} ed esce (\texttt{\small 22}).

\begin{figure}[!htbp]
  \footnotesize \centering
  \begin{minipage}[c]{\codesamplewidth}
    \includecodesample{listati/ReadMonitor.c}
  \end{minipage} 
  \normalsize 
  \caption{Codice del programma client del monitor delle proprietà di una
    directory, \file{ReadMonitor.c}.}
  \label{fig:ipc_dirmonitor_client}
\end{figure}

Il codice del client usato per leggere le informazioni mantenute nella memoria
condivisa è riportato in fig.~\ref{fig:ipc_dirmonitor_client}. Al solito si è
omessa la sezione di gestione delle opzioni e la funzione che stampa a video
le istruzioni; il codice completo è nei sorgenti allegati, nel file
\file{ReadMonitor.c}.

Una volta conclusa la gestione delle opzioni a riga di comando il programma
rigenera (\texttt{\small 7}) con \func{ftok} la stessa chiave usata dal server
per identificare il segmento di memoria condivisa ed il mutex, poi
(\texttt{\small 8}) richiede con \func{ShmFind} l'indirizzo della memoria
condivisa agganciando al contempo il segmento al processo, Infine
(\texttt{\small 17-20}) con \func{MutexFind} si richiede l'identificatore del
mutex.  Completata l'inizializzazione ed ottenuti i riferimenti agli oggetti
di intercomunicazione necessari viene eseguito il corpo principale del
programma (\texttt{\small 21-33}); si comincia (\texttt{\small 22})
acquisendo il mutex con \func{MutexLock}; qui avviene il blocco del processo
se la memoria condivisa non è disponibile.  Poi (\texttt{\small 23-31}) si
stampano i vari valori mantenuti nella memoria condivisa attraverso l'uso di
\var{shmptr}.  Infine (\texttt{\small 41}) con \func{MutexUnlock} si rilascia
il mutex, prima di uscire.

Verifichiamo allora il funzionamento dei nostri programmi; al solito, usando
le funzioni di libreria occorre definire opportunamente
\code{LD\_LIBRARY\_PATH}; poi si potrà lanciare il server con:
\begin{Console}
[piccardi@gont sources]$ \textbf{./dirmonitor ./}
\end{Console}
%$
ed avendo usato \func{daemon} il comando ritornerà immediatamente. Una volta
che il server è in esecuzione, possiamo passare ad invocare il client per
verificarne i risultati, in tal caso otterremo:
\begin{Console}
[piccardi@gont sources]$ \textbf{./readmon} 
Ci sono 68 file dati
Ci sono 3 directory
Ci sono 0 link
Ci sono 0 fifo
Ci sono 0 socket
Ci sono 0 device a caratteri
Ci sono 0 device a blocchi
Totale  71 file, per 489831 byte
\end{Console}
%$
ed un rapido calcolo (ad esempio con \code{ls -a | wc} per contare i file) ci
permette di verificare che il totale dei file è giusto. Un controllo con
\cmd{ipcs} ci permette inoltre di verificare la presenza di un segmento di
memoria condivisa e di un semaforo:
\begin{Console}
[piccardi@gont sources]$ \textbf{ipcs}
------ Shared Memory Segments --------
key        shmid      owner      perms      bytes      nattch     status      
0xffffffff 54067205   piccardi  666        4096       1                       

------ Semaphore Arrays --------
key        semid      owner      perms      nsems     
0xffffffff 229376     piccardi  666        1         

------ Message Queues --------
key        msqid      owner      perms      used-bytes   messages    
\end{Console}
%$

Se a questo punto aggiungiamo un file, ad esempio con \code{touch prova},
potremo verificare che, passati nel peggiore dei casi almeno 10 secondi (o
l'eventuale altro intervallo impostato per la rilettura dei dati) avremo:
\begin{Console}
[piccardi@gont sources]$ \textbf{./readmon} 
Ci sono 69 file dati
Ci sono 3 directory
Ci sono 0 link
Ci sono 0 fifo
Ci sono 0 socket
Ci sono 0 device a caratteri
Ci sono 0 device a blocchi
Totale  72 file, per 489887 byte
\end{Console}
%$

A questo punto possiamo far uscire il server inviandogli un segnale di
\signal{SIGTERM} con il comando \code{killall dirmonitor}, a questo punto
ripetendo la lettura, otterremo un errore:
\begin{Console}
[piccardi@gont sources]$ \textbf{./readmon} 
Cannot find shared memory: No such file or directory
\end{Console}
%$
e inoltre potremo anche verificare che anche gli oggetti di intercomunicazione
visti in precedenza sono stati regolarmente  cancellati:
\begin{Console}
[piccardi@gont sources]$ \textbf{ipcs}
------ Shared Memory Segments --------
key        shmid      owner      perms      bytes      nattch     status      

------ Semaphore Arrays --------
key        semid      owner      perms      nsems     

------ Message Queues --------
key        msqid      owner      perms      used-bytes   messages    
\end{Console}
%$


%% Per capire meglio il funzionamento delle funzioni facciamo ancora una volta
%% riferimento alle strutture con cui il kernel implementa i segmenti di memoria
%% condivisa; uno schema semplificato della struttura è illustrato in
%% fig.~\ref{fig:ipc_shm_struct}. 

%% \begin{figure}[!htb]
%%   \centering
%%   \includegraphics[width=10cm]{img/shmstruct}
%%    \caption{Schema dell'implementazione dei segmenti di memoria condivisa in
%%     Linux.}
%%   \label{fig:ipc_shm_struct}
%% \end{figure}




\section{Tecniche alternative}
\label{sec:ipc_alternatives}

Come abbiamo detto in sez.~\ref{sec:ipc_sysv_generic}, e ripreso nella
descrizione dei singoli oggetti che ne fan parte, il \textit{SysV-IPC}
presenta numerosi problemi; in \cite{APUE}\footnote{in particolare nel
  capitolo 14.}  Stevens ne effettua una accurata analisi (alcuni dei concetti
sono già stati accennati in precedenza) ed elenca alcune possibili tecniche
alternative, che vogliamo riprendere in questa sezione.


\subsection{Alternative alle code di messaggi}
\label{sec:ipc_mq_alternative}
 
Le code di messaggi sono probabilmente il meno usato degli oggetti del
\textit{SysV-IPC}; esse infatti nacquero principalmente come meccanismo di
comunicazione bidirezionale quando ancora le \textit{pipe} erano
unidirezionali; con la disponibilità di \func{socketpair} (vedi
sez.~\ref{sec:ipc_socketpair}) o utilizzando una coppia di \textit{pipe}, si
può ottenere questo risultato senza incorrere nelle complicazioni introdotte
dal \textit{SysV-IPC}.

In realtà, grazie alla presenza del campo \var{mtype}, le code di messaggi
hanno delle caratteristiche ulteriori, consentendo una classificazione dei
messaggi ed un accesso non rigidamente sequenziale; due caratteristiche che
sono impossibili da ottenere con le \textit{pipe} e i socket di
\func{socketpair}.  A queste esigenze però si può comunque ovviare in maniera
diversa con un uso combinato della memoria condivisa e dei meccanismi di
sincronizzazione, per cui alla fine l'uso delle code di messaggi classiche è
relativamente poco diffuso.

% TODO: trattare qui, se non si trova posto migliore, copy_from_process e
% copy_to_process, introdotte con il kernel 3.2. Vedi
% http://lwn.net/Articles/405346/ e
% http://ozlabs.org/~cyeoh/cma/process_vm_readv.txt 


\subsection{I \textsl{file di lock}}
\label{sec:ipc_file_lock}

\index{file!di~lock|(}

Come illustrato in sez.~\ref{sec:ipc_sysv_sem} i semafori del \textit{SysV-IPC}
presentano una interfaccia inutilmente complessa e con alcuni difetti
strutturali, per questo quando si ha una semplice esigenza di sincronizzazione
per la quale basterebbe un semaforo binario (quello che abbiamo definito come
\textit{mutex}), per indicare la disponibilità o meno di una risorsa, senza la
necessità di un contatore come i semafori, si possono utilizzare metodi
alternativi.

La prima possibilità, utilizzata fin dalle origini di Unix, è quella di usare
dei \textsl{file di lock} (per i quali è stata anche riservata una opportuna
directory, \file{/var/lock}, nella standardizzazione del \textit{Filesystem
  Hierarchy Standard}). Per questo si usa la caratteristica della funzione
\func{open} (illustrata in sez.~\ref{sec:file_open_close}) che
prevede\footnote{questo è quanto dettato dallo standard POSIX.1, ciò non
  toglie che in alcune implementazioni questa tecnica possa non funzionare; in
  particolare per Linux, nel caso di NFS, si è comunque soggetti alla
  possibilità di una \textit{race condition}.} che essa ritorni un errore
quando usata con i flag di \const{O\_CREAT} e \const{O\_EXCL}. In tal modo la
creazione di un \textsl{file di lock} può essere eseguita atomicamente, il
processo che crea il file con successo si può considerare come titolare del
lock (e della risorsa ad esso associata) mentre il rilascio si può eseguire
con una chiamata ad \func{unlink}.

Un esempio dell'uso di questa funzione è mostrato dalle funzioni
\func{LockFile} ed \func{UnlockFile} riportate in fig.~\ref{fig:ipc_file_lock}
(sono contenute in \file{LockFile.c}, un altro dei sorgenti allegati alla
guida) che permettono rispettivamente di creare e rimuovere un \textsl{file di
  lock}. Come si può notare entrambe le funzioni sono elementari; la prima
(\texttt{\small 4-10}) si limita ad aprire il file di lock (\texttt{\small
  9}) nella modalità descritta, mentre la seconda (\texttt{\small 11-17}) lo
cancella con \func{unlink}.

\begin{figure}[!htbp]
  \footnotesize \centering
  \begin{minipage}[c]{\codesamplewidth}
    \includecodesample{listati/LockFile.c}
  \end{minipage} 
  \normalsize 
  \caption{Il codice delle funzioni \func{LockFile} e \func{UnlockFile} che
    permettono di creare e rimuovere un \textsl{file di lock}.}
  \label{fig:ipc_file_lock}
\end{figure}

Uno dei limiti di questa tecnica è che, come abbiamo già accennato in
sez.~\ref{sec:file_open_close}, questo comportamento di \func{open} può non
funzionare (la funzione viene eseguita, ma non è garantita l'atomicità
dell'operazione) se il filesystem su cui si va ad operare è su NFS; in tal
caso si può adottare una tecnica alternativa che prevede l'uso della
\func{link} per creare come \textsl{file di lock} un hard link ad un file
esistente; se il link esiste già e la funzione fallisce, significa che la
risorsa è bloccata e potrà essere sbloccata solo con un \func{unlink},
altrimenti il link è creato ed il lock acquisito; il controllo e l'eventuale
acquisizione sono atomici; la soluzione funziona anche su NFS, ma ha un altro
difetto è che è quello di poterla usare solo se si opera all'interno di uno
stesso filesystem.

In generale comunque l'uso di un \textsl{file di lock} presenta parecchi
problemi che non lo rendono una alternativa praticabile per la
sincronizzazione: anzitutto in caso di terminazione imprevista del processo,
si lascia allocata la risorsa (il \textsl{file di lock}) e questa deve essere
sempre cancellata esplicitamente.  Inoltre il controllo della disponibilità
può essere eseguito solo con una tecnica di \textit{polling}, ed è quindi
molto inefficiente.

La tecnica dei file di lock ha comunque una sua utilità, e può essere usata
con successo quando l'esigenza è solo quella di segnalare l'occupazione di una
risorsa, senza necessità di attendere che questa si liberi; ad esempio la si
usa spesso per evitare interferenze sull'uso delle porte seriali da parte di
più programmi: qualora si trovi un file di lock il programma che cerca di
accedere alla seriale si limita a segnalare che la risorsa non è disponibile.

\index{file!di~lock|)}


\subsection{La sincronizzazione con il \textit{file locking}}
\label{sec:ipc_lock_file}

Dato che i file di lock presentano gli inconvenienti illustrati in precedenza,
la tecnica alternativa di sincronizzazione più comune è quella di fare ricorso
al \textit{file locking} (trattato in sez.~\ref{sec:file_locking}) usando
\func{fcntl} su un file creato per l'occasione per ottenere un write lock. In
questo modo potremo usare il lock come un \textit{mutex}: per bloccare la
risorsa basterà acquisire il lock, per sbloccarla basterà rilasciare il
lock. Una richiesta fatta con un write lock metterà automaticamente il
processo in stato di attesa, senza necessità di ricorrere al \textit{polling}
per determinare la disponibilità della risorsa, e al rilascio della stessa da
parte del processo che la occupava si otterrà il nuovo lock atomicamente.

Questo approccio presenta il notevole vantaggio che alla terminazione di un
processo tutti i lock acquisiti vengono rilasciati automaticamente (alla
chiusura dei relativi file) e non ci si deve preoccupare di niente; inoltre
non consuma risorse permanentemente allocate nel sistema. Lo svantaggio è che,
dovendo fare ricorso a delle operazioni sul filesystem, esso è in genere
leggermente più lento.

\begin{figure}[!htbp]
  \footnotesize \centering
  \begin{minipage}[c]{\codesamplewidth}
    \includecodesample{listati/MutexLocking.c}
  \end{minipage} 
  \normalsize 
  \caption{Il codice delle funzioni che permettono per la gestione dei
    \textit{mutex} con il \textit{file locking}.}
  \label{fig:ipc_flock_mutex}
\end{figure}

Il codice delle varie funzioni usate per implementare un mutex utilizzando il
\textit{file locking} è riportato in fig.~\ref{fig:ipc_flock_mutex}; si è
mantenuta volutamente una struttura analoga alle precedenti funzioni che usano
i semafori, anche se le due interfacce non possono essere completamente
equivalenti, specie per quanto riguarda la rimozione del mutex.

La prima funzione (\texttt{\small 1-5}) è \func{CreateMutex}, e serve a
creare il mutex; la funzione è estremamente semplice, e si limita
(\texttt{\small 4}) a creare, con una opportuna chiamata ad \func{open}, il
file che sarà usato per il successivo \textit{file locking}, assicurandosi che
non esista già (nel qual caso segnala un errore); poi restituisce il file
descriptor che sarà usato dalle altre funzioni per acquisire e rilasciare il
mutex.

La seconda funzione (\texttt{\small 6-10}) è \func{FindMutex}, che, come la
precedente, è stata definita per mantenere una analogia con la corrispondente
funzione basata sui semafori. Anch'essa si limita (\texttt{\small 9}) ad
aprire il file da usare per il \textit{file locking}, solo che in questo caso
le opzioni di \func{open} sono tali che il file in questione deve esistere di
già.

La terza funzione (\texttt{\small 11-22}) è \func{LockMutex} e serve per
acquisire il mutex. La funzione definisce (\texttt{\small 14}) e inizializza
(\texttt{\small 16-19}) la struttura \var{lock} da usare per acquisire un
write lock sul file, che poi (\texttt{\small 21}) viene richiesto con
\func{fcntl}, restituendo il valore di ritorno di quest'ultima. Se il file è
libero il lock viene acquisito e la funzione ritorna immediatamente;
altrimenti \func{fcntl} si bloccherà (si noti che la si è chiamata con
\const{F\_SETLKW}) fino al rilascio del lock.

La quarta funzione (\texttt{\small 24-34}) è \func{UnlockMutex} e serve a
rilasciare il mutex. La funzione è analoga alla precedente, solo che in questo
caso si inizializza (\texttt{\small 28-31}) la struttura \var{lock} per il
rilascio del lock, che viene effettuato (\texttt{\small 33}) con la opportuna
chiamata a \func{fcntl}. Avendo usato il \textit{file locking} in semantica
POSIX (si riveda quanto detto sez.~\ref{sec:file_posix_lock}) solo il processo
che ha precedentemente eseguito il lock può sbloccare il mutex.

La quinta funzione (\texttt{\small 36-39}) è \func{RemoveMutex} e serve a
cancellare il mutex. Anche questa funzione è stata definita per mantenere una
analogia con le funzioni basate sui semafori, e si limita a cancellare
(\texttt{\small 38}) il file con una chiamata ad \func{unlink}. Si noti che in
questo caso la funzione non ha effetto sui mutex già ottenuti con precedenti
chiamate a \func{FindMutex} o \func{CreateMutex}, che continueranno ad essere
disponibili fintanto che i relativi file descriptor restano aperti. Pertanto
per rilasciare un mutex occorrerà prima chiamare \func{UnlockMutex} oppure
chiudere il file usato per il lock.

La sesta funzione (\texttt{\small 41-55}) è \func{ReadMutex} e serve a
leggere lo stato del mutex. In questo caso si prepara (\texttt{\small 46-49})
la solita struttura \var{lock} come l'acquisizione del lock, ma si effettua
(\texttt{\small 51}) la chiamata a \func{fcntl} usando il comando
\const{F\_GETLK} per ottenere lo stato del lock, e si restituisce
(\texttt{\small 52}) il valore di ritorno in caso di errore, ed il valore del
campo \var{l\_type} (che descrive lo stato del lock) altrimenti
(\texttt{\small 54}). Per questo motivo la funzione restituirà -1 in caso di
errore e uno dei due valori \const{F\_UNLCK} o \const{F\_WRLCK}\footnote{non
  si dovrebbe mai avere il terzo valore possibile, \const{F\_RDLCK}, dato che
  la nostra interfaccia usa solo i write lock. Però è sempre possibile che
  siano richiesti altri lock sul file al di fuori dell'interfaccia, nel qual
  caso si potranno avere, ovviamente, interferenze indesiderate.} in caso di
successo, ad indicare che il mutex è, rispettivamente, libero o occupato.

Basandosi sulla semantica dei file lock POSIX valgono tutte le considerazioni
relative al comportamento di questi ultimi fatte in
sez.~\ref{sec:file_posix_lock}; questo significa ad esempio che, al contrario
di quanto avveniva con l'interfaccia basata sui semafori, chiamate multiple a
\func{UnlockMutex} o \func{LockMutex} non si cumulano e non danno perciò
nessun inconveniente.


\subsection{Il \textit{memory mapping} anonimo}
\label{sec:ipc_mmap_anonymous}

\itindbeg{memory~mapping} 

Abbiamo già visto che quando i processi sono \textsl{correlati}, se cioè hanno
almeno un progenitore comune, l'uso delle \textit{pipe} può costituire una
valida alternativa alle code di messaggi; nella stessa situazione si può
evitare l'uso di una memoria condivisa facendo ricorso al cosiddetto
\textit{memory mapping} anonimo.

In sez.~\ref{sec:file_memory_map} abbiamo visto come sia possibile mappare il
contenuto di un file nella memoria di un processo, e che, quando viene usato
il flag \const{MAP\_SHARED}, le modifiche effettuate al contenuto del file
vengono viste da tutti i processi che lo hanno mappato. Utilizzare questa
tecnica per creare una memoria condivisa fra processi diversi è estremamente
inefficiente, in quanto occorre passare attraverso il disco. 

Però abbiamo visto anche che se si esegue la mappatura con il flag
\const{MAP\_ANONYMOUS} la regione mappata non viene associata a nessun file,
anche se quanto scritto rimane in memoria e può essere riletto; allora, dato
che un processo figlio mantiene nel suo spazio degli indirizzi anche le
regioni mappate, esso sarà anche in grado di accedere a quanto in esse è
contenuto.

In questo modo diventa possibile creare una memoria condivisa fra processi
diversi, purché questi abbiano almeno un progenitore comune che ha effettuato
il \textit{memory mapping} anonimo.\footnote{nei sistemi derivati da SysV una
  funzionalità simile a questa viene implementata mappando il file speciale
  \file{/dev/zero}. In tal caso i valori scritti nella regione mappata non
  vengono ignorati (come accade qualora si scriva direttamente sul file), ma
  restano in memoria e possono essere riletti secondo le stesse modalità usate
  nel \textit{memory mapping} anonimo.} Vedremo come utilizzare questa tecnica
più avanti, quando realizzeremo una nuova versione del monitor visto in
sez.~\ref{sec:ipc_sysv_shm} che possa restituisca i risultati via rete.

\itindend{memory~mapping}

% TODO: fare esempio di mmap anonima

% TODO: con il kernel 3.2 è stata introdotta un nuovo meccanismo di
% intercomunicazione veloce chiamato Cross Memory Attach, da capire se e come
% trattarlo qui, vedi http://lwn.net/Articles/405346/
% https://git.kernel.org/?p=linux/kernel/git/torvalds/linux-2.6.git;a=commitdiff;h=fcf634098c00dd9cd247447368495f0b79be12d1

% TODO: con il kernel 3.17 è stata introdotta una fuunzionalità di
% sigillatura dei file mappati in memoria e la system call memfd
% (capire se va messo qui o altrove) vedi: http://lwn.net/Articles/593918/
% col 5.1 aggiunta a memfd F_SEAL_FUTURE_WRITE, vedi 
% https://git.kernel.org/linus/ab3948f58ff8 e https://lwn.net/Articles/782511/


\section{L'intercomunicazione fra processi di POSIX}
\label{sec:ipc_posix}

Per superare i numerosi problemi del \textit{SysV-IPC}, evidenziati per i suoi
aspetti generali in coda a sez.~\ref{sec:ipc_sysv_generic} e per i singoli
oggetti nei paragrafi successivi, lo standard POSIX.1b ha introdotto dei nuovi
meccanismi di comunicazione, che vanno sotto il nome di POSIX IPC, definendo
una interfaccia completamente nuova, che tratteremo in questa sezione.


\subsection{Considerazioni generali}
\label{sec:ipc_posix_generic}

Oggi Linux supporta tutti gli oggetti definito nello standard POSIX per l'IPC,
ma a lungo non è stato così; la memoria condivisa è presente a partire dal
kernel 2.4.x, i semafori sono forniti dalla \acr{glibc} nella sezione che
implementa i \textit{thread} POSIX di nuova generazione che richiedono il
kernel 2.6, le code di messaggi sono supportate a partire dal kernel 2.6.6.

La caratteristica fondamentale dell'interfaccia POSIX è l'abbandono dell'uso
degli identificatori e delle chiavi visti nel \textit{SysV-IPC}, per passare ai
\itindex{POSIX~IPC~names} \textit{POSIX IPC names}, che sono sostanzialmente
equivalenti ai nomi dei file. Tutte le funzioni che creano un oggetto di IPC
POSIX prendono come primo argomento una stringa che indica uno di questi nomi;
lo standard è molto generico riguardo l'implementazione, ed i nomi stessi
possono avere o meno una corrispondenza sul filesystem; tutto quello che è
richiesto è che:
\begin{itemize*}
\item i nomi devono essere conformi alle regole che caratterizzano i
  \textit{pathname}, in particolare non essere più lunghi di \const{PATH\_MAX}
  byte e terminati da un carattere nullo.
\item se il nome inizia per una \texttt{/} chiamate differenti allo stesso
  nome fanno riferimento allo stesso oggetto, altrimenti l'interpretazione del
  nome dipende dall'implementazione.
\item l'interpretazione di ulteriori \texttt{/} presenti nel nome dipende
  dall'implementazione.
\end{itemize*}

Data la assoluta genericità delle specifiche, il comportamento delle funzioni
è subordinato in maniera quasi completa alla relativa implementazione, tanto
che Stevens in \cite{UNP2} cita questo caso come un esempio della maniera
standard usata dallo standard POSIX per consentire implementazioni non
standardizzabili. 

Nel caso di Linux, sia per quanto riguarda la memoria condivisa ed i semafori,
che per le code di messaggi, tutto viene creato usando come radici delle
opportune directory (rispettivamente \file{/dev/shm} e \file{/dev/mqueue}, per
i dettagli si faccia riferimento a sez.~\ref{sec:ipc_posix_shm},
sez.~\ref{sec:ipc_posix_sem} e sez.~\ref{sec:ipc_posix_mq}).  I nomi
specificati nelle relative funzioni devono essere nella forma di un
\textit{pathname} assoluto (devono cioè iniziare con ``\texttt{/}'') e
corrisponderanno ad altrettanti file creati all'interno di queste directory;
per questo motivo detti nomi non possono contenere altre ``\texttt{/}'' oltre
quella iniziale.

Il vantaggio degli oggetti di IPC POSIX è comunque che essi vengono inseriti
nell'albero dei file, e possono essere maneggiati con le usuali funzioni e
comandi di accesso ai file, che funzionano come su dei file normali. Questo
però è vero nel caso di Linux, che usa una implementazione che lo consente,
non è detto che altrettanto valga per altri kernel. In particolare, come si
può facilmente verificare con il comando \cmd{strace}, sia per la memoria
condivisa che per le code di messaggi varie \textit{system call} utilizzate da
Linux corrispondono in realtà a quelle ordinarie dei file, essendo detti
oggetti realizzati come tali usando degli specifici filesystem.

In particolare i permessi associati agli oggetti di IPC POSIX sono identici ai
permessi dei file, ed il controllo di accesso segue esattamente la stessa
semantica (quella illustrata in sez.~\ref{sec:file_access_control}), e non
quella particolare (si ricordi quanto visto in
sez.~\ref{sec:ipc_sysv_access_control}) che viene usata per gli oggetti del
SysV-IPC. Per quanto riguarda l'attribuzione dell'utente e del gruppo
proprietari dell'oggetto alla creazione di quest'ultimo essa viene effettuata
secondo la semantica SysV: corrispondono cioè a \ids{UID} e \ids{GID} effettivi
del processo che esegue la creazione.


\subsection{Code di messaggi Posix}
\label{sec:ipc_posix_mq}

Le code di messaggi POSIX sono supportate da Linux a partire dalla versione
2.6.6 del kernel. In generale, come le corrispettive del \textit{SysV-IPC}, le
code di messaggi sono poco usate, dato che i socket, nei casi in cui sono
sufficienti, sono più comodi, e che in casi più complessi la comunicazione può
essere gestita direttamente con mutex (o semafori) e memoria condivisa con
tutta la flessibilità che occorre.

Per poter utilizzare le code di messaggi, oltre ad utilizzare un kernel
superiore al 2.6.6 occorre utilizzare la libreria \file{librt} che contiene le
funzioni dell'interfaccia POSIX ed i programmi che usano le code di messaggi
devono essere compilati aggiungendo l'opzione \code{-lrt} al comando
\cmd{gcc}. In corrispondenza all'inclusione del supporto nel kernel ufficiale
le funzioni di libreria sono state inserite nella \acr{glibc}, e sono
disponibili a partire dalla versione 2.3.4 delle medesime.

La libreria inoltre richiede la presenza dell'apposito filesystem di tipo
\texttt{mqueue} montato sulla directory \file{/dev/mqueue}; questo può essere
fatto aggiungendo ad \conffile{/etc/fstab} una riga come:
\begin{FileExample}[label=/etc/fstab]
mqueue   /dev/mqueue       mqueue    defaults        0      0
\end{FileExample}
ed esso sarà utilizzato come radice sulla quale vengono risolti i nomi delle
code di messaggi che iniziano con una ``\texttt{/}''. Le opzioni di mount
accettate sono \texttt{uid}, \texttt{gid} e \texttt{mode} che permettono
rispettivamente di impostare l'utente, il gruppo ed i permessi associati al
filesystem.


La funzione di sistema che permette di aprire (e crearla se non esiste ancora)
una coda di messaggi POSIX è \funcd{mq\_open}, ed il suo prototipo è:

\begin{funcproto}{
\fhead{fcntl.h}
\fhead{sys/stat.h}
\fhead{mqueue.h}
\fdecl{mqd\_t mq\_open(const char *name, int oflag)}
\fdecl{mqd\_t mq\_open(const char *name, int oflag, unsigned long mode,
    struct mq\_attr *attr)}

\fdesc{Apre una coda di messaggi POSIX impostandone le caratteristiche.}
}

{La funzione ritorna il descrittore associato alla coda in caso di successo e
  $-1$ per un errore, nel qual caso \var{errno} assumerà uno dei valori:
  \begin{errlist}
  \item[\errcode{EACCES}] il processo non ha i privilegi per accedere alla
    coda secondo quanto specificato da \param{oflag} oppure \const{name}
    contiene più di una ``\texttt{/}''.
  \item[\errcode{EEXIST}] si è specificato \const{O\_CREAT} e \const{O\_EXCL}
    ma la coda già esiste.
  \item[\errcode{EINVAL}] il file non supporta la funzione, o si è specificato
    \const{O\_CREAT} con una valore non nullo di \param{attr} e valori non
    validi dei campi \var{mq\_maxmsg} e \var{mq\_msgsize}; questi valori
    devono essere positivi ed inferiori ai limiti di sistema se il processo
    non ha privilegi amministrativi, inoltre \var{mq\_maxmsg} non può comunque
    superare \const{HARD\_MAX}.
  \item[\errcode{ENOENT}] non si è specificato \const{O\_CREAT} ma la coda non
    esiste o si è usato il nome ``\texttt{/}''.
  \item[\errcode{ENOSPC}] lo spazio è insufficiente, probabilmente per aver
    superato il limite di \texttt{queues\_max}.
  \end{errlist}
  ed inoltre \errval{EMFILE}, \errval{ENAMETOOLONG}, \errval{ENFILE},
  \errval{ENOMEM} ed nel loro significato generico.  }
\end{funcproto}

La funzione apre la coda di messaggi identificata dall'argomento \param{name}
restituendo il descrittore ad essa associato, del tutto analogo ad un file
descriptor, con l'unica differenza che lo standard prevede un apposito tipo
\typed{mqd\_t}. Nel caso di Linux si tratta in effetti proprio di un normale
file descriptor; pertanto, anche se questo comportamento non è portabile, lo
si può tenere sotto osservazione con le funzioni dell'I/O multiplexing (vedi
sez.~\ref{sec:file_multiplexing}) come possibile alternativa all'uso
dell'interfaccia di notifica di \func{mq\_notify} (che vedremo a breve).

Se il nome indicato fa riferimento ad una coda di messaggi già esistente, il
descrittore ottenuto farà riferimento allo stesso oggetto, pertanto tutti i
processi che hanno usato \func{mq\_open} su quel nome otterranno un
riferimento alla stessa coda. Diventa così immediato costruire un canale di
comunicazione fra detti processi.

La funzione è del tutto analoga ad \func{open} ed analoghi sono i valori che
possono essere specificati per \param{oflag}, che deve essere specificato come
maschera binaria; i valori possibili per i vari bit sono quelli visti in
sez.~\ref{sec:file_open_close} (per questo occorre includere \texttt{fcntl.h})
dei quali però \func{mq\_open} riconosce solo i seguenti:
\begin{basedescript}{\desclabelwidth{2.2cm}\desclabelstyle{\nextlinelabel}}
\item[\const{O\_RDONLY}] Apre la coda solo per la ricezione di messaggi. Il
  processo potrà usare il descrittore con \func{mq\_receive} ma non con
  \func{mq\_send}.
\item[\const{O\_WRONLY}] Apre la coda solo per la trasmissione di messaggi. Il
  processo potrà usare il descrittore con \func{mq\_send} ma non con
  \func{mq\_receive}.
\item[\const{O\_RDWR}] Apre la coda solo sia per la trasmissione che per la
  ricezione. 
\item[\const{O\_CREAT}] Necessario qualora si debba creare la coda; la
  presenza di questo bit richiede la presenza degli ulteriori argomenti
  \param{mode} e \param{attr}.
\item[\const{O\_EXCL}] Se usato insieme a \const{O\_CREAT} fa fallire la
  chiamata se la coda esiste già, altrimenti esegue la creazione atomicamente.
\item[\const{O\_NONBLOCK}] Imposta la coda in modalità non bloccante, le
  funzioni di ricezione e trasmissione non si bloccano quando non ci sono le
  risorse richieste, ma ritornano immediatamente con un errore di
  \errcode{EAGAIN}.
\end{basedescript}

I primi tre bit specificano la modalità di apertura della coda, e sono fra
loro esclusivi. Ma qualunque sia la modalità in cui si è aperta una coda,
questa potrà essere riaperta più volte in una modalità diversa, e vi si potrà
sempre accedere attraverso descrittori diversi, esattamente come si può fare
per i file normali.

Se la coda non esiste e la si vuole creare si deve specificare
\const{O\_CREAT}, in tal caso occorre anche specificare i permessi di
creazione con l'argomento \param{mode};\footnote{fino al 2.6.14 per un bug i
  valori della \textit{umask} del processo non venivano applicati a questi
  permessi.} i valori di quest'ultimo sono identici a quelli usati per
\func{open} (per questo occorre includere \texttt{sys/stat.h}), anche se per
le code di messaggi han senso solo i permessi di lettura e scrittura.

Oltre ai permessi di creazione possono essere specificati anche gli attributi
specifici della coda tramite l'argomento \param{attr}; quest'ultimo è un
puntatore ad una apposita struttura \struct{mq\_attr}, la cui definizione è
riportata in fig.~\ref{fig:ipc_mq_attr}.

\begin{figure}[!htb]
  \footnotesize \centering
  \begin{minipage}[c]{0.90\textwidth}
    \includestruct{listati/mq_attr.h}
  \end{minipage} 
  \normalsize
  \caption{La struttura \structd{mq\_attr}, contenente gli attributi di una
    coda di messaggi POSIX.}
  \label{fig:ipc_mq_attr}
\end{figure}

Per la creazione della coda i campi della struttura che devono essere
specificati sono \var{mq\_maxmsg} e \var{mq\_msgsize}, che indicano
rispettivamente il numero massimo di messaggi che può contenere e la
dimensione massima di un messaggio. Il valore dovrà essere positivo e minore
dei rispettivi limiti di sistema altrimenti la funzione fallirà con un errore
di \errcode{EINVAL}.  Se \param{attr} è un puntatore nullo gli attributi della
coda saranno impostati ai valori predefiniti.

I suddetti limiti di sistema sono impostati attraverso una serie di file
presenti sotto \texttt{/proc/sys/fs/mqueue}, in particolare i file che
controllano i valori dei limiti sono:
\begin{basedescript}{\desclabelwidth{1.5cm}\desclabelstyle{\nextlinelabel}}
\item[\sysctlfiled{fs/mqueue/msg\_max}] Indica il valore massimo del numero di
  messaggi in una coda e agisce come limite superiore per il valore di
  \var{attr->mq\_maxmsg} in \func{mq\_open}. Il suo valore di default è 10. Il
  valore massimo è \constd{HARD\_MAX} che vale \code{(131072/sizeof(void *))},
  ed il valore minimo 1 (ma era 10 per i kernel precedenti il 2.6.28). Questo
  limite viene ignorato per i processi con privilegi amministrativi (più
  precisamente con la \textit{capability} \const{CAP\_SYS\_RESOURCE}) ma
  \const{HARD\_MAX} resta comunque non superabile.

\item[\sysctlfiled{fs/mqueue/msgsize\_max}] Indica il valore massimo della
  dimensione in byte di un messaggio sulla coda ed agisce come limite
  superiore per il valore di \var{attr->mq\_msgsize} in \func{mq\_open}. Il
  suo valore di default è 8192.  Il valore massimo è 1048576 ed il valore
  minimo 128 (ma per i kernel precedenti il 2.6.28 detti limiti erano
  rispettivamente \const{INT\_MAX} e 8192). Questo limite viene ignorato dai
  processi con privilegi amministrativi (con la \textit{capability}
  \const{CAP\_SYS\_RESOURCE}).

\item[\sysctlfiled{fs/mqueue/queues\_max}] Indica il numero massimo di code di
  messaggi creabili in totale sul sistema, il valore di default è 256 ma si
  può usare un valore qualunque fra $0$ e \const{INT\_MAX}. Il limite non
  viene applicato ai processi con privilegi amministrativi (cioè con la
  \textit{capability} \const{CAP\_SYS\_RESOURCE}).

\end{basedescript}

Infine sulle code di messaggi si applica il limite imposto sulla risorsa
\const{RLIMIT\_MSGQUEUE} (vedi sez.~\ref{sec:sys_resource_limit}) che indica
lo spazio massimo (in byte) occupabile dall'insieme di tutte le code di
messaggi appartenenti ai processi di uno stesso utente, che viene identificato
in base al \textit{real user ID} degli stessi.

Quando l'accesso alla coda non è più necessario si può chiudere il relativo
descrittore con la funzione \funcd{mq\_close}, il cui prototipo è:

\begin{funcproto}{
\fhead{mqueue.h}
\fdecl{int mq\_close(mqd\_t mqdes)}

\fdesc{Chiude una coda di messaggi.}
}

{La funzione ritorna $0$ in caso di successo e $-1$ per un errore, nel qual
  caso \var{errno} assumerà uno dei valori \errval{EBADF} o \errval{EINTR} nel
  loro significato generico.
}  
\end{funcproto}

La funzione è analoga a \func{close},\footnote{su Linux, dove le code sono
  implementate come file su un filesystem dedicato, è esattamente la stessa
  funzione, per cui non esiste una \textit{system call} autonoma e la funzione
  viene rimappata su \func{close} dalla \acr{glibc}.}  dopo la sua esecuzione
il processo non sarà più in grado di usare il descrittore della coda, ma
quest'ultima continuerà ad esistere nel sistema e potrà essere acceduta con
un'altra chiamata a \func{mq\_open}. All'uscita di un processo tutte le code
aperte, così come i file, vengono chiuse automaticamente. Inoltre se il
processo aveva agganciato una richiesta di notifica sul descrittore che viene
chiuso, questa sarà rilasciata e potrà essere richiesta da qualche altro
processo.

Quando si vuole effettivamente rimuovere una coda dal sistema occorre usare la
funzione di sistema \funcd{mq\_unlink}, il cui prototipo è:

\begin{funcproto}{
\fhead{mqueue.h}
\fdecl{int mq\_unlink(const char *name)}

\fdesc{Rimuove una coda di messaggi.}
}

{La funzione ritorna $0$ in caso di successo e $-1$ per un errore, nel qual
  caso \var{errno} assumerà gli uno dei valori:
  \begin{errlist}
    \item[\errcode{EACCES}] non si hanno i permessi per cancellare la coda.
    \item[\errcode{ENAMETOOLONG}] il nome indicato è troppo lungo.
    \item[\errcode{ENOENT}] non esiste una coda con il nome indicato.
  \end{errlist}
}  
\end{funcproto}

Anche in questo caso il comportamento della funzione è analogo a quello di
\func{unlink} per i file, la funzione rimuove la coda \param{name} (ed il
relativo file sotto \texttt{/dev/mqueue}), così che una successiva chiamata a
\func{mq\_open} fallisce o crea una coda diversa.

% TODO, verificare se mq_unlink è davvero una system call indipendente.

Come per i file ogni coda di messaggi ha un contatore di riferimenti, per cui
la coda non viene effettivamente rimossa dal sistema fin quando questo non si
annulla. Pertanto anche dopo aver eseguito con successo \func{mq\_unlink} la
coda resterà accessibile a tutti i processi che hanno un descrittore aperto su
di essa.  Allo stesso modo una coda ed i suoi contenuti resteranno disponibili
all'interno del sistema anche quando quest'ultima non è aperta da nessun
processo (questa è una delle differenze più rilevanti nei confronti di
\textit{pipe} e \textit{fifo}).  La sola differenza fra code di messaggi POSIX
e file normali è che, essendo il filesystem delle code di messaggi virtuale, e
basato su oggetti interni al kernel, il suo contenuto viene perduto con il
riavvio del sistema.

Come accennato ad ogni coda di messaggi è associata una struttura
\struct{mq\_attr}, che può essere letta e modificata attraverso le due
funzioni \funcd{mq\_getattr} e \funcd{mq\_setattr}, i cui prototipi sono:

\begin{funcproto}{
\fhead{mqueue.h}
\fdecl{int mq\_getattr(mqd\_t mqdes, struct mq\_attr *mqstat)}
\fdesc{Legge gli attributi di una coda di messaggi POSIX.}
\fdecl{int mq\_setattr(mqd\_t mqdes, const struct mq\_attr *mqstat,
    struct mq\_attr *omqstat)}
\fdesc{Modifica gli attributi di una coda di messaggi POSIX.}
}
{
Entrambe le funzioni ritornano $0$ in caso di successo e $-1$ per un errore,
  nel qual caso \var{errno} assumerà i valori \errval{EBADF}
    o \errval{EINVAL} nel loro significato generico.
}  
\end{funcproto}

La funzione \func{mq\_getattr} legge i valori correnti degli attributi della
coda \param{mqdes} nella struttura \struct{mq\_attr} puntata
da \param{mqstat}; di questi l'unico relativo allo stato corrente della coda è
\var{mq\_curmsgs} che indica il numero di messaggi da essa contenuti, gli
altri indicano le caratteristiche generali della stessa impostate in fase di
apertura.

La funzione \func{mq\_setattr} permette di modificare gli attributi di una
coda (indicata da \param{mqdes}) tramite i valori contenuti nella struttura
\struct{mq\_attr} puntata da \param{mqstat}, ma può essere modificato solo il
campo \var{mq\_flags}, gli altri campi vengono comunque ignorati.

In particolare i valori di \var{mq\_maxmsg} e \var{mq\_msgsize} possono essere
specificati solo in fase ci creazione della coda.  Inoltre i soli valori
possibili per \var{mq\_flags} sono 0 e \const{O\_NONBLOCK}, per cui alla fine
la funzione può essere utilizzata solo per abilitare o disabilitare la
modalità non bloccante. L'argomento \param{omqstat} viene usato, quando
diverso da \val{NULL}, per specificare l'indirizzo di una struttura su cui
salvare i valori degli attributi precedenti alla chiamata della funzione.

Per inserire messaggi su di una coda sono previste due funzioni di sistema,
\funcd{mq\_send} e \funcd{mq\_timedsend}. In realtà su Linux la \textit{system
  call} è soltanto \func{mq\_timedsend}, mentre \func{mq\_send} viene
implementata come funzione di libreria che si appoggia alla
precedente. Inoltre \func{mq\_timedsend} richiede che sia definita la macro
\macro{\_XOPEN\_SOURCE} ad un valore pari ad almeno \texttt{600} o la macro
\macro{\_POSIX\_C\_SOURCE} ad un valore uguale o maggiore di \texttt{200112L}.
I rispettivi prototipi sono:

\begin{funcproto}{
\fhead{mqueue.h}
\fdecl{int mq\_send(mqd\_t mqdes, const char *msg\_ptr, size\_t msg\_len,
    unsigned int msg\_prio)}
\fdesc{Esegue l'inserimento di un messaggio su una coda.}
\fhead{mqueue.h}
\fhead{time.h}
\fdecl{int mq\_timedsend(mqd\_t mqdes, const char *msg\_ptr, size\_t
    msg\_len, \\ 
\phantom{int mq\_timedsend(}unsigned int msg\_prio, const struct timespec
*abs\_timeout)} 
\fdesc{Esegue l'inserimento di un messaggio su una coda entro un tempo
  specificato}
}

{Entrambe le funzioni ritornano $0$ in caso di successo e $-1$ per un errore,
  nel qual caso \var{errno} assumerà uno dei valori:
  \begin{errlist}
    \item[\errcode{EAGAIN}] si è aperta la coda con \const{O\_NONBLOCK}, e la
      coda è piena.
    \item[\errcode{EBADF}] si specificato un file descriptor non valido.
    \item[\errcode{EINVAL}] si è specificato un valore nullo per
      \param{msg\_len}, o un valore di \param{msg\_prio} fuori dai limiti, o
      un valore non valido per \param{abs\_timeout}.
    \item[\errcode{EMSGSIZE}] la lunghezza del messaggio \param{msg\_len}
      eccede il limite impostato per la coda.
    \item[\errcode{ETIMEDOUT}] l'inserimento del messaggio non è stato
      effettuato entro il tempo stabilito (solo \func{mq\_timedsend}).
  \end{errlist}
  ed inoltre \errval{EBADF} e \errval{EINTR} nel loro significato generico.
}
\end{funcproto}

Entrambe le funzioni richiedono un puntatore ad un buffer in memoria
contenente il testo del messaggio da inserire nella coda \param{mqdes}
nell'argomento \param{msg\_ptr}, e la relativa lunghezza in \param{msg\_len}.
Se quest'ultima eccede la dimensione massima specificata da \var{mq\_msgsize}
le funzioni ritornano immediatamente con un errore di \errcode{EMSGSIZE}.

L'argomento \param{msg\_prio} indica la priorità dell'argomento che essendo
definito come \ctyp{unsigned int} è sempre un intero positivo. I messaggi di
priorità maggiore vengono inseriti davanti a quelli di priorità inferiore, e
quindi saranno riletti per primi. A parità del valore della priorità il
messaggio sarà inserito in coda a tutti quelli che hanno la stessa priorità
che quindi saranno letti con la politica di una \textit{fifo}. Il valore della
priorità non può eccedere il limite di sistema \constd{MQ\_PRIO\_MAX}, che al
momento è pari a 32768.

Qualora la coda sia piena, entrambe le funzioni si bloccano, a meno che non
sia stata selezionata in fase di apertura della stessa la modalità non
bloccante o non si sia impostato il flag \const{O\_NONBLOCK} sul file
descriptor della coda, nel qual caso entrambe ritornano con un codice di
errore di \errcode{EAGAIN}.

La sola differenza fra le due funzioni è che \func{mq\_timedsend}, passato il
tempo massimo impostato con l'argomento \param{abs\_timeout}, ritorna con un
errore di \errcode{ETIMEDOUT}, se invece il tempo è già scaduto al momento
della chiamata e la coda è piena la funzione ritorna immediatamente. Il valore
di \param{abs\_timeout} deve essere specificato come tempo assoluto tramite
una struttura \struct{timespec} (vedi fig.~\ref{fig:sys_timespec_struct})
indicato in numero di secondi e nanosecondi a partire dal 1 gennaio 1970.

Come per l'inserimento, anche per l'estrazione dei messaggi da una coda sono
previste due funzioni di sistema, \funcd{mq\_receive} e
\funcd{mq\_timedreceive}. Anche in questo caso su Linux soltanto
\func{mq\_timedreceive} è effettivamente, una \textit{system call} e per
usarla devono essere definite le opportune macro come per
\func{mq\_timedsend}. I rispettivi prototipi sono:

\begin{funcproto}{
\fhead{mqueue.h} 
\fdecl{ssize\_t mq\_receive(mqd\_t mqdes, char *msg\_ptr, size\_t
    msg\_len, unsigned int *msg\_prio)} 
\fdesc{Effettua la ricezione di un messaggio da una coda.}
\fhead{mqueue.h} 
\fhead{time.h} 
\fdecl{ssize\_t mq\_timedreceive(mqd\_t mqdes, char *msg\_ptr, size\_t
    msg\_len,\\ 
\phantom{ssize\_t  mq\_timedreceive(}unsigned int *msg\_prio, const struct timespec
*abs\_timeout)} 
\fdesc{Riceve un messaggio da una coda entro un limite di tempo.}
}
{Entrambe le funzioni ritornano il numero di byte del messaggio in caso di
  successo e $-1$ per un errore, nel qual caso \var{errno} assumerà uno dei
  valori:
  \begin{errlist}
    \item[\errcode{EAGAIN}] si è aperta la coda con \const{O\_NONBLOCK}, e la
      coda è vuota.
    \item[\errcode{EINVAL}] si è specificato un valore nullo per
      \param{msg\_ptr}, o un valore non valido per \param{abs\_timeout}.
    \item[\errcode{EMSGSIZE}] la lunghezza del messaggio sulla coda eccede il
      valore \param{msg\_len} specificato per la ricezione.
    \item[\errcode{ETIMEDOUT}] la ricezione del messaggio non è stata
      effettuata entro il tempo stabilito.
  \end{errlist}
  ed inoltre \errval{EBADF} o \errval{EINTR} nel loro significato generico.  }
\end{funcproto}

La funzione estrae dalla coda \param{mqdes} il messaggio a priorità più alta,
o il più vecchio fra quelli della stessa priorità. Una volta ricevuto il
messaggio viene tolto dalla coda e la sua dimensione viene restituita come
valore di ritorno; si tenga presente che 0 è una dimensione valida e che la
condizione di errore è indicata soltanto da un valore di
$-1$.\footnote{Stevens in \cite{UNP2} fa notare che questo è uno dei casi in
  cui vale ciò che lo standard \textsl{non} dice, una dimensione nulla
  infatti, pur non essendo citata, non viene proibita.}

Se la dimensione specificata da \param{msg\_len} non è sufficiente a contenere
il messaggio, entrambe le funzioni, al contrario di quanto avveniva nelle code
di messaggi di SysV, ritornano un errore di \errcode{EMSGSIZE} senza estrarre
il messaggio.  È pertanto opportuno eseguire sempre una chiamata a
\func{mq\_getattr} prima di eseguire una ricezione, in modo da ottenere la
dimensione massima dei messaggi sulla coda, per poter essere in grado di
allocare dei buffer sufficientemente ampi per la lettura.

Se si specifica un puntatore per l'argomento \param{msg\_prio} il valore della
priorità del messaggio viene memorizzato all'indirizzo da esso indicato.
Qualora non interessi usare la priorità dei messaggi si può specificare
\var{NULL}, ed usare un valore nullo della priorità nelle chiamate a
\func{mq\_send}.

Si noti che con le code di messaggi POSIX non si ha la possibilità di
selezionare quale messaggio estrarre con delle condizioni sulla priorità, a
differenza di quanto avveniva con le code di messaggi di SysV che permettono
invece la selezione in base al valore del campo \var{mtype}. 

Qualora la coda sia vuota entrambe le funzioni si bloccano, a meno che non si
sia selezionata la modalità non bloccante; in tal caso entrambe ritornano
immediatamente con l'errore \errcode{EAGAIN}. Anche in questo caso la sola
differenza fra le due funzioni è che la seconda non attende indefinitamente e
passato il tempo massimo \param{abs\_timeout} ritorna comunque con un errore
di \errcode{ETIMEDOUT}.

Uno dei problemi sottolineati da Stevens in \cite{UNP2}, comuni ad entrambe le
tipologie di code messaggi, è che non è possibile per chi riceve identificare
chi è che ha inviato il messaggio, in particolare non è possibile sapere da
quale utente esso provenga. Infatti, in mancanza di un meccanismo interno al
kernel, anche se si possono inserire delle informazioni nel messaggio, queste
non possono essere credute, essendo completamente dipendenti da chi lo invia.
Vedremo però come, attraverso l'uso del meccanismo di notifica, sia possibile
superare in parte questo problema.

Una caratteristica specifica delle code di messaggi POSIX è la possibilità di
usufruire di un meccanismo di notifica asincrono; questo può essere attivato
usando la funzione \funcd{mq\_notify}, il cui prototipo è:

\begin{funcproto}{
\fhead{mqueue.h}
\fdecl{int mq\_notify(mqd\_t mqdes, const struct sigevent *notification)}

\fdesc{Attiva il meccanismo di notifica per una coda.}
}

{La funzione ritorna $0$ in caso di successo e $-1$ per un errore, nel qual
  caso \var{errno} assumerà uno dei valori:
  \begin{errlist}
    \item[\errcode{EBADF}] il descrittore non fa riferimento ad una coda di
      messaggi.
    \item[\errcode{EBUSY}] c'è già un processo registrato per la notifica.
    \item[\errcode{EINVAL}] si è richiesto un meccanismo di notifica invalido
      o specificato nella notifica con i segnali il valore di un segnale non
      esistente.
  \end{errlist}
  ed inoltre \errval{ENOMEM} nel suo significato generico.
}  
\end{funcproto}

Il meccanismo di notifica permette di segnalare in maniera asincrona ad un
processo la presenza di dati sulla coda indicata da \param{mqdes}, in modo da
evitare la necessità di bloccarsi nell'attesa. Per far questo un processo deve
registrarsi con la funzione \func{mq\_notify}, ed il meccanismo è disponibile
per un solo processo alla volta per ciascuna coda.

Il comportamento di \func{mq\_notify} dipende dai valori passati con
l'argomento \param{notification}, che è un puntatore ad una apposita struttura
\struct{sigevent}, (definita in fig.~\ref{fig:struct_sigevent}) introdotta
dallo standard POSIX.1b per gestire la notifica di eventi; per altri dettagli
su di essa si può rivedere quanto detto in sez.~\ref{sec:sig_timer_adv} a
proposito dell'uso della stessa struttura per la notifica delle scadenze dei
\textit{timer}.

Attraverso questa struttura si possono impostare le modalità con cui viene
effettuata la notifica nel campo \var{sigev\_notify}, che può assumere i
valori di tab.~\ref{tab:sigevent_sigev_notify}; fra questi la pagina di
manuale riporta soltanto i primi tre, ed inizialmente era possibile solo
\const{SIGEV\_SIGNAL}. Il metodo consigliato è quello di usare
\const{SIGEV\_SIGNAL} usando il campo \var{sigev\_signo} per indicare il quale
segnale deve essere inviato al processo. Inoltre il campo \var{sigev\_value} è
un puntatore ad una struttura \struct{sigval} (definita in
fig.~\ref{fig:sig_sigval}) che permette di restituire al gestore del segnale
un valore numerico o un indirizzo,\footnote{per il suo uso si riveda la
  trattazione fatta in sez.~\ref{sec:sig_real_time} a proposito dei segnali
  \textit{real-time}.} posto che questo sia installato nella forma estesa
vista in sez.~\ref{sec:sig_sigaction}.

La funzione registra il processo chiamante per la notifica se
\param{notification} punta ad una struttura \struct{sigevent} opportunamente
inizializzata, o cancella una precedente registrazione se è \val{NULL}. Dato
che un solo processo alla volta può essere registrato, la funzione fallisce
con \errcode{EBUSY} se c'è un altro processo già registrato. Questo significa
anche che se si registra una notifica con \const{SIGEV\_NONE} il processo non
la riceverà, ma impedirà anche che altri possano registrarsi per poterlo fare.
Si tenga presente inoltre che alla chiusura del descrittore associato alla
coda (e quindi anche all'uscita del processo) ogni eventuale registrazione di
notifica presente viene cancellata.

La notifica del segnale avviene all'arrivo di un messaggio in una coda vuota
(cioè solo se sulla coda non ci sono messaggi) e se non c'è nessun processo
bloccato in una chiamata a \func{mq\_receive}, in questo caso infatti il
processo bloccato ha la precedenza ed il messaggio gli viene immediatamente
inviato, mentre per il meccanismo di notifica tutto funziona come se la coda
fosse rimasta vuota.

Quando un messaggio arriva su una coda vuota al processo che si era registrato
viene inviato il segnale specificato da \code{notification->sigev\_signo}, e
la coda diventa disponibile per una ulteriore registrazione.  Questo comporta
che se si vuole mantenere il meccanismo di notifica occorre ripetere la
registrazione chiamando nuovamente \func{mq\_notify} all'interno del gestore
del segnale di notifica. A differenza della situazione simile che si aveva con
i segnali non affidabili (l'argomento è stato affrontato in
\ref{sec:sig_semantics}) questa caratteristica non configura una \textit{race
  condition} perché l'invio di un segnale avviene solo se la coda è vuota;
pertanto se si vuole evitare di correre il rischio di perdere eventuali
ulteriori segnali inviati nel lasso di tempo che occorre per ripetere la
richiesta di notifica basta avere cura di eseguire questa operazione prima di
estrarre i messaggi presenti dalla coda.

L'invio del segnale di notifica avvalora alcuni campi di informazione
restituiti al gestore attraverso la struttura \struct{siginfo\_t} (definita in
fig.~\ref{fig:sig_siginfo_t}). In particolare \var{si\_pid} viene impostato al
valore del \ids{PID} del processo che ha emesso il segnale, \var{si\_uid}
all'\textsl{user-ID} effettivo, \var{si\_code} a \const{SI\_MESGQ}, e
\var{si\_errno} a 0. Questo ci dice che, se si effettua la ricezione dei
messaggi usando esclusivamente il meccanismo di notifica, è possibile ottenere
le informazioni sul processo che ha inserito un messaggio usando un gestore
per il segnale in forma estesa, di nuovo si faccia riferimento a quanto detto
al proposito in sez.~\ref{sec:sig_sigaction} e sez.~\ref{sec:sig_real_time}.


\subsection{Memoria condivisa}
\label{sec:ipc_posix_shm}

La memoria condivisa è stato il primo degli oggetti di IPC POSIX inserito nel
kernel ufficiale; il supporto a questo tipo di oggetti è realizzato attraverso
il filesystem \texttt{tmpfs}, uno speciale filesystem che mantiene tutti i
suoi contenuti in memoria, che viene attivato abilitando l'opzione
\texttt{CONFIG\_TMPFS} in fase di compilazione del kernel.

Per potere utilizzare l'interfaccia POSIX per la memoria condivisa la
\acr{glibc} (le funzioni sono state introdotte con la versione 2.2) richiede
di compilare i programmi con l'opzione \code{-lrt}; inoltre è necessario che
in \file{/dev/shm} sia montato un filesystem \texttt{tmpfs}; questo di norma
viene fatto aggiungendo una riga del tipo di:
\begin{FileExample}[label=/etc/fstab]
tmpfs   /dev/shm        tmpfs   defaults        0      0
\end{FileExample}
ad \conffile{/etc/fstab}. In realtà si può montare un filesystem
\texttt{tmpfs} dove si vuole, per usarlo come RAM disk, con un comando del
tipo:
\begin{Example}
mount -t tmpfs -o size=128M,nr_inodes=10k,mode=700 tmpfs /mytmpfs
\end{Example}

Il filesystem riconosce, oltre quelle mostrate, le opzioni \texttt{uid} e
\texttt{gid} che identificano rispettivamente utente e gruppo cui assegnarne
la titolarità, e \texttt{nr\_blocks} che permette di specificarne la
dimensione in blocchi, cioè in multipli di \constd{PAGECACHE\_SIZE} che in
questo caso è l'unità di allocazione elementare.

La funzione che permette di aprire un segmento di memoria condivisa POSIX, ed
eventualmente di crearlo se non esiste ancora, è \funcd{shm\_open}; il suo
prototipo è:

\begin{funcproto}{
\fhead{sys/mman.h}
\fhead{sys/stat.h}
\fhead{fcntl.h}
\fdecl{int shm\_open(const char *name, int oflag, mode\_t mode)}

\fdesc{Apre un segmento di memoria condivisa.}
}

{La funzione ritorna un file descriptor in caso di successo e $-1$ per un
  errore, nel qual caso \var{errno} assumerà uno dei valori:
  \begin{errlist}
  \item[\errcode{EACCES}] non si hanno i permessi di aprire il segmento nella
    modalità scelta o si richiesto \const{O\_TRUNC} per un segmento su cui non
    si ha il permesso di scrittura.
  \item[\errcode{EINVAL}] si è utilizzato un nome non valido.
  \end{errlist}
  ed inoltre \errval{EEXIST}, \errval{EMFILE}, \errval{ENAMETOOLONG},
  \errval{ENFILE} e \errval{ENOENT} nello stesso significato che hanno per
  \func{open}.
}  
\end{funcproto}


La funzione apre un segmento di memoria condivisa identificato dal nome
\param{name}. Come già spiegato in sez.~\ref{sec:ipc_posix_generic} questo
nome può essere specificato in forma standard solo facendolo iniziare per
``\texttt{/}'' e senza ulteriori ``\texttt{/}''. Linux supporta comunque nomi
generici, che verranno interpretati prendendo come radice
\file{/dev/shm}.\footnote{occorre pertanto evitare di specificare qualcosa del
  tipo \file{/dev/shm/nome} all'interno di \param{name}, perché questo
  comporta, da parte delle funzioni di libreria, il tentativo di accedere a
  \file{/dev/shm/dev/shm/nome}.}

La funzione è del tutto analoga ad \func{open} ed analoghi sono i valori che
possono essere specificati per \param{oflag}, che deve essere specificato come
maschera binaria comprendente almeno uno dei due valori \const{O\_RDONLY} e
\const{O\_RDWR}; i valori possibili per i vari bit sono quelli visti in
sez.~\ref{sec:file_open_close} dei quali però \func{shm\_open} riconosce solo
i seguenti:
\begin{basedescript}{\desclabelwidth{2.0cm}\desclabelstyle{\nextlinelabel}}
\item[\const{O\_RDONLY}] Apre il file descriptor associato al segmento di
  memoria condivisa per l'accesso in sola lettura.
\item[\const{O\_RDWR}] Apre il file descriptor associato al segmento di
  memoria condivisa per l'accesso in lettura e scrittura.
\item[\const{O\_CREAT}] Necessario qualora si debba creare il segmento di
  memoria condivisa se esso non esiste; in questo caso viene usato il valore
  di \param{mode} per impostare i permessi, che devono essere compatibili con
  le modalità con cui si è aperto il file.
\item[\const{O\_EXCL}] Se usato insieme a \const{O\_CREAT} fa fallire la
  chiamata a \func{shm\_open} se il segmento esiste già, altrimenti esegue la
  creazione atomicamente.
\item[\const{O\_TRUNC}] Se il segmento di memoria condivisa esiste già, ne
  tronca le dimensioni a 0 byte.
\end{basedescript}

In caso di successo la funzione restituisce un file descriptor associato al
segmento di memoria condiviso con le stesse modalità di \func{open} viste in
sez.~\ref{sec:file_open_close}. Inoltre sul file descriptor viene sempre
impostato il flag \const{FD\_CLOEXEC}.  Chiamate effettuate da diversi
processi usando lo stesso nome restituiranno file descriptor associati allo
stesso segmento, così come, nel caso di file ordinari, essi sono associati
allo stesso inode. In questo modo è possibile effettuare una chiamata ad
\func{mmap} sul file descriptor restituito da \func{shm\_open} ed i processi
vedranno lo stesso segmento di memoria condivisa.

Quando il nome non esiste si può creare un nuovo segmento specificando
\const{O\_CREAT}; in tal caso il segmento avrà (così come i nuovi file)
lunghezza nulla. Il nuovo segmento verrà creato con i permessi indicati
da \param{mode} (di cui vengono usati solo i 9 bit meno significativi, non si
applicano pertanto i permessi speciali di sez.~\ref{sec:file_special_perm})
filtrati dal valore dell'\textit{umask} del processo. Come gruppo ed utente
proprietario del segmento saranno presi quelli facenti parte del gruppo
\textit{effective} del processo chiamante.

Dato che un segmento di lunghezza nulla è di scarsa utilità, una vola che lo
si è creato per impostarne la dimensione si dovrà poi usare \func{ftruncate}
(vedi sez.~\ref{sec:file_file_size}) prima di mapparlo in memoria con
\func{mmap}.  Si tenga presente che una volta chiamata \func{mmap} si può
chiudere il file descriptor ad esso associato (semplicemente con
\func{close}), senza che la mappatura ne risenta, e che questa può essere
rimossa usando \func{munmap}.

Come per i file, quando si vuole rimuovere completamente un segmento di
memoria condivisa occorre usare la funzione \funcd{shm\_unlink}, il cui
prototipo è:

\begin{funcproto}{
\fhead{sys/mman.h}
\fdecl{int shm\_unlink(const char *name)}

\fdesc{Rimuove un segmento di memoria condivisa.}
}

{La funzione ritorna $0$ in caso di successo e $-1$ per un errore, nel qual
  caso \var{errno} assumerà uno dei valori:
  \begin{errlist}
  \item[\errcode{EACCES}] non si è proprietari del segmento.
  \end{errlist}
  ed inoltre \errval{ENAMETOOLONG} e \errval{ENOENT}, nel loro significato
  generico.
}  
\end{funcproto}

La funzione è del tutto analoga ad \func{unlink}, e si limita a cancellare il
nome del segmento da \file{/dev/shm}, senza nessun effetto né sui file
descriptor precedentemente aperti con \func{shm\_open}, né sui segmenti già
mappati in memoria; questi verranno cancellati automaticamente dal sistema
solo con le rispettive chiamate a \func{close} e \func{munmap}.  Una volta
eseguita questa funzione però, qualora si richieda l'apertura di un segmento
con lo stesso nome, la chiamata a \func{shm\_open} fallirà, a meno di non aver
usato \const{O\_CREAT}, in quest'ultimo caso comunque si otterrà un file
descriptor che fa riferimento ad un segmento distinto da eventuali precedenti.

Dato che i segmenti di memoria condivisa sono trattati come file del
filesystem \texttt{tmpfs}, si possono usare su di essi, con lo stesso
significato che assumono sui file ordinari, anche funzioni come quelle delle
famiglie \func{fstat}, \func{fchown} e \func{fchmod}. Inoltre a partire dal
kernel 2.6.19 per i permessi sono supportate anche le ACL illustrate in
sez.~\ref{sec:file_ACL}.

Come esempio dell'uso delle funzioni attinenti ai segmenti di memoria
condivisa POSIX, vediamo come è possibile riscrivere una interfaccia
semplificata analoga a quella vista in fig.~\ref{fig:ipc_sysv_shm_func} per la
memoria condivisa in stile SysV. Il codice completo, di cui si sono riportate
le parti essenziali in fig.~\ref{fig:ipc_posix_shmmem}, è contenuto nel file
\file{SharedMem.c} dei sorgenti allegati.

\begin{figure}[!htb]
  \footnotesize \centering
  \begin{minipage}[c]{\codesamplewidth}
    \includecodesample{listati/MemShared.c}
  \end{minipage} 
  \normalsize 
  \caption{Il codice delle funzioni di gestione dei segmenti di memoria
    condivisa POSIX.}
  \label{fig:ipc_posix_shmmem}
\end{figure}

La prima funzione (\texttt{\small 1-24}) è \func{CreateShm} che, dato un nome
nell'argomento \var{name} crea un nuovo segmento di memoria condivisa,
accessibile in lettura e scrittura, e ne restituisce l'indirizzo. Anzitutto si
definiscono (\texttt{\small 8}) i flag per la successiva (\texttt{\small 9})
chiamata a \func{shm\_open}, che apre il segmento in lettura e scrittura
(creandolo se non esiste, ed uscendo in caso contrario) assegnandogli sul
filesystem i permessi specificati dall'argomento \var{perm}. 

In caso di errore (\texttt{\small 10-12}) si restituisce un puntatore nullo,
altrimenti si prosegue impostando (\texttt{\small 14}) la dimensione del
segmento con \func{ftruncate}. Di nuovo (\texttt{\small 15-16}) si esce
immediatamente restituendo un puntatore nullo in caso di errore. Poi si passa
(\texttt{\small 18}) a mappare in memoria il segmento con \func{mmap}
specificando dei diritti di accesso corrispondenti alla modalità di apertura.
Di nuovo si restituisce (\texttt{\small 19-21}) un puntatore nullo in caso di
errore, altrimenti si inizializza (\texttt{\small 22}) il contenuto del
segmento al valore specificato dall'argomento \var{fill} con \func{memset}, e
se ne restituisce (\texttt{\small 23}) l'indirizzo.

La seconda funzione (\texttt{\small 25-40}) è \func{FindShm} che trova un
segmento di memoria condiviso esistente, restituendone l'indirizzo. In questo
caso si apre (\texttt{\small 31}) il segmento con \func{shm\_open} richiedendo
che il segmento sia già esistente, in caso di errore (\texttt{\small 31-33})
si ritorna immediatamente un puntatore nullo.  Ottenuto il file descriptor del
segmento lo si mappa (\texttt{\small 35}) in memoria con \func{mmap},
restituendo (\texttt{\small 36-38}) un puntatore nullo in caso di errore, o
l'indirizzo (\texttt{\small 39}) dello stesso in caso di successo.

La terza funzione (\texttt{\small 40-45}) è \func{RemoveShm}, e serve a
cancellare un segmento di memoria condivisa. Dato che al contrario di quanto
avveniva con i segmenti del \textit{SysV-IPC} gli oggetti allocati nel kernel
vengono rilasciati automaticamente quando nessuna li usa più, tutto quello che
c'è da fare (\texttt{\small 44}) in questo caso è chiamare \func{shm\_unlink},
restituendo al chiamante il valore di ritorno.




\subsection{Semafori}
\label{sec:ipc_posix_sem}

Fino alla serie 2.4.x del kernel esisteva solo una implementazione parziale
dei semafori POSIX che li realizzava solo a livello di \textit{thread} e non
di processi,\footnote{questo significava che i semafori erano visibili solo
  all'interno dei \textit{thread} creati da un singolo processo, e non
  potevano essere usati come meccanismo di sincronizzazione fra processi
  diversi.} fornita attraverso la sezione delle estensioni \textit{real-time}
della \acr{glibc} (quelle che si accedono collegandosi alla libreria
\texttt{librt}). Esisteva inoltre una libreria che realizzava (parzialmente)
l'interfaccia POSIX usando le funzioni dei semafori di \textit{SysV-IPC}
(mantenendo così tutti i problemi sottolineati in
sez.~\ref{sec:ipc_sysv_sem}).

A partire dal kernel 2.5.7 è stato introdotto un meccanismo di
sincronizzazione completamente nuovo, basato sui cosiddetti \textit{futex} (la
sigla sta per \textit{fast user mode mutex}) con il quale è stato possibile
implementare una versione nativa dei semafori POSIX.  Grazie a questo con i
kernel della serie 2.6 e le nuove versioni della \acr{glibc} che usano questa
nuova infrastruttura per quella che viene chiamata \textit{New Posix Thread
  Library}, sono state implementate anche tutte le funzioni dell'interfaccia
dei semafori POSIX.

Anche in questo caso è necessario appoggiarsi alla libreria per le estensioni
\textit{real-time} \texttt{librt}, questo significa che se si vuole utilizzare
questa interfaccia, oltre ad utilizzare gli opportuni file di definizione,
occorrerà compilare i programmi con l'opzione \texttt{-lrt} o con
\texttt{-lpthread} se si usano questi ultimi. 

La funzione che permette di creare un nuovo semaforo POSIX, creando il
relativo file, o di accedere ad uno esistente, è \funcd{sem\_open}, questa
prevede due forme diverse a seconda che sia utilizzata per aprire un semaforo
esistente o per crearne uno nuovi, i relativi prototipi sono:

\begin{funcproto}{
\fhead{semaphore.h}
\fhead{sys/stat.h}
\fhead{fcntl.h}
\fdecl{sem\_t *sem\_open(const char *name, int oflag)}
\fdecl{sem\_t *sem\_open(const char *name, int oflag, mode\_t mode,
    unsigned int value)} 

\fdesc{Crea un semaforo o ne apre uno esistente.}
}
{La funzione ritorna l'indirizzo del semaforo in caso di successo e
  \constd{SEM\_FAILED} per un errore, nel qual caso \var{errno} assumerà uno
  dei valori:
  \begin{errlist}
    \item[\errcode{EACCES}] il semaforo esiste ma non si hanno permessi
      sufficienti per accedervi.
    \item[\errcode{EEXIST}] si sono specificati \const{O\_CREAT} e
      \const{O\_EXCL} ma il semaforo esiste.
    \item[\errcode{EINVAL}] il valore di \param{value} eccede
      \constd{SEM\_VALUE\_MAX} o il nome è solo ``\texttt{/}''.
    \item[\errcode{ENAMETOOLONG}] si è utilizzato un nome troppo lungo.
    \item[\errcode{ENOENT}] non si è usato \const{O\_CREAT} ed il nome
      specificato non esiste.
  \end{errlist}
  ed inoltre \errval{EMFILE}, \errval{ENFILE} ed \errval{ENOMEM} nel loro
  significato generico.

}
\end{funcproto}

L'argomento \param{name} definisce il nome del semaforo che si vuole
utilizzare, ed è quello che permette a processi diversi di accedere allo
stesso semaforo. Questo deve essere specificato nella stessa forma utilizzata
per i segmenti di memoria condivisa, con un nome che inizia con ``\texttt{/}''
e senza ulteriori ``\texttt{/}'', vale a dire nella forma
\texttt{/nome-semaforo}.

Con Linux i file associati ai semafori sono mantenuti nel filesystem virtuale
\texttt{/dev/shm}, e gli viene assegnato automaticamente un nome nella forma
\texttt{sem.nome-semaforo}, si ha cioè una corrispondenza per cui
\texttt{/nome-semaforo} viene rimappato, nella creazione tramite
\func{sem\_open}, su \texttt{/dev/shm/sem.nome-semaforo}. Per questo motivo la
dimensione massima per il nome di un semaforo, a differenza di quanto avviene
per i segmenti di memoria condivisa, è pari a \const{NAME\_MAX}$ - 4$.

L'argomento \param{oflag} è quello che controlla le modalità con cui opera la
funzione, ed è passato come maschera binaria; i bit corrispondono a quelli
utilizzati per l'analogo argomento di \func{open}, anche se dei possibili
valori visti in sez.~\ref{sec:file_open_close} sono utilizzati soltanto
\const{O\_CREAT} e \const{O\_EXCL}.

Se si usa \const{O\_CREAT} si richiede la creazione del semaforo qualora
questo non esista, ed in tal caso occorre utilizzare la seconda forma della
funzione, in cui si devono specificare sia un valore iniziale con l'argomento
\param{value},\footnote{e si noti come così diventa possibile, differenza di
  quanto avviene per i semafori del \textit{SysV-IPC}, effettuare in maniera
  atomica creazione ed inizializzazione di un semaforo usando una unica
  funzione.} che una maschera dei permessi con l'argomento
\param{mode}; se il semaforo esiste già questi saranno semplicemente
ignorati. Usando il flag \const{O\_EXCL} si richiede invece la verifica che il
semaforo non esista, ed usandolo insieme ad \const{O\_CREAT} la funzione
fallisce qualora un semaforo con lo stesso nome sia già presente.

Si tenga presente che, come accennato in sez.~\ref{sec:ipc_posix_generic}, i
semafori usano la semantica standard dei file per quanto riguarda i controlli
di accesso, questo significa che un nuovo semaforo viene sempre creato con
l'\ids{UID} ed il \ids{GID} effettivo del processo chiamante, e che i permessi
indicati con \param{mode} vengono filtrati dal valore della \textit{umask} del
processo.  Inoltre per poter aprire un semaforo è necessario avere su di esso
sia il permesso di lettura che quello di scrittura.

La funzione restituisce in caso di successo un puntatore all'indirizzo del
semaforo con un valore di tipo \ctyp{sem\_t *}, è questo valore che dovrà
essere passato alle altre funzioni per operare sul semaforo stesso, e non sarà
più necessario fare riferimento al nome, che potrebbe anche essere rimosso con
\func{sem\_unlink}.

Una volta che si sia ottenuto l'indirizzo di un semaforo, sarà possibile
utilizzarlo; se si ricorda quanto detto all'inizio di
sez.~\ref{sec:ipc_sysv_sem}, dove si sono introdotti i concetti generali
relativi ai semafori, le operazioni principali sono due, quella che richiede
l'uso di una risorsa bloccando il semaforo e quella che rilascia la risorsa
liberando il semaforo. La prima operazione è effettuata dalla funzione
\funcd{sem\_wait}, il cui prototipo è:

\begin{funcproto}{
\fhead{semaphore.h}
\fdecl{int sem\_wait(sem\_t *sem)}

\fdesc{Blocca un semaforo.}
}

{La funzione ritorna $0$ in caso di successo e $-1$ per un errore, nel qual
  caso \var{errno} assumerà uno dei valori:
  \begin{errlist}
    \item[\errcode{EINTR}] la funzione è stata interrotta da un segnale.
    \item[\errcode{EINVAL}] il semaforo \param{sem} non esiste.
  \end{errlist}
}  
\end{funcproto}

La funzione cerca di decrementare il valore del semaforo indicato dal
puntatore \param{sem}, se questo ha un valore positivo, cosa che significa che
la risorsa è disponibile, la funzione ha successo, il valore del semaforo
viene diminuito di 1 ed essa ritorna immediatamente consentendo la
prosecuzione del processo.

Se invece il valore è nullo la funzione si blocca (fermando l'esecuzione del
processo) fintanto che il valore del semaforo non ritorna positivo (cosa che a
questo punto può avvenire solo per opera di altro processo che rilascia il
semaforo con una chiamata a \func{sem\_post}) così che poi essa possa
decrementarlo con successo e proseguire.

Si tenga presente che la funzione può sempre essere interrotta da un segnale,
nel qual caso si avrà un errore di \errval{EINTR}; inoltre questo avverrà
comunque, anche qualora si fosse richiesta la gestione con la semantica BSD,
installando il gestore del suddetto segnale con l'opzione \const{SA\_RESTART}
(vedi sez.~\ref{sec:sig_sigaction}) per riavviare le \textit{system call}
interrotte.

Della funzione \func{sem\_wait} esistono due varianti che consentono di
gestire diversamente le modalità di attesa in caso di risorsa occupata, la
prima di queste è \funcd{sem\_trywait}, che serve ad effettuare un tentativo
di acquisizione senza bloccarsi; il suo prototipo è:

\begin{funcproto}{
\fhead{semaphore.h} 
\fdecl{int sem\_trywait(sem\_t *sem)}

\fdesc{Tenta di bloccare un semaforo.}
}
{La funzione ritorna $0$ in caso di successo e $-1$ per un errore, nel qual
  caso \var{errno} assumerà uno dei valori:
  \begin{errlist}
    \item[\errcode{EAGAIN}] il semaforo non può essere acquisito senza
      bloccarsi. 
    \item[\errcode{EINVAL}] l'argomento \param{sem} non indica un semaforo
      valido.
  \end{errlist}
}
\end{funcproto}

La funzione è identica a \func{sem\_wait} ed se la risorsa è libera ha lo
stesso effetto, vale a dire che in caso di semaforo diverso da zero la
funzione lo decrementa e ritorna immediatamente; la differenza è che nel caso
in cui il semaforo è occupato essa non si blocca e di nuovo ritorna
immediatamente, restituendo però un errore di \errval{EAGAIN}, così che il
programma possa proseguire.

La seconda variante di \func{sem\_wait} è una estensione specifica che può
essere utilizzata soltanto se viene definita la macro \macro{\_XOPEN\_SOURCE}
ad un valore di almeno 600 o la macro \macro{\_POSIX\_C\_SOURCE} ad un valore
uguale o maggiore di \texttt{200112L} prima di includere
\headfiled{semaphore.h}, la funzione è \funcd{sem\_timedwait}, ed il suo
prototipo è:

\begin{funcproto}{
\fhead{semaphore.h}
\fdecl{int sem\_timedwait(sem\_t *sem, const struct timespec
    *abs\_timeout)}

\fdesc{Blocca un semaforo.}
}

{La funzione ritorna $0$ in caso di successo e $-1$ per un errore, nel qual
  caso \var{errno} assumerà uno dei valori:
  \begin{errlist}
    \item[\errcode{EINTR}] la funzione è stata interrotta da un segnale.
    \item[\errcode{EINVAL}] l'argomento \param{sem} non indica un semaforo
      valido.
    \item[\errcode{ETIMEDOUT}] è scaduto il tempo massimo di attesa. 
  \end{errlist}
}  
\end{funcproto}

Anche in questo caso il comportamento della funzione è identico a quello di
\func{sem\_wait}, ma è possibile impostare un tempo limite per l'attesa
tramite la struttura \struct{timespec} (vedi
fig.~\ref{fig:sys_timespec_struct}) puntata
dall'argomento \param{abs\_timeout}, indicato in secondi e nanosecondi a
partire dalla cosiddetta \textit{Epoch} (00:00:00, 1 January 1970
UTC). Scaduto il limite la funzione ritorna anche se non è possibile acquisire
il semaforo fallendo con un errore di \errval{ETIMEDOUT}.

La seconda funzione principale utilizzata per l'uso dei semafori è quella che
viene usata per rilasciare un semaforo occupato o, in generale, per aumentare
di una unità il valore dello stesso anche qualora non fosse occupato (si
ricordi che in generale un semaforo viene usato come indicatore di un numero
di risorse disponibili). Detta funzione è \funcd{sem\_post} ed il suo
prototipo è:

\begin{funcproto}{
\fhead{semaphore.h}
\fdecl{int sem\_post(sem\_t *sem)}

\fdesc{Rilascia un semaforo.}
}

{La funzione ritorna $0$ in caso di successo e $-1$ per un errore, nel qual
  caso \var{errno} assumerà uno dei valori:
  \begin{errlist}
    \item[\errcode{EINVAL}] l'argomento \param{sem} non indica un semaforo
      valido.
    \item[\errcode{EOVERFLOW}] si superato il massimo valore di un semaforo.
  \end{errlist}
}  
\end{funcproto}

La funzione incrementa di uno il valore corrente del semaforo indicato
dall'argomento \param{sem}, se questo era nullo la relativa risorsa risulterà
sbloccata, cosicché un altro processo (o \textit{thread}) eventualmente
bloccato in una \func{sem\_wait} sul semaforo possa essere svegliato e rimesso
in esecuzione.  Si tenga presente che la funzione è sicura per l'uso
all'interno di un gestore di segnali (si ricordi quanto detto in
sez.~\ref{sec:sig_signal_handler}).

Se invece di operare su un semaforo se ne volesse semplicemente leggere il
valore, si potrà usare la funzione \funcd{sem\_getvalue}, il cui prototipo è:

\begin{funcproto}{
\fhead{semaphore.h}
\fdecl{int sem\_getvalue(sem\_t *sem, int *sval)}

\fdesc{Richiede il valore di un semaforo.}
}
{La funzione ritorna $0$ in caso di successo e $-1$ per un errore, nel qual
  caso \var{errno} assumerà uno dei valori:
  \begin{errlist}
    \item[\errcode{EINVAL}] l'argomento \param{sem} non indica un semaforo
      valido.
  \end{errlist}
}  
\end{funcproto}

La funzione legge il valore del semaforo indicato dall'argomento \param{sem} e
lo restituisce nella variabile intera puntata dall'argomento
\param{sval}. Qualora ci siano uno o più processi bloccati in attesa sul
semaforo lo standard prevede che la funzione possa restituire un valore nullo
oppure il numero di processi bloccati in una \func{sem\_wait} sul suddetto
semaforo; nel caso di Linux vale la prima opzione.

Questa funzione può essere utilizzata per avere un suggerimento sullo stato di
un semaforo, ovviamente non si può prendere il risultato riportato in
\param{sval} che come indicazione, il valore del semaforo infatti potrebbe
essere già stato modificato al ritorno della funzione.

% TODO verificare comportamento sem_getvalue

Una volta che non ci sia più la necessità di operare su un semaforo se ne può
terminare l'uso con la funzione \funcd{sem\_close}, il cui prototipo è:

\begin{funcproto}{
\fhead{semaphore.h}
\fdecl{int sem\_close(sem\_t *sem)}

\fdesc{Chiude un semaforo.}
}

{La funzione ritorna $0$ in caso di successo e $-1$ per un errore, nel qual
  caso \var{errno} assumerà uno dei valori:
  \begin{errlist}
    \item[\errcode{EINVAL}] l'argomento \param{sem} non indica un semaforo
      valido.
  \end{errlist}
}  
\end{funcproto}

La funzione chiude il semaforo indicato dall'argomento \param{sem}, che non
potrà più essere utilizzato nelle altre funzioni. La chiusura comporta anche
che tutte le risorse che il sistema poteva avere assegnato al processo
nell'uso del semaforo vengono immediatamente rilasciate. Questo significa che
un eventuale altro processo bloccato sul semaforo a causa della acquisizione
dello stesso da parte del processo che chiama \func{sem\_close} potrà essere
immediatamente riavviato.

Si tenga presente poi che come avviene per i file, all'uscita di un processo
anche tutti i semafori che questo aveva aperto vengono automaticamente chiusi.
Questo comportamento risolve il problema che si aveva con i semafori del
\textit{SysV IPC} (di cui si è parlato in sez.~\ref{sec:ipc_sysv_sem}) per i
quali le risorse possono restare bloccate. Si tenga infine presente che, a
differenza di quanto avviene per i file, in caso di una chiamata ad
\func{execve} tutti i semafori vengono chiusi automaticamente.

Come per i semafori del \textit{SysV-IPC} anche quelli POSIX hanno una
persistenza di sistema; questo significa che una volta che si è creato un
semaforo con \func{sem\_open} questo continuerà ad esistere fintanto che il
kernel resta attivo (vale a dire fino ad un successivo riavvio) a meno che non
lo si cancelli esplicitamente. Per far questo si può utilizzare la funzione
\funcd{sem\_unlink}, il cui prototipo è:

\begin{funcproto}{
\fhead{semaphore.h}
\fdecl{int sem\_unlink(const char *name)}

\fdesc{Rimuove un semaforo.}
}

{La funzione ritorna $0$ in caso di successo e $-1$ per un errore, nel qual
  caso \var{errno} assumerà uno dei valori:
  \begin{errlist}
    \item[\errcode{EACCES}] non si hanno i permessi necessari a cancellare il
      semaforo.
    \item[\errcode{ENAMETOOLONG}] il nome indicato è troppo lungo.
    \item[\errcode{ENOENT}] il semaforo \param{name} non esiste.
  \end{errlist}
}  
\end{funcproto}

La funzione rimuove il semaforo indicato dall'argomento \param{name}, che
prende un valore identico a quello usato per creare il semaforo stesso con
\func{sem\_open}. Il semaforo viene rimosso dal filesystem immediatamente; ma
il semaforo viene effettivamente cancellato dal sistema soltanto quando tutti
i processi che lo avevano aperto lo chiudono. Si segue cioè la stessa
semantica usata con \func{unlink} per i file, trattata in dettaglio in
sez.~\ref{sec:link_symlink_rename}.

Una delle caratteristiche peculiari dei semafori POSIX è che questi possono
anche essere utilizzati anche in forma anonima, senza necessità di fare
ricorso ad un nome sul filesystem o ad altri indicativi.  In questo caso si
dovrà porre la variabile che contiene l'indirizzo del semaforo in un tratto di
memoria che sia accessibile a tutti i processi in gioco.  La funzione che
consente di inizializzare un semaforo anonimo è \funcd{sem\_init}, il cui
prototipo è:

\begin{funcproto}{
\fhead{semaphore.h}
\fdecl{int sem\_init(sem\_t *sem, int pshared, unsigned int value)}
\fdesc{Inizializza un semaforo anonimo.}
}

{La funzione ritorna $0$ in caso di successo e $-1$ per un errore, nel qual
  caso \var{errno} assumerà uno dei valori:
  \begin{errlist}
    \item[\errcode{EINVAL}] il valore di \param{value} eccede
      \const{SEM\_VALUE\_MAX}.
    \item[\errcode{ENOSYS}] il valore di \param{pshared} non è nullo ed il
      sistema non supporta i semafori per i processi.
  \end{errlist}
}  
\end{funcproto}

La funzione inizializza un semaforo all'indirizzo puntato dall'argomento
\param{sem}, e come per \func{sem\_open} consente di impostare un valore
iniziale con \param{value}. L'argomento \param{pshared} serve ad indicare se
il semaforo deve essere utilizzato dai \textit{thread} di uno stesso processo
(con un valore nullo) o condiviso fra processi diversi (con un valore non
nullo).

Qualora il semaforo debba essere condiviso dai \textit{thread} di uno stesso
processo (nel qual caso si parla di \textit{thread-shared semaphore}),
occorrerà che \param{sem} sia l'indirizzo di una variabile visibile da tutti i
\textit{thread}, si dovrà usare cioè una variabile globale o una variabile
allocata dinamicamente nello \textit{heap}.

Qualora il semaforo debba essere condiviso fra più processi (nel qual caso si
parla di \textit{process-shared semaphore}) la sola scelta possibile per
renderlo visibile a tutti è di porlo in un tratto di memoria condivisa. Questo
potrà essere ottenuto direttamente sia con \func{shmget} (vedi
sez.~\ref{sec:ipc_sysv_shm}) che con \func{shm\_open} (vedi
sez.~\ref{sec:ipc_posix_shm}), oppure, nel caso che tutti i processi in gioco
abbiano un genitore comune, con una mappatura anonima con \func{mmap} (vedi
sez.~\ref{sec:file_memory_map}) a cui essi poi potranno accedere (si ricordi
che i tratti di memoria condivisa vengono mantenuti nei processi figli
attraverso la funzione \func{fork}).

Una volta inizializzato il semaforo anonimo con \func{sem\_init} lo si potrà
utilizzare nello stesso modo dei semafori normali con \func{sem\_wait} e
\func{sem\_post}. Si tenga presente però che inizializzare due volte lo stesso
semaforo può dar luogo ad un comportamento indefinito. 

Qualora non si intenda più utilizzare un semaforo anonimo questo può essere
eliminato dal sistema; per far questo di deve utilizzare una apposita
funzione, \funcd{sem\_destroy}, il cui prototipo è:

\begin{funcproto}{
\fhead{semaphore.h}
\fdecl{int sem\_destroy(sem\_t *sem)}
\fdesc{Elimina un semaforo anonimo.}
}
{La funzione ritorna $0$ in caso di successo e $-1$ per un errore, nel qual
  caso \var{errno} assumerà uno dei valori:
  \begin{errlist}
    \item[\errcode{EINVAL}] l'argomento \param{sem} non indica un semaforo
      valido.
  \end{errlist}
}  
\end{funcproto}

La funzione prende come unico argomento l'indirizzo di un semaforo che deve
essere stato inizializzato con \func{sem\_init}; non deve quindi essere
applicata a semafori creati con \func{sem\_open}. Inoltre si deve essere
sicuri che il semaforo sia effettivamente inutilizzato, la distruzione di un
semaforo su cui sono presenti processi (o \textit{thread}) in attesa (cioè
bloccati in una \func{sem\_wait}) provoca un comportamento indefinito.

Si tenga presente infine che utilizzare un semaforo che è stato distrutto con
\func{sem\_destroy} di nuovo può dare esito a comportamenti indefiniti.  Nel
caso ci si trovi in una tale evenienza occorre reinizializzare il semaforo una
seconda volta con \func{sem\_init}.

Come esempio di uso sia della memoria condivisa che dei semafori POSIX si sono
scritti due semplici programmi con i quali è possibile rispettivamente
monitorare il contenuto di un segmento di memoria condivisa e modificarne il
contenuto. 

\begin{figure}[!htbp]
  \footnotesize \centering
  \begin{minipage}[c]{\codesamplewidth}
    \includecodesample{listati/message_getter.c}
  \end{minipage} 
  \normalsize 
  \caption{Sezione principale del codice del programma
    \file{message\_getter.c}.}
  \label{fig:ipc_posix_sem_shm_message_server}
\end{figure}

Il corpo principale del primo dei due, il cui codice completo è nel file
\file{message\_getter.c} dei sorgenti allegati, è riportato in
fig.~\ref{fig:ipc_posix_sem_shm_message_server}; si è tralasciata la parte che
tratta la gestione delle opzioni a riga di comando (che consentono di
impostare un nome diverso per il semaforo e il segmento di memoria condivisa)
ed il controllo che al programma venga fornito almeno un argomento, contenente
la stringa iniziale da inserire nel segmento di memoria condivisa.

Lo scopo del programma è quello di creare un segmento di memoria condivisa su
cui registrare una stringa, e tenerlo sotto osservazione stampando la stessa
una volta al secondo. Si utilizzerà un semaforo per proteggere l'accesso in
lettura alla stringa, in modo che questa non possa essere modificata
dall'altro programma prima di averla finita di stampare.

La parte iniziale del programma contiene le definizioni (\texttt{\small 1-8})
del gestore del segnale usato per liberare le risorse utilizzate, delle
variabili globali contenenti i nomi di default del segmento di memoria
condivisa e del semaforo (il default scelto è \texttt{messages}), e delle
altre variabili utilizzate dal programma.

Come prima istruzione (\texttt{\small 10}) si è provveduto ad installare un
gestore di segnale che consentirà di effettuare le operazioni di pulizia
(usando la funzione \func{Signal} illustrata in
fig.~\ref{fig:sig_Signal_code}), dopo di che (\texttt{\small 12-16}) si è
creato il segmento di memoria condivisa con la funzione \func{CreateShm} che
abbiamo appena trattato in sez.~\ref{sec:ipc_posix_shm}, uscendo con un
messaggio in caso di errore. 

Si tenga presente che la funzione \func{CreateShm} richiede che il segmento
non sia già presente e fallirà qualora un'altra istanza, o un altro programma
abbia già allocato un segmento con quello stesso nome. Per semplicità di
gestione si è usata una dimensione fissa pari a 256 byte, definita tramite la
costante \texttt{MSGMAXSIZE}.

Il passo successivo (\texttt{\small 17-21}) è quello della creazione del
semaforo che regola l'accesso al segmento di memoria condivisa con
\func{sem\_open}; anche in questo caso si gestisce l'uscita con stampa di un
messaggio in caso di errore. Anche per il semaforo, avendo specificato la
combinazione di flag \code{O\_CREAT|O\_EXCL} come secondo argomento, si esce
qualora fosse già esistente; altrimenti esso verrà creato con gli opportuni
permessi specificati dal terzo argomento, (indicante lettura e scrittura in
notazione ottale). Infine il semaforo verrà inizializzato ad un valore nullo
(il quarto argomento), corrispondete allo stato in cui risulta bloccato.

A questo punto (\texttt{\small 22}) si potrà inizializzare il messaggio posto
nel segmento di memoria condivisa usando la stringa passata come argomento al
programma. Essendo il semaforo stato creato già bloccato non ci si dovrà
preoccupare di eventuali \textit{race condition} qualora il programma di
modifica del messaggio venisse lanciato proprio in questo momento.  Una volta
inizializzato il messaggio occorrerà però rilasciare il semaforo
(\texttt{\small 24-27}) per consentirne l'uso; in tutte queste operazioni si
provvederà ad uscire dal programma con un opportuno messaggio in caso di
errore.

Una volta completate le inizializzazioni il ciclo principale del programma
(\texttt{\small 29-47}) viene ripetuto indefinitamente (\texttt{\small 29})
per stampare sia il contenuto del messaggio che una serie di informazioni di
controllo. Il primo passo (\texttt{\small 30-34}) è quello di acquisire (con
\func{sem\_getvalue}, con uscita in caso di errore) e stampare il valore del
semaforo ad inizio del ciclo; seguito (\texttt{\small 35-36}) dal tempo
corrente.

\begin{figure}[!htb]
  \footnotesize \centering
  \begin{minipage}[c]{\codesamplewidth}
    \includecodesample{listati/HandSigInt.c}
  \end{minipage} 
  \normalsize 
  \caption{Codice del gestore di segnale del programma
    \file{message\_getter.c}.}
  \label{fig:ipc_posix_sem_shm_message_server_handler}
\end{figure}

Prima della stampa del messaggio invece si deve acquisire il semaforo
(\texttt{\small 30-33}) per evitare accessi concorrenti alla stringa da parte
del programma di modifica. Una volta eseguita la stampa (\texttt{\small 41})
il semaforo dovrà essere rilasciato (\texttt{\small 42-45}). Il passo finale
(\texttt{\small 46}) è attendere per un secondo prima di eseguire da capo il
ciclo. 

Per uscire in maniera corretta dal programma sarà necessario fermarlo con una
interruzione da tastiera (\texttt{C-c}), che corrisponde all'invio del segnale
\signal{SIGINT}, per il quale si è installato (\texttt{\small 10}) una
opportuna funzione di gestione, riportata in
fig.~\ref{fig:ipc_posix_sem_shm_message_server_handler}. La funzione è molto
semplice e richiama le funzioni di rimozione sia per il segmento di memoria
condivisa che per il semaforo, garantendo così che possa essere riaperto
ex-novo senza errori in un futuro riutilizzo del comando.

\begin{figure}[!htb]
  \footnotesize \centering
  \begin{minipage}[c]{\codesamplewidth}
    \includecodesample{listati/message_setter.c}
  \end{minipage} 
  \normalsize 
  \caption{Sezione principale del codice del programma
    \file{message\_setter.c}.}
  \label{fig:ipc_posix_sem_shm_message_setter}
\end{figure}

Il secondo programma di esempio è \file{message\_setter.c}, di cui si è
riportato il corpo principale in
fig.~\ref{fig:ipc_posix_sem_shm_message_setter},\footnote{al solito il codice
  completo è nel file dei sorgenti allegati.} dove si è tralasciata, non
essendo significativa per quanto si sta trattando, la parte relativa alla
gestione delle opzioni a riga di comando e degli argomenti, che sono identici
a quelli usati da \file{message\_getter}, con l'unica aggiunta di un'opzione
``\texttt{-t}'' che consente di indicare un tempo di attesa (in secondi) in
cui il programma si ferma tenendo bloccato il semaforo.

Una volta completata la gestione delle opzioni e degli argomenti (ne deve
essere presente uno solo, contenente la nuova stringa da usare come
messaggio), il programma procede (\texttt{\small 10-14}) con l'acquisizione
del segmento di memoria condivisa usando la funzione \func{FindShm} (trattata
in sez.~\ref{sec:ipc_posix_shm}) che stavolta deve già esistere.  Il passo
successivo (\texttt{\small 16-19}) è quello di aprire il semaforo, e a
differenza di \file{message\_getter}, in questo caso si richiede a
\func{sem\_open} che questo esista, passando uno zero come secondo ed unico
argomento.

Una volta completate con successo le precedenti inizializzazioni, il passo
seguente (\texttt{\small 21-24}) è quello di acquisire il semaforo, dopo di
che sarà possibile eseguire la sostituzione del messaggio (\texttt{\small 25})
senza incorrere in possibili \textit{race condition} con la stampa dello
stesso da parte di \file{message\_getter}.

Una volta effettuata la modifica viene stampato (\texttt{\small 26}) il tempo
di attesa impostato con l'opzione ``\texttt{-t}'' dopo di che (\texttt{\small
  27}) viene eseguita la stessa, senza rilasciare il semaforo che resterà
quindi bloccato (causando a questo punto una interruzione delle stampe
eseguite da \file{message\_getter}). Terminato il tempo di attesa si rilascerà
(\texttt{\small 29-32}) il semaforo per poi uscire.

Per verificare il funzionamento dei programmi occorrerà lanciare prima
\file{message\_getter} (lanciare per primo \file{message\_setter} darebbe
luogo ad un errore, non essendo stati creati il semaforo ed il segmento di
memoria condivisa) che inizierà a stampare una volta al secondo il contenuto
del messaggio ed i suoi dati, con qualcosa del tipo:
\begin{Console}
piccardi@hain:~/gapil/sources$  \textbf{./message_getter messaggio}
sem=1, Fri Dec 31 14:12:41 2010
message: messaggio
sem=1, Fri Dec 31 14:12:42 2010
message: messaggio
...
\end{Console}
%$
proseguendo indefinitamente fintanto che non si prema \texttt{C-c} per farlo
uscire. Si noti come il valore del semaforo risulti sempre pari ad 1 (in
quanto al momento esso sarà sempre libero). 

A questo punto si potrà lanciare \file{message\_setter} per cambiare il
messaggio, nel nostro caso per rendere evidente il funzionamento del blocco
richiederemo anche una attesa di 3 secondi, ed otterremo qualcosa del tipo:
\begin{Console}
piccardi@hain:~/gapil/sources$ \textbf{./message_setter -t 3 ciao}
Sleeping for 3 seconds
\end{Console}
%$
dove il programma si fermerà per 3 secondi prima di rilasciare il semaforo e
terminare. L'effetto di tutto ciò si potrà vedere nell'\textit{output} di
\file{message\_getter}, che verrà interrotto per questo stesso tempo, prima di
ricominciare con il nuovo testo:
\begin{Console}
...
sem=1, Fri Dec 31 14:16:27 2010
message: messaggio
sem=1, Fri Dec 31 14:16:28 2010
message: messaggio
sem=0, Fri Dec 31 14:16:29 2010
message: ciao
sem=1, Fri Dec 31 14:16:32 2010
message: ciao
sem=1, Fri Dec 31 14:16:33 2010
message: ciao
...
\end{Console}
%$

E si noterà come nel momento in cui si lancia \file{message\_setter} le stampe
di \file{message\_getter} si bloccheranno, come corretto, dopo aver registrato
un valore nullo per il semaforo.  Il programma infatti resterà bloccato nella
\func{sem\_wait} (quella di riga (\texttt{\small 37}) in
fig.~\ref{fig:ipc_posix_sem_shm_message_server}) fino alla scadenza
dell'attesa di \file{message\_setter} (con l'esecuzione della \func{sem\_post}
della riga (\texttt{\small 29}) di
fig.~\ref{fig:ipc_posix_sem_shm_message_setter}), e riprenderanno con il nuovo
testo alla terminazione di quest'ultimo.


% LocalWords:  like fifo System POSIX RPC Calls Common Object Request Brocker
% LocalWords:  Architecture descriptor kernel unistd int filedes errno EMFILE
% LocalWords:  ENFILE EFAULT BUF sez fig fork Stevens siblings EOF read SIGPIPE
% LocalWords:  EPIPE shell CGI Gateway Interface HTML JPEG URL mime type gs dup
% LocalWords:  barcode PostScript race condition stream BarCodePage WriteMess
% LocalWords:  size PS switch wait popen pclose stdio const char command NULL
% LocalWords:  EINVAL cap fully buffered Ghostscript l'Encapsulated epstopsf of
% LocalWords:  PDF EPS lseek ESPIPE PPM Portable PixMap format pnmcrop PNG pnm
% LocalWords:  pnmmargin png BarCode inode filesystem l'inode mknod mkfifo RDWR
% LocalWords:  ENXIO deadlock client reinviate fortunes fortunefilename daemon
% LocalWords:  FortuneServer FortuneParse FortuneClient pid libgapil  LD librt
% LocalWords:  PATH linker pathname ps tmp killall fortuned crash socket domain
% LocalWords:  socketpair BSD sys protocol sv EAFNOSUPPORT EPROTONOSUPPORT AF
% LocalWords:  EOPNOTSUPP SOCK Process Comunication ipc perm key exec pipefd SZ
% LocalWords:  header ftok proj stat libc SunOS glibc XPG dell'inode number uid
% LocalWords:  cuid cgid gid tab MSG shift group umask seq MSGMNI SEMMNI SHMMNI
% LocalWords:  shmmni msgmni sem sysctl IPCMNI IPCTestId msgget EACCES EEXIST
% LocalWords:  CREAT EXCL EIDRM ENOENT ENOSPC ENOMEM novo proc MSGMAX msgmax ds
% LocalWords:  MSGMNB msgmnb linked list msqid msgid linux msg qnum lspid lrpid
% LocalWords:  rtime ctime qbytes first last cbytes msgctl semctl shmctl ioctl
% LocalWords:  cmd struct buf EPERM RMID msgsnd msgbuf msgp msgsz msgflg EAGAIN
% LocalWords:  NOWAIT EINTR mtype mtext long message sizeof LENGTH ts sleep BIG
% LocalWords:  msgrcv ssize msgtyp NOERROR EXCEPT ENOMSG multiplexing select ls
% LocalWords:  poll polling queue MQFortuneServer write init HandSIGTERM  l'IPC
% LocalWords:  MQFortuneClient mqfortuned mutex risorse' inter semaphore semget
% LocalWords:  nsems SEMMNS SEMMSL semid otime semval sempid semncnt semzcnt nr
% LocalWords:  SEMVMX SEMOPM semop SEMMNU SEMUME SEMAEM semnum union semun arg
% LocalWords:  ERANGE SETALL SETVAL GETALL array GETNCNT GETPID GETVAL GETZCNT
% LocalWords:  sembuf sops unsigned nsops UNDO flg nsop num undo pending semadj
% LocalWords:  sleeper scheduler running next semundo MutexCreate semunion lock
% LocalWords:  MutexFind wrapper MutexRead MutexLock MutexUnlock unlock locking
% LocalWords:  MutexRemove shmget SHMALL SHMMAX SHMMIN shmid shm segsz atime FD
% LocalWords:  dtime lpid cpid nattac shmall shmmax SHMLBA SHMSEG EOVERFLOW brk
% LocalWords:  memory shmat shmdt void shmaddr shmflg SVID RND RDONLY rounded
% LocalWords:  SIGSEGV nattch exit SharedMem ShmCreate memset fill ShmFind home
% LocalWords:  ShmRemove DirMonitor DirProp chdir GaPiL shmptr ipcs NFS SETPIPE
% LocalWords:  ComputeValues ReadMonitor touch SIGTERM dirmonitor unlink fcntl
% LocalWords:  LockFile UnlockFile CreateMutex FindMutex LockMutex SETLKW GETLK
% LocalWords:  UnlockMutex RemoveMutex ReadMutex UNLCK WRLCK RDLCK mapping MAP
% LocalWords:  SHARED ANONYMOUS thread patch names strace system call userid Di
% LocalWords:  groupid Michal Wronski Krzysztof Benedyczak wrona posix mqueue
% LocalWords:  lmqueue gcc mount mqd name oflag attr maxmsg msgsize receive ptr
% LocalWords:  send WRONLY NONBLOCK close mqdes EBADF getattr setattr mqstat to
% LocalWords:  omqstat curmsgs flags timedsend len prio timespec abs EMSGSIZE
% LocalWords:  ETIMEDOUT timedreceive getaddr notify sigevent notification l'I
% LocalWords:  EBUSY sigev SIGNAL signo value sigval siginfo all'userid MESGQ
% LocalWords:  Konstantin Knizhnik futex tmpfs ramfs cache shared swap CONFIG
% LocalWords:  lrt blocks PAGECACHE TRUNC CLOEXEC mmap ftruncate munmap FindShm
% LocalWords:  CreateShm RemoveShm LIBRARY Library libmqueue FAILED has fclose
% LocalWords:  ENAMETOOLONG qualchenome RESTART trywait XOPEN SOURCE timedwait
% LocalWords:  process getvalue sval execve pshared ENOSYS heap PAGE destroy it
% LocalWords:  xffffffff Arrays owner perms Queues used bytes messages device
% LocalWords:  Cannot find such Segments getter Signal MSGMAXSIZE been stable
% LocalWords:  for now it's break Berlin sources Let's an accidental feature fs
% LocalWords:  Larry Wall Escape the Hell William ipctestid Identifier segment
% LocalWords:  violation dell'I SIGINT setter Fri Dec Sleeping seconds ECHILD
% LocalWords:  SysV capability short RESOURCE INFO UNDEFINED EFBIG semtimedop
% LocalWords:  scan HUGETLB huge page NORESERVE copy RLIMIT MEMLOCK REMAP UTC
% LocalWords:  readmon Hierarchy defaults queues MSGQUEUE effective fstat
% LocalWords:  fchown fchmod Epoch January


%%% Local Variables: 
%%% mode: latex
%%% TeX-master: "gapil"
%%% End: 
